\documentclass{amsart}
\usepackage[utf8]{inputenc}
\usepackage{amsmath,amssymb,amsthm, graphics, comment,bm}
\usepackage{mathtools}
\usepackage{amsthm}
\usepackage{diagmac2}
\usepackage{textcomp}
\usepackage{color}
\usepackage[T1]{fontenc}
\usepackage[alphabetic]{amsrefs}
\usepackage[all,cmtip,color,matrix,arrow]{xy}
\usepackage{yfonts}
\usepackage{tikz}
\usepackage{tikz-cd}
\usepackage{stmaryrd}
\usepackage{mathrsfs}
\usepackage{enumerate}
\usepackage{hyperref}
\hypersetup{
	colorlinks=true,
	linkcolor=blue,
	filecolor=magenta,      
	urlcolor=cyan,
}
\usepackage[mathscr]{euscript}
\usepackage[all,cmtip]{xy}
\usepackage{cleveref}

\setlength{\textwidth}{\paperwidth}
\addtolength{\textwidth}{-2in}
\calclayout

\newcommand{\dotr}{\mbox{$\boldsymbol{\cdot}$}}
\newcommand{\tilda}{\accentset{\sim}}
\newcommand{\Exal}{\operatorname{Exal}}
\newcommand{\Exalstack}{\underline{\operatorname{Exal}}}
\newcommand{\bgm}{B\mathbb{G}_m}
\newcommand{\Diff}{\operatorname{Diff}}
\newcommand{\noname}{[\mathbb{A}^1/\mathbb{G}_m]}
\newcommand{\bA}{\mathbb{A}}
\newcommand{\bC}{\mathbb{C}}
\newcommand{\bF}{\mathbb{F}}
\newcommand{\bG}{\mathbb{G}}
\newcommand{\bL}{\mathbb{L}}
\newcommand{\bN}{\mathbb{N}}
\newcommand{\bP}{\mathbb{P}}
\newcommand{\bQ}{\mathbb{Q}}
\newcommand{\bR}{\mathbb{R}}
\newcommand{\bT}{\mathbb{T}}
\newcommand{\bZ}{\mathbb{Z}}
\newcommand{\kC}{\mathfrak{C}}
\newcommand{\kM}{\mathfrak{M}}
\newcommand{\km}{\mathfrak{m}}
\newcommand{\sA}{\mathscr{A}}
\newcommand{\sB}{\mathscr{B}}
\newcommand{\sC}{\mathscr{C}}
\newcommand{\sD}{\mathscr{D}}
\newcommand{\sE}{\mathscr{E}}
\newcommand{\sF}{\mathscr{F}}
\newcommand{\sG}{\mathscr{G}}
\newcommand{\sH}{\mathscr{H}}
\newcommand{\sI}{\mathscr{I}}
\newcommand{\sK}{\mathscr{K}}
\newcommand{\sL}{\mathscr{L}}
\newcommand{\sM}{\mathscr{M}}
\newcommand{\oH}{\operatorname{H}}
\newcommand{\sN}{\mathscr{N}}
\newcommand{\sO}{\mathscr{O}}
\newcommand{\sP}{\mathscr{P}}
\newcommand{\sR}{\mathscr{R}}
\newcommand{\sS}{\mathscr{S}}
\newcommand{\sU}{\mathscr{U}}
\newcommand{\sX}{\mathscr{X}}
\newcommand{\sZ}{\mathscr{Z}}
\newcommand{\cA}{\mathcal{A}}
\newcommand{\cB}{\mathcal{B}}
\newcommand{\cC}{\mathcal{C}}
\newcommand{\cD}{\mathcal{D}}
\newcommand{\cE}{\mathcal{E}}
\newcommand{\cF}{\mathcal{F}}
\newcommand{\cG}{\mathcal{G}}
\newcommand{\cH}{\mathcal{H}}
\newcommand{\cI}{\mathcal{I}}
\newcommand{\cK}{\mathcal{K}}
\newcommand{\cL}{\mathcal{L}}
\newcommand{\cM}{\mathcal{M}}
\newcommand{\cN}{\mathcal{N}}
\newcommand{\cO}{\mathcal{O}}
\newcommand{\cP}{\mathcal{P}}
\newcommand{\cQ}{\mathcal{Q}}
\newcommand{\cR}{\mathcal{R}}
\newcommand{\cS}{\mathcal{S}}
\newcommand{\cT}{\mathcal{T}}
\newcommand{\cU}{\mathcal{U}}
\newcommand{\cV}{\mathcal{V}}
\newcommand{\cW}{\mathcal{W}}
\newcommand{\cX}{\mathcal{X}}
\newcommand{\cY}{\mathcal{Y}}
\newcommand{\cZ}{\mathcal{Z}}

\newcommand{\qB}{R}
\newcommand{\qR}{\cB}

\newcommand{\Ass}{\operatorname{Ass}}
\newcommand{\E}{\operatorname{Ext}}
\newcommand{\spec}{\operatorname{Spec}}
\newcommand{\Kg}{\mathcal{K}}
\newcommand{\sY}{\mathscr{Y}}
\newcommand{\supp}{\operatorname{Supp}}
\newcommand{\Ker}{\operatorname{Ker}}
\newcommand{\lcm}{\operatorname{lcm}}
\newcommand{\Tr}{\operatorname{Tr}}
\newcommand{\Sym}{\operatorname{Sym}}
\newcommand{\Hom}{\operatorname{Hom}}
\newcommand{\Gr}{\operatorname{Gr}}
\newcommand{\Proj}{\operatorname{Proj}}
\newcommand{\Div}{\operatorname{Div}}
\newcommand{\Pic}{\operatorname{Pic}}
\newcommand{\diag}{\operatorname{diag}}
\newcommand{\End}{\operatorname{End}}
\newcommand{\Gm}{\bG_{{\rm m}}}
\newcommand{\Gmdelta}{\bG_{{\rm m,\Delta}}}
\newcommand{\GL}{\operatorname{GL}}
\newcommand{\SL}{\operatorname{SL}}
\newcommand{\Hilb}{\operatorname{Hilb}}
\newcommand{\Aut}{\operatorname{Aut}}
\newcommand{\PGL}{\operatorname{PGL}}
\newcommand{\im}{\operatorname{im}}
\newcommand{\Id}{\operatorname{Id}}
\newcommand{\GIT}{\operatorname{GIT}}
\newcommand{\Ind}{\operatorname{Ind}}
\newcommand{\sgn}{\operatorname{sgn}}
\newcommand{\Imm}{\operatorname{Imm}}
\newcommand{\Stab}{\operatorname{Stab}}
\newcommand{\Bl}{\operatorname{Bl}}
\newcommand{\Isom}{\operatorname{Isom}}
\newcommand{\uIsom}{\underline{\Isom}}
\newcommand{\Ext}{\operatorname{E}}
\newcommand{\I}{\operatorname{I}}
\newcommand{\bmu}{\bm{\mu}}
\newcommand{\dm}{{\rm DM}}
\newcommand{\id}{\rm{id}}
\newcommand{\ps}{{\rm{ps}}}
\newcommand{\Mcal}{\mathcal{M}}
\newcommand{\Mbar}{\overline{\mathcal{M}}}
\newcommand{\Mtilde}{\widetilde{\mathcal{M}}}
\newcommand{\Ctilde}{\widetilde{\mathcal{C}}}
\newcommand{\Hass}{{\rm H}}
\newcommand{\git}{{\rm GIT}}
\DeclareMathOperator{\EB}{EB}
\newcommand{\CY}{{\rm CY}}

\newtheorem{theorem}{Theorem}[section]
\newtheorem{Teo}[theorem]{Theorem}
\newtheorem*{Teo*}{Theorem}
\newtheorem{Lemma}[theorem]{Lemma}

\newtheorem{Cor}[theorem]{Corollary}

\newtheorem{Prop}[theorem]{Proposition}
\newtheorem*{Ques*}{Question}
\newtheorem{Specialcase}[theorem]{Special case of cite}
\newtheorem{Claim}{Claim}
\newtheorem{Step}{Step}
\newtheorem{Step_1}{Blow-up}
\newtheorem{Step_2}{Step}
\newtheorem{Step_3}{Step}

\theoremstyle{definition}
\newtheorem{Oss}[theorem]{Remark}
\newtheorem{Assumptions}[theorem]{Assumption}
\newtheorem*{Oss'}{Remark}
\newtheorem{EG}[theorem]{Example}
\newtheorem{Def}[theorem]{Definition}
\newtheorem*{Def*}{Definition}
\newtheorem{Notation}[theorem]{Notation}
\newtheorem{Assumption}[theorem]{Assumption}
\newtheorem{Setup}[theorem]{Setup}
\newtheorem{Remark}[theorem]{Remark}
\newtheorem{Conjecture}[theorem]{Conjecture}
\newcommand{\marg}[1]{\normalsize{{\color{red}\footnote{{\color{blue}#1}}}{\marginpar[{\color{red}\hfill\tiny\thefootnote$\rightarrow$}]{{\color{red}$\leftarrow$\tiny\thefootnote}}}}}
\newcommand{\Gio}[1]{\marg{(Giovanni) #1}}

\newcommand{\A}[1]{\marg{(Andrea) #1}}
\theoremstyle{definition}
\newtheorem{example}[theorem]{Example}


\begin{document}
\title[Stable maps to quotient stacks with a properly stable point]{Stable maps to quotient stacks with a properly stable point}
	\author[A. Di Lorenzo]{Andrea Di Lorenzo}
	\address[A. Di Lorenzo]{University of Pisa, Italy}
	\email{andrea.dilorenzo@unipi.it}
	\author[G. Inchiostro]{Giovanni Inchiostro}
	\address[G. Inchiostro]{University of Washington, Seattle, Washington, USA}
	\email{ginchios@uw.edu}
	\maketitle
	\begin{abstract}
		We compactify the moduli stack of maps from curves to certain quotient stacks $\cX=[W/G]$ with a projective good moduli space, extending previous results from quasimap theory. For doing so, we introduce a new birational transformation for algebraic stacks, the \textit{extended weighted blow-up}, to prove that any algebraic stack with a properly stable point can be enlarged so that it contains an open substack which is proper and Deligne-Mumford.  
        {As a first application, we use our main theorem to construct a compact moduli stack for certain fibered log-Calabi-Yau pairs. We further apply our result to construct a compactification of the space of maps to $\cX$ when $\cX$ is respectively: a quotient by a torus of a proper Deligne-Mumford stack; a GIT compactification of the stack of binary forms of degree $2n$; a GIT compactification of the stack of $2n$-marked smooth rational curves, and a GIT compactification of the stack of smooth plane cubics.}
        In the appendix, we give a criterion for when a morphism of algebraic stacks is an extended weighted blow-up, and we use it in order to give a modular proof of a conjecture of Hassett on weighted pointed rational curves.
	\end{abstract}
 \section{Introduction}
 Given the construction of a moduli space of curves, it is natural to ask if one can construct a moduli space of pairs consisting of a curve $C$ together with some extra data.
 When this extra data is parametrized by an algebraic stack $\cX$, this corresponds to constructing a moduli space of \textit{maps} $(C, f:C\to \cX)$ where $C$ is a curve and $f$ is a map to $\cX$.
Therefore, it is natural to try to construct a \textit{compact} moduli space, which generically parametrizes morphisms from families of smooth curves to an algebraic stack.

This problem has been solved for $\cX$ a projective variety by Kontsevich \cites{Kon, FP}, a Deligne-Mumford stack with a projective coarse moduli space by Abramovich and Vistoli \cite{AV_compactifying}, for $\cX=\cB \Gm$, $\cX=[X/\Gm]$ with $X$ projective or $\cX=\cB \GL_n$ by Caporaso, Frenkel-Teleman-Tolland, Cooper and Pandharipande \cites{Cap, FTT,cooper2022git, Pand}. Finally, when $\cX=[\spec(A)/G]$ and there is a character $\theta:G\to \Gm$ such that
$[\spec(A)(\Bbbk_\theta)^{ss}/G]$ is a Deligne-Mumford stack, the problem is again fully solved via quasimap theory by Cheong, Ciocan-Fontanine, Kim and Maulik \cites{ciocan2010moduli, ciocan2014stable, orbifold_qmap_theory}. Similarly, Halpern-Leistner and Herrero solve the problem for a fixed smooth $C$ \cite{HLH23}.
\subsection{Main results} We focus on stacks which are global quotients and admitting a projective good moduli space, as many moduli of varieties are as such. \Cref{thm_intro_simplified} is a simplified version of our main results:\begin{Teo}\label{thm_intro_simplified}
Let $\cX=[W/G]$ be a quotient stack with $G$ reductive, with a projective good moduli space $\cX\to X$, and admitting a dense open subset $U\subseteq X$ such that $\cX\times_XU$ is Deligne-Mumford. Then:
\begin{enumerate}
\item there is a proper Deligne-Mumford stack $\cQ_{g}(\widetilde{\cX},\widetilde{\cX}_\dm,\beta)$ which generically parametrizes morphisms $\phi:C\to \cX$ of class $\beta$ satisfying a stability condition, from a smooth curve $C$ of genus $g$; and
\item the boundary of $\cQ_g(\widetilde{\cX},\widetilde{\cX}_\dm,\beta)$ correspond to maps $\phi:\sC\to \widetilde{\cX}$ of class $\beta$, %\textcolor{blue}{such that $\phi^{-1}(\cX_{\dm})\neq\emptyset$}G: non credo sia vero, per sicurezza lo levo
and satisfying a stability condition, from a twisted curve $\sC$ to an enlargement $\cX\subseteq \widetilde{\cX}$ of $\cX$. 
\end{enumerate} 
 The enlargement $\cX\subseteq \widetilde{\cX}$ is a dense open embedding, with $\widetilde{\cX}$ a global quotient stack, with the same good moduli space of $\cX$. Moreover, if $\cX$ smooth and contains a dense open substack which is proper and Deligne-Mumford, then no enlargement is needed: one can take $\widetilde{\cX}=\cX$.
 \end{Teo}
 The stability condition is introduced in \Cref{def_stable_qmap}, whereas the enlargement $\cX\subseteq \widetilde{\cX}$ is obtained via \Cref{thm_intro_you_can_embed_a_Stack_in_one_with_open_dm_by_taking_deformations_to_nc}; we will explain it in more detail in \Cref{subsection_strategy_of_proof}.
 For the moreover part, we remark that there are weaker assumptions that one can make on $\cX$ so that no enlargement is needed; see \Cref{teo_intro_cX_contains_open_proper_dm} and \Cref{cor_when_X_dm_is_ss_locus_for_a_lb}.

 \subsection{First applications} \Cref{section_examples_qmaps} and \Cref{section_examples_extended} are devoted to applications of \Cref{thm_intro_simplified}.
 
{In \Cref{section_examples_qmaps} we report two examples when the enlargement $\cX\subseteq \widetilde{\cX}$ of \Cref{thm_intro_simplified} is not needed:
\begin{enumerate}
\item when $\cX$ is a certain moduli of boundary-polarized Calabi-Yau pairs. This application is particularly relevant, and we discuss it further in \Cref{subsectuon_k_dim_1} below, and
    \item when $\cX$ admits a morphism $\cX\to \cB\Gm^n$ representable in Deligne-Mumford stacks (for example, all stacks of the form $[W/\Gm^n]$ admitting a projective good moduli space).
    \end{enumerate}  }

{In \Cref{section_examples_extended} instead we discuss an application of \Cref{thm_intro_simplified} when $\cX$ does not contain a proper Deligne-Mumford stack: we take $\cX$ to be the GIT compactification of the moduli stack of smooth plane cubics. 
In this case, we describe what happens if we enlarge $\cX\subseteq \widetilde{\cX}$ using extended weighted blowups, and one can still give a modular interpretation of a morphism $\phi:C\to \widetilde{\cX}$, in terms of plane cubics, by composing $\phi$ with the projection $\widetilde{\cX}\to \cX$.

\subsection{Applications to moduli of varieties of Kodaira dimension one}\label{subsectuon_k_dim_1} One of the main application of \Cref{thm_intro_simplified} is for constructing certain moduli of varieties.
 While there has been tremendous work for constructing moduli spaces for varieties of general type \cites{KSB,al1,HMX18,kol19} culminating in \cite{kollarmodulibook}, and similarly for Fano varieties \cites{blum2019uniqueness,alper2020reductivity,blum2021properness,blum2022openness,xu2020uniqueness,MR4445441}, the situation for varieties of intermediate Kodaira dimension is less clear.
 
 The main application of \Cref{thm_intro_simplified}, and specifically \Cref{teo_intro_cX_contains_open_proper_dm}, is \Cref{teo_intro_kodairadim1}: the construction of a compact moduli stack of certain pairs $(Y,cD)$ of Kodaira dimension one, which are fibrations $f\colon (Y,cD)\to C$ with fibers in a moduli space of log Calabi-Yau pairs. The quasimap condition and the stability condition of \Cref{def_stable_qmap} are translated in terms of the birational geometry of the corresponding fibered variety: labels (Q), (S) and (N) in the statement of \Cref{teo_intro_kodairadim1} stand for quasimap condition, stability condition and numerical condition.
\begin{Teo}\label{teo_intro_kodairadim1}
    We denote by $\cX:=\cD\cP^{\CY}_m$ the moduli spaces of \cites{blum2024good} which parametrize boundary polarized Calabi-Yau pairs $(S,cD)$ of dimension at most 2 admitting a $\bQ$-Gorenstein smoothing, and with $c<1$. Then the algebraic stack $\cX$ contains a proper Deligne-Mumford stack $\cX_\dm$, denoted in \cite{blum2024good} by $\cD\cP^{\operatorname{KSBA}}_m$, such that the  assumptions of \Cref{teo_intro_cX_contains_open_proper_dm} apply for the inclusion $\cX_\dm\subseteq \cX$. In particular:
    \begin{enumerate}
        \item the stack $\cQ_g(\cX,\cX_\dm,\beta)$ compactifies the space of maps $\pi\colon(Y,cD)\to C$ with fibers in $\cX$ such that:
        \begin{itemize}
            \item[(Q)] the curve $C$ is smooth and the generic fiber of $\pi$ is has klt singularities,
            \item[(S)] either $\omega_C$ is ample, or not all the fibers of $\pi$ are $S$-equivalent,
            \item[(N)] the family $\pi$ comes from a map $C\to \cX$ of class $\beta$ from a curve of genus $g$.
        \end{itemize}
        \item the boundary of $\cQ_g(\cX,\cX_\dm,\beta)$ parametrizes families $\pi:(\sY,c\sD)\to \sC$ of pairs in $\cX$ with fibers in $\cX$, fibered over a twisted curve $\sC$, such that:
        \begin{itemize}
            \item[(Q)] the set $\Delta:=\{p\in \sC:(\sY_p,(c+\epsilon)\sD_p)$ does \underline{not} have semi-log-canonical singularities for any $0<\epsilon \ll 1\}$ is a finite union of smooth points $\sC$,
            \item[(S)] if $\sR\subseteq \sC$ is an irreducible component such that $\deg(\omega_\sC |_\sR)< 0$, then not all the fibers of $\pi|_\sR\colon(\sY|_\sR,c\sD|_\sR)\to \sR$ are $S$-equivalent; whereas if $\deg(\omega_\sC |_\sR)=0$, then not all fibers of $\pi|_\sR$ are isomorphic, and
            \item[(N)] the family $\pi$ comes from a map $\sC\to \cX$ of class $\beta$ from a twisted curve of genus $g$.
        \end{itemize}
    \end{enumerate}
\end{Teo}
\begin{Remark}Being $S$-equivalent is a weaker condition than being isomorphic; one can consult \cite{ascher2023moduli}*{Definition 6.7} for the definition.
    Moreover, from \cite{ascher2023moduli}*{Proposition 14.17}, one can rephrase condition (S) in point (1) of \Cref{teo_intro_kodairadim1} by asking that the moduli part in the canonical bundle formula for $\pi$ has positive degree. 
\end{Remark}

\subsection{Strategy of proof}\label{subsection_strategy_of_proof}
 The strategy for proving \Cref{thm_intro_simplified} is as follows. We begin by showing that if $\cX$ is an algebraic stack which contains a \textit{proper} Deligne-Mumford stack $\cX_\dm\subseteq \cX$ as an open substack, then (in some cases) one can build on results from quasimap theory to construct a moduli stack of maps to $\cX$. More specifically, we prove the following.
 
\begin{Teo}\label{teo_intro_cX_contains_open_proper_dm} Let $\cX$ be an algebraic stack with a projective good moduli space $\cX\to X$, which is a global quotient $[W/G]$ by a reductive group $G$, and which contains an open substack $\cX_\dm\subseteq \cX$ which is a proper Deligne-Mumford stack.
Assume that $\cX_\dm$ is the $\cL$-semistable locus of $\cX$ over $X$ for a line bundle $\cL$. 

Then there is an algebraic stack, which we denote by $\cQ_{g,n}(\cX,\cX_\dm,\beta)$, which is proper, Deligne-Mumford, it parametrizes stable quasimaps $\phi:(\sC,x_1,\ldots,x_n)\to \cX$, defined in \Cref{def_stable_qmap}, from an $n$-pointed twisted curve $\sC$ of genus $g$ and class $\beta$ to $\cX$. If $\cX$ is smooth, the moduli space $\cQ_{g,n}(\cX,\cX_\dm,\beta)$ carries a perfect obstruction theory.
\end{Teo}
The condition that $\cX_\dm$ is the semistable locus for a line bundle is not a heavy restriction. For example, when $\cX$ is non-singular, every open substack of $\cX$ which is proper and Deligne-Mumford is the semistable locus of $\cX$ over $X$ for a line bundle $\cL$. On the other hand, the condition that $\cX$ contains an open which is Deligne-Mumford and proper is heavier: it is not hard to find algebraic stacks which are global quotients and have a projective good moduli space, which do not have an open substack which is proper and Deligne-Mumford; simple examples are given by the GIT moduli space of plane cubics, or the GIT moduli space of $2n$ unordered points on $\bP^1$. \Cref{teo_intro_cX_contains_open_proper_dm} alone does not lead to a moduli of maps to these stacks. 

\Cref{teo_intro_cX_contains_open_proper_dm} can be interpreted as a relative version of \cites{ciocan2010moduli,orbifold_qmap_theory}, relative over $X$. Indeed, the results in \textit{loc. cit.} apply to stacks of the form $[\spec(A)/G]$ with the semistable locus given by a character of $G$. The main observation is that a good moduli space is, \'etale locally on $X$, of the form $[\spec(A)/G]$, so one can leverage the results in \cites{ciocan2010moduli,orbifold_qmap_theory} to obtain \Cref{teo_intro_cX_contains_open_proper_dm}. The definition of \textit{stable quasimap} is given in \Cref{def_stable_qmap}, and is a combination of the definition of quasimap of \cite{orbifold_qmap_theory} with the definition of twisted map of \cite{dli1}. The class $\beta$ is a numeric invariant needed to prove that $\cQ_g(\cX,\cX_\dm,\beta)$ is bounded, it is defined in \Cref{def_beta}.

To compactify the space of maps to algebraic stacks which do \textit{not} contain an open substack that is proper and Deligne-Mumford, {a natural approach would be to find a bigger stack $\widetilde{\cX}\supset \cX$ with an open substack $\widetilde{\cX}_\dm$ which is proper and Deligne-Mumford:}
%It would be even better to find an enlargement which still has the same good moduli space $X$ of $\cX$, as then the stability conditions would be mostly unaffected \Gio{lo ho commentato siccome non cosa sia la stab condition se non abbiamo l'aperto dm. in other terms, non mi parsa la frase sopra. We prove that one can always find such an enlargement.}
\begin{Teo}\label{thm_intro_you_can_embed_a_Stack_in_one_with_open_dm_by_taking_deformations_to_nc}
    Let $\cX$ be an algebraic stack with a good moduli space $p:\cX\to X$, and with a dense open $U\subseteq X$ such that $\cX\times_XU$ is Deligne-Mumford. Then there is an algebraic stack $\widetilde{\cX}$ with an open embedding $i:\cX\hookrightarrow \widetilde{\cX}$ such that:
    \begin{enumerate}
        \item $\widetilde{\cX}$ has a good moduli space which is isomorphic to $X$ via the inclusion $i$,
        \item there is a line bundle $\cL_{DM}$ on $\widetilde{\cX}$ such that $\widetilde{\cX}(\cL_{DM})^{ss}_X$ is  Deligne-Mumford and proper over $X$,
        \item if $\cX$ is a global quotient, then $\widetilde{\cX}$ can be chosen to be a global quotient,
        \item there is a morphism $\pi:\widetilde{\cX}\to \cX$ which is an isomorphism over $(p\circ \pi)^{-1}(U)$,
        \item the morphism $\pi\circ i$ is isomorphic to the identity.
    \end{enumerate}
    In particular, from \Cref{teo_intro_cX_contains_open_proper_dm}, if $\cX$ is a global quotient and one fixes an integer $g$ and a class $\beta$, there is a moduli space $\cQ_g(\widetilde{\cX},\widetilde{\cX}(\cL_{DM})^{ss}_X,\beta)$ of stable quasimaps to $\widetilde{\cX}$ which is a proper Deligne-Mumford stack. 
\end{Teo}
In particular, for every projective variety $\overline{W}$ with the action of a reductive group $G$ and with an ample $G$-linearized line bundle $L$, one can use \Cref{teo_intro_cX_contains_open_proper_dm} and \Cref{thm_intro_you_can_embed_a_Stack_in_one_with_open_dm_by_taking_deformations_to_nc} to compactify the space of maps from smooth curves $\phi\colon C\to[\overline{W}^{ss}(L)/G]$ such that the map $X\to W/\!\!/G$ is Kontsevich stable.

The construction of the enlargement $\widetilde{\cX}$ of $\cX$ is \textit{explicit}, using \cite{ER21}. Specifically, the proper Deligne-Mumford stack $\widetilde{\cX}(\cL_{DM})^{ss}_X$ is obtained by the Kirwan desingularization procedure of \cites{ER21,kirwan}. For doing so, we define \textit{extended weighted blow-ups}, which are a generalization of weighted blow-ups of \cite{QR}, and
the enlargement $\widetilde{\cX}$ of $\cX$ is obtained by performing a sequence of extended weighted blow-ups.

\subsection{Extended weighted blow-ups}
Extended weighted blow-ups, defined in \Cref{def_ext_w_blowup}, are birational transformations needed to construct the enlargement $\cX\subseteq \widetilde{\cX}$ of \Cref{thm_intro_you_can_embed_a_Stack_in_one_with_open_dm_by_taking_deformations_to_nc}. The main properties of an extended weighted blowup $\pi:\EB_{I^\bullet}\cX\to \cX$ are:
\begin{enumerate}
    \item the morphism $\pi$ induces an isomorphism on good moduli spaces,
    \item the algebraic stack $\EB_{I^\bullet}\cX$ admits an open embedding $i:\cX\to \EB_{I^\bullet}\cX$, and
    \item there is an open substack $\cU\subseteq \EB_{I^\bullet}\cX$ which is the weighted blow-up of $\cX$ at the weighted ideal sequence $I^\bullet$. 
\end{enumerate}
In the appendix to this paper, we establish a criterion (see \Cref{prop:criterion}) for telling when a morphism of algebraic stacks is an extended weighted blow-up. For $\cX$ the GIT compactification of the moduli stack of $2n$ distinct points on $\bP^1$, we use this criterion to show that a \textit{modular} enlargement $\cX\subseteq \widetilde{\cX}$ is an extended weighted blow-up,
and we use it to give a modular proof of the following conjecture of Brendan Hassett:
\begin{Conjecture}\label{conjecture_hassett}
    Let $\overline{\cM}_{0,(\frac{1}{n}  + \epsilon,\ldots,\frac{1}{n}  + \epsilon)}$ the moduli space of weighted $2n$-pointed stable curves of genus 0 and with weights $\frac{1}{n}+\epsilon$. Let $\widetilde{M}_{0,2n}$ the coarse moduli space of $[\overline{\cM}_{0,(\frac{1}{n}  + \epsilon,\ldots,\frac{1}{n}  + \epsilon)}/S_{2n}]$ where the action permutes the $2n$ points. Let $\bP(\oH^0(\bP^1,\cO_{\bP^1}(2n)))/\!\!/\PGL_2$ the GIT moduli space of $2n$ unordered points on $\bP^1$, linearized via $\cO(2)$. Then there is a map $\widetilde{M}_{0,2n}\to \bP(\oH^0(\bP^1,\cO_{\bP^1}(2n)))/\!\!/\PGL_2$ which is a weighted blow-up.
\end{Conjecture}
\Cref{conjecture_hassett} was proven in \cite{moduli_pointed_curves_via_git} by showing that the corresponding statement for ordered points is true, namely that there is a map $\overline{\cM}_{0,(\frac{1}{n}  + \epsilon,\ldots,\frac{1}{n}  + \epsilon)}\to (\bP^1)^n/\!\!/\PGL_2$ that is a weighted blow-up (in this case, it is the Kirwan desingularization). The conjecture was then proved by observing that the previous map is $S_{2n}$-equivariant, and the $S_{2n}$-quotient induces $\widetilde{M}_{0,2n}\to \bP(\oH^0(\bP^1,\cO_{\bP^1}(2n)))/\!\!/\PGL_2$ which is a weighted blow-up.
It is natural to wonder if one can construct the map above from the modular properties of the quotient stack $[\bP(\oH^0(\bP^1,\cO_{\bP^1}(2n)))(\cO(2))^{ss}/\PGL_2]$, which we will denote by $\cC^{\git}_{2n}$, directly. Specifically, the question is whether one can construct a moduli stack of rational $2n$-marked curves $\Mcal$, having $\widetilde{M}_{0,2n}$ as coarse moduli space, and with a map $\cM\to \cC^{\git}_{2n}$ which induces the desired weighted blow-up on good moduli spaces. This is the content of the next theorem, proved in \Cref{appendix}.

\begin{Teo}\label{thm_intro_hassett}
Consider pairs $(P,D)$ where:
        \begin{itemize}
            \item either $P\cong \bP^1$ or $P$ is a twisted conic, namely a twisted curve that is the nodal union of two root stacks of $\bP^1$ at $0$, glued along the $\cB\bmu_2$ gerbe, 
            \item $D$ is supported on the smooth locus of $P$, has degree $2n$ and each point of $P$ has multiplicity at most $n$ in $D$ (in other terms, $(P,\frac{1}{n}D)$ is semi log-canonical), and
            \item the divisor $\omega_P(\frac{1}{n}D)$ is nef.
        \end{itemize}
        Then:
        \begin{enumerate}
            \item there is an algebraic stack $\Ctilde^{\CY}_{2n}$ which is a moduli stack for pairs as above,
            \item there is a morphism $\Phi:\Ctilde^{\CY}_{2n}\to \cC^{\git}_{2n}$ which is an extended weighted blow-up with center the polystable but not stable point of $\cC^{\git}_{2n}$,
            \item if we denote by $X$ the good moduli space of $\cC^{\git}_{2n}$, there is a line bundle $\cL$ on $\Ctilde^{\CY}_{2n}$ such that the open substack $\Ctilde^{\CY}_{2n}(\cL)^{ss}_X\subseteq\widetilde{\cC}^\CY_{2n}$ is a proper Deligne-Mumford stack and its coarse moduli space is a weighted blow-up of $C_{2n}^{\git}$ at the unique polystable but not stable point, and
            \item the pairs in $\Ctilde^{\CY}_{2n}(\cL)^{ss}_{X}$ are pairs $(P,D)$ as above such that each point of $P$ has multiplicity at most $n-1$ in $D$. 
        \end{enumerate}
        In particular, the weighted blow-up constructed in \cite{moduli_pointed_curves_via_git} is induced on good moduli spaces by the morphism $\Ctilde_{2n}^{\CY}(\cL)^{ss}_{X}\hookrightarrow\Ctilde_{2n}^{\CY}\xrightarrow{\Phi}\cC^{\git}_{2n}$.
        \end{Teo}
        \Cref{remark_why_we_need_twisted_nodes} explains why one needs to consider twisted curves rather than schematic curves.


 \subsection{Organization of the paper} The paper is organized as follows. In \Cref{section_background} we report a few results on algebraic stacks which will be needed in the rest of the paper. In particular, in \Cref{subsection_relative_GIT} we report the main definitions from relative GIT from \cite{Alp}, and we prove \Cref{prop_stable_locus_of_pi*L^n_otimes_G} which will be used a few times in the paper. 
 
 In \Cref{sec:qmaps} we give the definition of quasimaps with target an algebraic stack which has an open substack that is proper and Deligne-Mumford, and then we proceed to proving \Cref{teo_intro_cX_contains_open_proper_dm}. Many of the ideas and proofs are generalizations of ideas in \cites{orbifold_qmap_theory, ciocan2014stable}, with the exception of \Cref{section_boundedness_qmaps} which requires more delicate arguments. 

 In \Cref{section_examples_qmaps} we report a few {applications} of \Cref{teo_intro_cX_contains_open_proper_dm}, the most interesting one being a moduli of certain fibrations in Calabi-Yau pairs.
 
 \Cref{section_extended_blowups} has the definition of extended weighted blowup, which we use to prove \Cref{thm_intro_you_can_embed_a_Stack_in_one_with_open_dm_by_taking_deformations_to_nc}. 
 
In \Cref{section_examples_extended} we apply our main theorem when the target stack $\cX$ is the GIT compactification of the moduli stack of smooth plane cubics.

In \Cref{appendix} we apply the aforementioned criterion for being an extended blow-up for proving \Cref{thm_intro_hassett}.

 \subsection{Acknowledgements} We are thankful to Luca Battistella, Dori Bejleri, Danier Halpern-Leistner, Felix Janda, Nick Kuhn, Denis Nesterov and Dhruv Ranganathan for helpful conversations. We are thankful to Andres 
Fernandez Herrero for giving us detailed feedback on an earlier version of this manuscript, and for suggesting an improvement on \Cref{lemma_one_can_check_rep_at_poli_pts}.
 
 \subsection{Conventions.} We work over a field, denoted by $k$, which is of characteristic 0 and algebraically closed. Unless otherwise stated, all the stacks have affine diagonal, are of finite type over $k$, and the stabilizers of the points are reductive (but a priori not connected).
 We denote the set of characters (resp. cocharacters) of a group $G$ by $\mathbf{X}(G)$ (resp. $\mathbf{X}(G)^*)$, and the natural pairing between them by \[<\cdot, \cdot >:\text{ }\mathbf{X}(G)^*\times \mathbf{X}(G)\to \bZ. \] 
 Given a character $\chi\in \mathbf{X}(G)$, we denote by $\Bbbk_\chi$ the representation of $G$ with character $\chi$.
 We will denote by $\bmu_n$ the group of $n$-th roots of unity of $k$ and by $\Theta:=[\bA^1/\Gm]$. Unless otherwise stated, all curves are assumed to be connected and proper.
 \section{Background}\label{section_background}
In this section we report a few results on algebraic stacks and geometric invariant theory which will be useful for the rest of the manuscript.
 \subsection{Results on algebraic stacks}\label{subsection_algebraic_stacks}
This subsection consists of a list of results that will be useful later. We choose to isolate them from the rest of the paper because we are using them multiple times, and because some of these results could be useful in general. 

\begin{comment}
The following is a slight generalization of \cite{AHH}*{Lemma 3.24}.
\begin{Lemma}\label{lemma_a_pt_in_the_closure_of_another_pt_can_be_reached_with_theta}
    Let $\cX$ be an algebraic stack, let $x$ and $x_0$ be two $k$-points in $\cX$ with $x_0\in\overline{\{x\}}$. Then there is a morphism $\Theta\to \cX$ sending $1\mapsto x$ and $0\mapsto x_0.$
\end{Lemma}
\begin{proof}
    From \cite{luna}*{Theorem 1.1} there is an affine scheme with an action of $G_{x_0}:=\Aut_\cX(x_0)$, with a fixed point $y_0$ and an \'etale morphism $\iota:[\spec(A)/G_{x_0}]\to \cX$ sending $y_0\mapsto x_0$. Since $\iota$ is open and quasi-finite, there is a point $y\in [\spec(A)/G_{x_0}]$ such that $\iota(y) = x$ and such that $y_0\in \overline{\{y\}}$. Since $y_0$ is now a \textit{closed} point of $[\spec(A)/G_{x_0}]$, we can apply \cite{AHH}*{Lemma 3.24}.
\end{proof}
\end{comment}


\begin{Lemma}\label{lemma_generalization_purity_tsm1}
    Assume that $S$ is a separated Deligne-Mumford stack of dimension 2, with finite quotient singularities. Let $\cX$ be an algebraic stack with a good moduli space $\cX\to X$, and let $s\in S$ be a closed point of $S$. Assume that there is a diagram as follows:
    \[
    \xymatrix{S\smallsetminus\{s\}\ar[d]\ar[r] & \cX\ar[d]\\S\ar[r]^f & X.}
    \]
    {Then there is an unique $S_2$ Deligne-Mumford stack $\cS$ with $\cS\to S$ a relative coarse moduli space that is an isomorphism on $S\smallsetminus \{s\}$, and such that there is an extension $\phi:\cS\to \cX$ which is representable. Such an extension is unique, the stack $\cS$ has finite quotient singularities, and if $S$ is smooth then $\cS\to S$ is an isomorphism.}  
\end{Lemma}

\begin{proof}From descent and from the uniqueness part of the statement, it suffices to treat the case when $S$ is a scheme. First recall that if $\cS\to S$ is a coarse moduli space morphism from a Deligne-Mumford stack $\cS$, then $\cS\to S$ is separated. Indeed, from \cite{luna}*{Theorem 4.12}  \'etale locally  it is of the form $[\spec(A)/\Gamma]\to \spec(A^\Gamma)$ for a finite group $\Gamma$,  and the latter is a separated morphism.

Since $S$ has finite quotient singularities, from \cite{Vis}*{Proposition 2.8} there is a smooth Deligne-Mumford stack $\cY$ with coarse moduli space $\cY\to S$ that is an isomorphism on the smooth locus of $S$.
There is a map $\cY\smallsetminus\{s\}\to \cX$, and to show that it extends to $\cY\to \cX$ it suffices to show that it extends uniquely up to passing to an \'etale neighbourhood of $s\in \cY$. We choose an \'etale neighbourhood of $s$ by pulling back an \'etale neighbourhood of $f(s)\in X$. From \cite{AHH}*{Proposition 2.10}, there is a cartesian square as follows:
\[
\xymatrix{[\spec(A)/\GL_n]\ar[r]\ar[d] & \cX\ar[d] \\ f(s)\in \spec(A^{\GL_n})\ar[r] & X.}
\]
Now the desired extension $\cY\to [\spec(A)/\GL_n]$ exists and is unique from \cite{dli1}*{Proposition 4.1}, so from descent there is a unique extension $\cY\to \cX$ of $\cY\smallsetminus \{s\}\to \cX$.

The morphism $\cY\to \cX$ might not be representable, but we can take the relative coarse moduli space $\cY\to \cS\to \cX$ of $\cY\to \cX$. Then $\cS$ has finite quotient singularities, as it is a relative coarse moduli space of a smooth Deligne-Mumford stack, and has a representable morphism $\cS\to 
\cX$. 

We now show that $\cS$ is unique, by contradiction. Let $\cT\to \cX$ be another representable extension, where $\cT$ is an $S_2$ Deligne-Mumford stack with coarse space $S$ and such that $\cT\smallsetminus\{s\}\simeq S\smallsetminus\{s\}$. Then by the same argument as above (except replacing $\cX$ with $\cT$), there is a morphism $\cY\to\cT$ extending the section $S\smallsetminus\{s\}\to \cT\smallsetminus\{s\}$. As $\cY\to S$ and $\cT\to S$ are coarse moduli spaces, $\cY\to \cT$ is a relative coarse moduli space from \Cref{lemma_map_between_S2_DM_stacks_that_is_isom_in_condim_1_with_same_cms_is_relative_cms}. Then $\cT\cong \cS$ as they are both relative coarse moduli spaces of $\cY\to \cX$. 
The map $\phi$ is unique from \Cref{lemma_two_maps_that_agree_in_codim_2_are_the_same}.
\end{proof}
\begin{Lemma}\label{lemma_map_between_S2_DM_stacks_that_is_isom_in_condim_1_with_same_cms_is_relative_cms}
    {Let $\phi:\cX\to \cY$ be a morphism of separated Deligne-Mumford stacks with coarse moduli spaces $X$ and $Y$ respectively, and such that the morphism $X\to Y$ is an isomorphism. Assume that $\cX$ and $\cY$ are $S_2$, and that $\phi$ is an isomorphism over an open dense $V\subseteq \cY$ with complement of codimension at least two. Then the natural map $\cO_\cY\to \phi_*\cO_\cX$ is an isomorphism. In particular, $\phi$ is a relative coarse moduli space.} 
\end{Lemma}
\begin{proof}
{Observe first that $\cX\to X$ is proper: it is universally closed from \cite{AHH}*{Theorem A.8}, of finite type by our conventions, and separated as $\cX$ is separated and $\cY$ has separated (in particular, proper) diagonal.
Since statement is local on $\cY$, we can replace $\cY$ with an \'etale cover of it $U\to \cY$. Let $\cX_U:=U\times_\cY \cX$ and let $\cX_U\to X_U$ its coarse moduli space. Observe that $X_U$ is $S_2$. Indeed, it suffices to check it is such \'etale locally. From \cite{AV_compactifying}*{Lemma 2.2.3} every point $p\in X_U$ admits an \'etale neighbourhood $\spec(A^\Gamma)\to X_U$ such that the following diagram is cartesian, with $\Gamma$ a finite group \[\xymatrix{[\spec(A)/\Gamma]\ar[r] \ar[d] & \cX_U \ar[d] \\ \spec(A^\Gamma)\ar[r] &X_U.}\] Since $\cX$ is $S_2$, also $\cX_U$ and $\spec(A)$ are $S_2$,
so $\spec(A^\Gamma)$ and $X_U$ will be $S_2$ from \cite{KM98}*{Proposition 5.4}. Now, $X_U\to U$ is quasi-finite as $\cX\to \cY$ is such, since $X\to Y$ is an isomorphism. The map $X_U\to U$ is proper as $\cX\to \cY$ is such, so it is finite, hence affine. But $\cO_{X_U}(X_U)$ and $\cO_U(U)$ agree as both $U$ and $X_U$ are $S_2$, the map $\cX_U \times_\cY V\to U\times_\cY V$ is an isomorphism by assumption, and $U\times_\cY V$ has complement of codimension at least two.}
\end{proof}
\begin{Lemma}\label{lemma_two_maps_that_agree_in_codim_2_are_the_same}
    Let $\cX$ be an algebraic stack with affine diagonal, let $S$ be an $S_2$ algebraic stack, and let $U\subset S$ an open subset with complement of codimension at least two. Assume we are given two morphisms $f,g: S\to \cX$ which are isomorphic when restricted on $U$. Then the isomorphism extends uniquely to $S$.
\end{Lemma}
\begin{proof}
Using descent, can assume that $S=\spec(A)$ is an affine scheme.
    Consider $S\to \cX\times \cX$ given by $(f,g)$, and let $I:=\cX\times_{\cX\times \cX}S\to S$ be the second projection. Sine the diagonal of $\cX$ is affine, $I\to S$ is affine, so $I=\spec(B)$. There is an open $U\subseteq S$ where $f$ and $g$ are isomorphic, and thus a map $U\to\spec(B)$ inducing $B\to \cO_S(U)$. As $S$ is $S_2$, we have that $\cO_S(U)=\cO_S(S)$ and as $S$ is affine $\cO_S(S) = A$, so we have a map $B\to A$. This induces the desired map $S\to I$ extending $U\to I$.  
\end{proof}
\begin{Lemma}\label{lemma_locus_Where_gms_is_isom_is_open}
    Let $\cX\to X$ a good moduli space, with $X$ separated. Then the set $\cS:=\{x\in X:\pi^{-1}(x)\to x$ is an isomorphism\} is open in $X$.
\end{Lemma}
\begin{proof}
    Indeed, first observe that the points in $\cS$ are \textit{stable} points in the sense of \cite{ER21}, and from \cite{ER21}*{Proposition 2.6} the set of stable points $X^{st}$ in $X$ is open. So up to replacing $X$ with $X^{st}$ and $\cX$ with $\cX\times_{X}X^{st}$, it suffices to check that the points where $\pi$ is an isomorphism is open, with the additional assumption that all the points on $X$ are stable. Then, as the dimension of the automorphisms is constant for stable points, in order for $\cS$ to be non-empty the stack $\cX$ has to be Deligne-Mumford: the dimension of the automorphism group of every point has to be 0. Then $\cX$ is also separated, as $S$-complete Deligne-Mumford stacks are separated. The result for separated Deligne-Mumford stacks is \cite{conrad2007arithmetic}*{Theorem 2.2.5 (2)}. 
\end{proof}
\begin{Lemma}\label{lemma_rep_after_fiber_product}
    Let $f:\cX\to \cY$ and $g:\cZ\to \cY$ be morphisms of algebraic stacks, with $\cF:=\cX\times_\cY\cZ$. Let $x\in \cF$ such that $\Aut_{\cF}(x)=\{1\}$. Then $f$ is representable at $\pi_1(x)$, where $\pi_1:\cF\to \cX$ is the first projection.
\end{Lemma}
\begin{proof}
    This follows from the universal property of the fiber product, arguing as in \cite{Inc}*{Lemma 2.11}.
\end{proof}

\begin{Lemma}\label{lemma_one_can_check_rep_at_poli_pts}
    Let $f:\cX\to \cY$ be a $S$-complete morphism, with $\cX$ quasicompact. Let $p$ be a point in $\cX$ and $x$ a closed $k$-point in the closure of $p$. If $f$ is representable at $x$, then it is representable at $p$.
\end{Lemma}
If $\cX$ admits a separated good moduli space, then $f$ is automatically $S$-complete. So to check that a morphism of algebraic stacks with a separated good moduli space is representable, it suffices to check that it is representable at the closed points of the domain. The following proof was suggested to us by Andres Fernandez Herrero, improving an argument we had in a previous version of this paper.

\begin{proof}It suffices to show that the relative inertia $I_{\cX/\cY}$ is trivial at $p$. Consider $i:\spec(A)\to \cY$ a smooth morphism such that the map $x\to \cX\to \cY$ lifts to $x\to \spec(A)\to \cY$. Let $g\colon\cX':=\cX\times_{\cY}\spec(A) \to \cX$ be the induced map. As the relative inertia commutes with smooth base change, we have that $I_{\cX'/\spec(A)}$ is trivial at $x'$ for every $x'\in g^{-1}(x)$, and it suffices to check that there is a point $p'$ in the fiber of $p$ such that $I_{\cX'/\spec(A)}$ is trivial at $p'$. We can choose $x'$ and $p'$ so that $x'$ is a closed point and is in the closure of $p'$.  

Observe that $\cX'$ is S-complete, since being S-complete is stable under base change from \cite{AHH}*{Remark 3.40} and $\spec(A)$ is S-complete. Then $\cX'$ is locally linearly reductive from \cite{alper2019cartan} and from \cite{luna}. Then $\cX'$ has unpunctured inertia from \cite{AHH}*{Theorem 5.3}.

Consider a morphism $g:\spec(A)\to \cX'$, where $A$ is a Henselian local ring, and assume that $g$ sends the generic point of $\spec(A)$ to $p'$ and the closed one to $x'$. Consider the pull-back of the inertia $g^*\cI_{\cX'}$ on $\spec(A)$, which is an affine group scheme over $\spec(A)$. Our goal is to show that it is the trivial group scheme. By upper-semicontinuity of fiber dimension for group schemes, since $\cI_{\cX'}$ is trivial at $x'$ then $g^*\cI_{\cX'}\to \spec(A)$ is quasifinite. Since we are on a field of characteristic 0, the group
scheme is unramified (that is, the fibers of $g^*\cI_{\cX'}\to \spec(A)$ are \'etale). Then from the structure of affine unramified morphisms over strictly Henselian rings \cite{stacks-project}*{\href{https://stacks.math.columbia.edu/tag/00UY}{Tag 00UY}}, we have that $g^*\cI_{\cX'}=\cI_1\sqcup \cI_2$ where $\cI_1$ is a disjoint union of finitely many closed
subschemes of $\spec(A)$ and the image of $\cI_2\to \spec(A)$ does not contain the closed point of $\spec(A)$. By the definition of unpuncuted inertia and since $\cX'$ has unpuncuted inertia, we have $\cI_2=\emptyset$. Since the special fiber of $g^*\cI_{\cX'}$ is trivial, $\cI_1$ must consist of a single closed $Z \to \spec(A)$. Observe that $Z\cong \spec(R)$ (it is the identity section), so we have that $g^*\cI_{\cX'}$ is trivial as desired.
\end{proof}

\begin{Remark}
    It might be tempting to try to prove \Cref{lemma_one_can_check_rep_at_poli_pts} by proving that the locus where $f$ is representable is open. This is false: the action of $\Gm\rtimes \bmu_2$ on $\bA^2$, where $\Gm$ acts with weights $(1,-1)$ and $\bmu_2$ swaps the two axis, has $\{(x,y):xy=0\}\smallsetminus\{(0,0)\}$ as  locus with trivial stabilizers. So the locus in $[\bA^2/\Gm\rtimes \bmu_2]$ where $[\bA^2/\Gm\rtimes \bmu_2]\to \spec(k)$ is representable is not open. Similarly, the stack $\cX:=[\bA^2\smallsetminus\{(0,0)\}/\Gm\rtimes \bmu_2]$ is Deligne-Mumford, non-separated, and the locus where $\cX\to\spec(k)$ is representable is not open. 
    \end{Remark}
    Finally, we will need the following lemma.
\begin{Lemma}\label{lm:stein of coh aff is gms}
    Let $f\colon \cX \to Y$ be a cohomologically affine morphism. Then $\pi\colon \cX\to X=\spec_Y(f_*\cO_{\cX})$ is a good moduli space.
\end{Lemma}
\begin{proof}
    Consider the canonical factorization $\cX\overset{\pi}{\to}X \overset{g}{\to} Y$. The functor $g_*\colon \operatorname{QCoh}(X) \to \operatorname{QCoh}(Y)$ is an equivalence of categories, because $g$ is affine. We deduce that $\pi_*$ is exact. We conclude by observing that by construction $\pi_*\cO_{\cX}\simeq \cO_X$.
\end{proof}
 \subsection{Relative GIT}\label{subsection_relative_GIT}
 The goal of this subsection is to prove \Cref{prop_stable_locus_of_pi*L^n_otimes_G}, which is the (cohomologically) affine version of \cite{ER21}*{Proposition 3.18}, and is used a few times throughout the paper. For doing so, we begin by recalling the following definition.
 
 \begin{Def}[\cite{Alp}*{\S 11}]
     Let $p:\cX\to S$ be a morphism from an algebraic stack $\cX$ to a scheme $S$, and let $\cL$ be a line bundle on $\cX$. We define the $\cL$-\textit{semistable locus} of $\cX$ over $S$, denoted by $\cX(\cL)^{ss}_S$, as the locus of points $x\in \cX$ such that there exist $U\subset S$ an open neighbourhood of $p(x)$ and a section $s\in H^0(U, p_*(\cL^{\otimes n}))=H^0(p^{-1}U,\cL^{\otimes n})$ for some $n>0$ such that
     \begin{enumerate}
         \item the section $s$ does not vanish at $x$, and
         \item the induced morphism $p^{-1}(U)_s \to U$ is cohomologically affine.
     \end{enumerate} 
 \end{Def}

 \begin{Remark}
     As observed in \cite{ER21}*{Remark 3.16} and the paragraph that follows, the previous definition can be extended to the case in which $S$ is an algebraic stack, as being cohomologically affine commutes with flat descent for morphisms of stacks with quasi-affine diagonal from \cite{Alp}*{Proposition 3.9.(vii)}. More explicitly, for every flat morphism $S'\to S$, inducing $g:\cX':=\cX\times_SS'\to \cX$ we have that $g^{-1}(\cX(\cL)^{ss}_S) = \cX'(g^*\cL)^{ss}_{S'}$.
 \end{Remark}
 \begin{Remark}
     Let $p:\cX\to S$ be a morphism from an algebraic stack $\cX$ to a scheme $S$, and let $\cL$ be a line bundle on $\cX$. Then from \cite{Alp}*{Theorem 11.5} the stack $\cX(\cL)^{ss}_S$ admits a good moduli space over $S$.
 \end{Remark}

\begin{Lemma}\label{lm:character bundle}
    Let $\cX=[\spec(A)/G]$ be an algebraic stack with good moduli space $X=\spec(A^G)$ and $\cL$ a line bundle on $\cX$. Then there exists an affine $G\times \Gm$-scheme $\spec(A')$ over $\spec(A)$ such that
    \begin{enumerate}
        \item there is an isomorphism $\cX\simeq[\spec(A')/G\times \Gm]$, and
        \item the line bundle $\cL$ is isomorphic to $\Bbbk_\chi$ for some character $\chi\colon G\times \Gm\to\Gm$.
    \end{enumerate}
\end{Lemma}
\begin{proof}
    Let $\cY\to \cX$ be the $\Gm$-torsor determined by $\cL$. Then the cartesian product $\cY \times_{\cX} \spec(A)$ is (1) a $\Gm$-torsor over $\spec(A)$ and (2) a $G$-torsor over $\cY$. From (1) we deduce that $\cY \times_{\cX} \spec(A)$ is affine over $\spec(A)$, hence isomorphic to $\spec(A')$ for some $A$-algebra $A'$; from (2) we deduce that $\cY=[\spec(A')/G]$, hence $\cX=[\spec(A')/G\times \Gm]$.
    By construction, given the character $\chi\colon G\times \Gm \to \Gm$ coming from the second projection, we have that $\cL$ is the pull-back of $\Bbbk_{\Id}$.
\end{proof}
Let $\Theta:=[\bA^1/\Gm]$, with geometric points $1$ and $0$.
Recall that given a stack $\cX$ and a line bundle $\cL$ on $\cX$, for any geometric point $p$ in $\cX$ and any morphism $\lambda\colon\Theta\to\cX$ which maps $1\to p$, we can define
\[\mu^{\cL}(p,\lambda)=-\text{weight of }(\lambda^*\cL)|_{[0/\Gm]}.\]
For $\cX=[\spec(A)/G]$ with $G$ reductive, we have \cite{Alp21}*{Proposition 6.9.1}
\[ \Hom(\Theta,\cX) = \left\{ x \in \spec(A)\text{, }\lambda\colon\Gm\to G \text{ such that there exists }\lim_{t\to 0}\lambda(t)\cdot x\right\}.\]
Moreover, for a point $x\in \spec(A)$ which is mapped to $p$, a character $\chi:G\to\Gm$ and $\lambda:\Gm\to G$ a cocharacter such that $\lim_{t\to 0}\lambda(t)\cdot p = y \in \spec(A)$, we have
\[ \mu^\chi(x,\lambda) = \mu^{\Bbbk_\chi}(p,f)\]
where $f(1)=p$ and $f(0)=q$ is the image of $y$ via $\spec(A)\to\cX$. As it is common, we will refer to both these functions as Hilbert-Mumford functions.
\begin{Lemma}\label{lemma_relative_GIT_for_gms_can_be_checked_using_HM}
     Let $\cX$ be an algebraic stack with a good moduli space $X$, a line bundle $\cL$ and $x\in \cX$ a point. Then the following conditions are equivalent:
     \begin{enumerate}
         \item $x\in \cX(\cL)^{ss}_X$, and
         \item for every morphism $\lambda:\Theta\to \cX$ sending $1\mapsto p$, the weight of $\lambda^*\cL$ on $\cB\Gm$ is non-positive.
     \end{enumerate}
     
     Moreover, if $\cX=[\spec(A)/G]$ and $\cL=\Bbbk_\chi$ for a $G$-character $\chi$, the conditions above are equivalent to the existence of $f\in A$ which is $\chi$-semi-invariant and which does not vanish on any point $q\to \spec(A)\to [\spec(A)/G]$ isomorphic to $p$.
 \end{Lemma}
 \begin{proof}
     Since conditions (1) and (2) are \'etale local on $X$, from \cite{luna}*{Theorem 4.12} we can assume that $\cX=[\spec(B)/\Gamma]$ and $X=\spec(B^\Gamma)$. 
     By \Cref{lm:character bundle}, we can assume that $\cL\simeq\Bbbk_\chi$ comes for a $\Gamma$-character $\chi$.
Now, a point $x$ is $\Bbbk_\chi$-semistable in $[\spec(B)/\Gamma]$ if and only if the corresponding point $p\in\spec(B)$ is $\chi$-semistable. The desired statement follows from the Hilbert-Mumford criterion for line bundles linearized by characters on affine rings \cite{Hos}*{Proposition 2.5}. 
 \end{proof}
Recall that, given a reductive group $G$, we can always define a norm $|\cdot |$ on $\Hom(\Gm,G)$ which is invariant under conjugation \cite{Hos}*{\S 2.2}.
 In what follows, when we talk of a \emph{normalized} Hilbert-Mumford function on a quotient stack $[W/G]$ with $G$ reductive, we are implicitly fixing a norm $|\cdot |$ on the lattice of cocharacters of $G$. 
This defines an element $\oH^4(\cB G,\bQ)$, hence an element $b\in \oH^4([W/G],\bQ)$ via pullback. Given $f\colon \Theta \to \cX$, we set $|f|$ to be $f^*b\in \oH^4(\Theta,\bQ)\simeq\bQ$.
 
\begin{Remark}\label{remark_norm_on_maps_from_theta_has_same_info_as_norm_from_1ps}
    If $\lambda\colon \Gm \to G$ such that $\lim_{t\to 0}\lambda(t)\cdot p=q$, the norm of the induced $f\colon\Theta\to [W/G]$ which maps $1\mapsto p$ and $0\mapsto q$ coincides with the previously defined norm of $\lambda$ \cite{HalInst}*{Example 4.1.17}. Moreover, from \cite{HalInst}*{Example 4.1.17}, for each $f\colon\Theta\to [W/G]$ there is a one parameter subgroup $\lambda\colon\Gm\to G$ such that $\lim_{t\to 0}\lambda(t)p = q$ and $\frac{\mu^{\Bbbk_\chi}(p,f)}{|f|} = \frac{\mu^{\chi}(p,\lambda)}{|\lambda|}$.
\end{Remark} 
 \begin{Lemma}\label{lm:bounded}
     Let $G$ be a reductive group and let $\cX=[\spec(A)/G]$ be a quotient stack. Let $\chi\colon G\to \Gm$ be a character.
     Then:
     \begin{enumerate}
         \item given an unstable point $p$, there exists $f_{\min}\colon\Theta\to\cX$ such that
         \[\frac{\mu^{\Bbbk_\chi}(p,f_{\min})}{|f_{\min}|} = \inf_{f(1)\simeq p} \left(\frac{\mu^{\Bbbk_\chi}(p,f)}{|f|}\right);\]
         \item  if we denote by $\cX^c$ the points $p\in \cX$ such that there is $f:\Theta\to \cX$ such that $f(1)\simeq p$ and $f(0)\not \simeq p$, the function \[\cX^c\to \mathbb{R},\quad p\mapsto \inf\left(\frac{\mu^{\Bbbk_\chi}(p,f)}{|f|}\text{ such that } f:\Theta\to \cX, f(1)\simeq p\text{ and }f(1)\not \simeq f(0)\right)\] takes finitely many values
     \end{enumerate}
 \end{Lemma}
 \begin{proof}
     If $G$ is a reductive group acting on an affine scheme $\spec(A)$ and $\chi\colon G\to\Gm$ a character, for any unstable point $x$ there exists a cocharacter $\lambda_{min}:\Gm\to G$ such that the normalized Hilbert-Mumford function
    \[ (x,\lambda) \longmapsto \frac{\mu^{\chi}(x,\lambda)}{|\lambda|} \]
    reaches its minimum value.
    Moreover, the set of minimal values that the normalized Hilbert-Mumford function reaches over the set of points $p\in \spec(A)$ whose $G$-orbit is not closed is finite. These statements follow from \cite{Hos}*{\S 2.2}, especially the discussion after \cite{Hos}*{Lemma 2.13}; we sketch the salient steps. 
    
    First observe that from \Cref{remark_norm_on_maps_from_theta_has_same_info_as_norm_from_1ps}
    it suffices to consider morphisms $\Theta\to \cX$ which come from one parameter subgroups. If $\lambda:\Gm\to G$ is a one parameter subgroup and $g\in G$, then $$\frac{\mu^{\chi}(x,\lambda)}{|\lambda|} = \frac{\mu^{\chi}(gx,g\cdot\lambda\cdot g^{-1})}{|g\cdot\lambda\cdot g^{-1}|},$$ where we denoted by $g\cdot\lambda\cdot g^{-1}$ the one parameter subgroup $\Gm\to G$ that sends $t\mapsto g\lambda(t)g^{-1}$.
        All one parameter subgroups are contained in a maximal torus of $G$ and all maximal tori are conjugated, so it suffices to check the desired statement for one parameter subgroups contained in a fixed maximal torus of $T$ of $G$. We can $T$-equivariantly embed $\spec(A)$ in $\bA^n$, so it suffices to check the desired statement for a $T$-action on $\bA^n$, which we can further assume to be diagonal as each $T$-representation of $T$ is diagonalizable. In this case, one can check the desired statement directly.

    
    To conclude, observe that for any $\Bbbk_\chi$-unstable point $p$ in $\cX$ and any point $x$ of $\spec(A)$ mapping to $p$, and any cocharacter $\lambda$ inducing $f\colon\Theta\to \cX$, we have $\mu^{\chi}(x,\lambda)=\mu^{\Bbbk_\chi}(p,f)$ and $|\lambda|=|f|$.
 \end{proof}


\begin{Prop}\label{prop_stable_locus_of_pi*L^n_otimes_G}
    For $i=1,2$, let $\cX_i$ be an algebraic stack with the same good moduli space $X$. Suppose that there is a morphism $\pi:\cX_2\to \cX_1$ inducing an isomorphism on good moduli spaces and such that the natural map $\cO_{\cX_1}\to \pi_*\cO_{\cX_2}$ is an isomorphism. Let $\cL$ be a line bundle on $\cX_1$ and let $\cM$ be a line bundle on $\cX_2$. Set $\cU_1:=\cX_1(\cL)^{ss}_{X}$ and denote $\cU_1\to U_1$ its good moduli space. Let $\cU_2:=(\pi^{-1}(\cU_1))(\cM)^{ss}_{U_1}$.
    
    Then, for $m\gg 0$, there is an equality $\cU_2 = \cX_2(\pi^*\cL^{\otimes m}\otimes \cM)^{ss}_{X}$. In other terms, the locus in $\cX_2$ which maps to $\cU_1$ and is $\cM$-semistable over $U_1$, agrees with the $\pi^*\cL^{\otimes m}\otimes \cM$-semistable locus of $\cX_2$ over $X$. 
\end{Prop}
\Cref{prop_stable_locus_of_pi*L^n_otimes_G} is the affine version of \cite{ER21}*{Proposition 3.18}, with a similar proof. The only exception is Step 2, which adapts some of the arguments in \cite{reichstein1989stability}*{Proof of Theorem 2.1}.
\begin{proof}
We will use several times the fact that the map $\pi$ is cohomologically affine, since it induces an isomorphism on good moduli spaces \cite{Alp}*{Proposition 3.13}.

\textbf{Step 1.} We show that the desired statement follows from the case $\cX_2=\cX_1$.

First observe that $\xi:\pi^{-1}(\cX_1(\cL)^{ss}_X)\to \cX_1(\cL)^{ss}_X\to U_1$ is still a good moduli space as it is cohomologically affine since it is a composition of cohomologically affine morphisms, and $\cO_{U_1}\to \xi_*\cO_{\pi^{-1}
(\cX_1(\cL)^{ss}_X)}$ is an isomorphism.  It suffices to check that \begin{equation}\label{eq_2}\cX_2(\pi^*\cL)^{ss}_X = \pi^{-1}(\cX_1(\cL)^{ss}_X).\end{equation} Indeed, if this is the case, and if we know that the desired statement holds in the case $\cX_2=\cX_1$, then \[
\cU_2 = (\pi^{-1}(\cX_1(\cL)^{ss}_X))(\cM)^{ss}_{U_1} = (\cX_2(\pi^*\cL)^{ss}_X)(\cM)^{ss}_{U_1} =\cX_2((\pi^*\cL)^{\otimes m}\otimes\cM)^{ss}_X
\]
where the first equality is the definition of $\cU_2$, the second one is \Cref{eq_2} together with the fact that $\xi$ is still a good moduli space, the third one is the desired statement in the case $\cX_1=\cX_2$.

It suffices that we show the two inclusions in \Cref{eq_2} set theoretically, since both stacks are open substacks of $\cX_2$. If $x_1\in \cX_1(\cL)^{ss}_X$ then, if $p:\cX_1\to X$ is the good moduli space, there is $W\subseteq X$ an open subscheme and $s\in p_*(\cL^{\otimes n})(W)$ such that $s$ does not vanish at $x_1$ and the locus $p^{-1}(W)_s\subseteq \cX_1$ where $s$ does not vanish is cohomologically affine over $W$. Since $\cX_2\to\cX_1$ is cohomologically affine, also \[(p\circ\pi)^{-1}(W)_{\pi^*s} = \pi^{-1}(p^{-1}(W)_s)\] is cohomologically affine over $W$, and $\pi^*s$ is also a section of $\pi^*\cL$. Since $\pi^*s$ does not vanish along $\pi^{-1}(x_1)$ we have that $\pi^{-1}(x_1)\subseteq \cX_2(\pi^*\cL)^{ss}_X$.

For the other inclusion, let $x_2\in \cX_2(\pi^*\cL)^{ss}_X$. Then there is $W$ an open neighbourhood of $p(\pi(x_2))$ and a section $\widetilde{s}\in p_*\pi_*\pi^*\cL(W)$ such that $(p\circ\pi)^{-1}(W)_{\widetilde{s}}\to W$ is cohomologically affine, and $\widetilde{s}$ does not vanish at $x_2$. Observe that, since $\pi_*\cO_{\cX_2}=\cO_{\cX_1}$, we have that
$\widetilde{s}=\pi^*s$ for a section $s\in p_*\cL(W)$, and \[(p\circ\pi)^{-1}(W)_{\widetilde{s}}=(p\circ\pi)^{-1}(W)_{\pi^*s} = \pi^{-1}(p^{-1}(W)_s)\]
where the first equality follows since $\widetilde{s}=\pi^*s$ and the second equality follows by definition of non-vanishing locus.
We need to show that $p^{-1}(W)_s\to W$ is cohomologically affine. So let $\cF\to \cG$ be a surjective morphism of quasi-coherent sheaves on $p^{-1}(W)_s$, we need to show that $p_*\cF\to p_*\cG$ remains surjective. For doing so, it suffices to check that the natural map \[\cH\to\pi_*\pi^*\cH\] is an isomorphism for any quasi-coherent sheaf $\cH$. Indeed, if this is the case then $\pi^*\cF\to \pi^*\cG$ will be surjective, which implies that $p_*\pi_*\pi^*\cF\to p_*\pi_*\pi^*\cG$ will be surjective since $p\circ\pi:(p\circ\pi^{-1})(W_s)\to W$ is cohomologically affine, and it will agree with $p_*\cF\to p_*\cG$. 

Consider a smooth atlas $Z\to p^{-1}(W)_s$ from an affine scheme $Z$. If we denote by \[\cY:=\pi^{-1}(p^{-1}(W)_s)\times_{p^{-1}(W)_s}Z,\] then the second projection $\mu:\cY\to Z$ is a good moduli space. Indeed, it is cohomologically affine since $\cX_2\to \cX_1$ is cohomologically affine and being cohomologically affine commutes with base change for a quasi-affine morphism \cite{Alp}*{Proposition 3.10.(vi)}. The morphism $Z\to p^{-1}(W)_s$ is quasi-affine (in fact, affine) as $Z$ is affine and $p^{-1}(W)_s$ has affine diagonal. Moreover $\mu_*\cO_\cY=\cO_{Z}$ since $\pi_*\cO_{\cX_2}=\cO_{\cX_1}$. Then \cite{Alp}*{Proposition 4.5} applies. This finishes Step 1. 
\color{black}

Therefore we can assume that $\cX_1=\cX_2$ and we will adopt the following notation:
\begin{itemize}
    \item we will denote $\cX_1=\cX_2$ by $\cX$, with the good moduli space morphism $p:\cX\to X$,
    \item we will not write $\pi$ any more as $\pi=\Id$ when $\cX_1=\cX_2$.
\end{itemize}


Let $\cV_m:=\cX(\cL^{\otimes m}\otimes \cM)^{ss}_{X}$.  The proof is in three steps, but first we perform some reductions. To start, observe that our claim is \'etale local on the good moduli space $X$, hence:
\begin{enumerate}
    \item we can assume that $\cX=[\spec(A)/G]$ from \cite{AHREtale};
    \item we can assume that the line bundles $\cL$ (respectively $\cM$) come from a character $\theta$ of $G$ (a character $\chi$ of $G$) by applying \Cref{lm:character bundle} twice;
    \item semistability with respect to $\cL$ and $\cM$ can be checked using the Hilbert-Mumford criterion by \Cref{lemma_relative_GIT_for_gms_can_be_checked_using_HM}; the same conclusion holds for $\cG_m := \cL^{\otimes m}\otimes \cM$;
    \item by \Cref{lm:bounded}, for $\cN=\cL, \cM$, we can assume that for every $\cN$-unstable point $q$ there exists $\lambda_{\min}$ which minimizes $\mu^{\cN}(q,\lambda)/|\lambda|$; furthermore, the set of such minimal values is finite.
\end{enumerate}
\textbf{Step 2.} We prove that there is $m_0$ such that, for every $m\ge m_0$, if $q\in \cX\smallsetminus \cU_1$, then $q\notin \cV_m$. 
For this, it is enough to prove that there exists an $m_0$ such that, for every $m\geq m_0$ and $q\notin \cU_1$, there exists a $\lambda\colon\Theta\to\cX$ which maps $1\mapsto q$ and $\lambda(0)\not\simeq q$ such that  $\mu^{\cG_m}(q,\lambda) < 0$. 

The idea is to factor $\mu^{\cG_m}(q,\lambda) = \mu^{\cL^{\otimes m}\otimes \cM}(q,\lambda) = m\mu^{\cL}(q,\lambda) + \mu^{ \cM}(q,\lambda)$, pick a $\lambda$ which makes $\mu^{\cL}(q, \lambda) $ negative, and pick $m$ big enough to cancel the contribution of $\mu^{ \cM}(q,\lambda)$. In order to perform this strategy uniformly on $\cX\smallsetminus \cU_1$, we use the normalized Hilbert-Mumford function.

For each point $q\not\in \cU_1$ there is a morphism $\lambda\colon\Theta\to \cX$ such that $\lambda(1)=q\not\simeq\lambda(0)$. Then the results from \Cref{lm:bounded} apply to points in $\cX\smallsetminus\cU_1$ and the normalized Hilbert-Mumford function $\frac{\mu^{\cG_m}(q,\lambda)}{|\lambda|}$. So let
 
\begin{align*}d=\max_{q \notin \cU_1} \left\{ \inf_{\lambda\colon 1 \to q} \frac{\mu^{\cL}(q,\lambda)} {|\lambda|} \right\},\quad e=\max_{q\notin \cU_1} \left\{\inf_{\lambda\colon 1 \to q} \frac{\mu^{\cM^{\vee}}(q,\lambda)}{|\lambda|} \right\}
    =\max_{q\notin \cU_1} \left\{\sup_{\lambda\colon 1 \to q} \frac{\mu^{\cM}(q,\lambda)}{|\lambda|} \right\}.
\end{align*}
Observe that with the reductions that we performed, we have $d<0$ and $e<\infty$.

Let $m_0>0$ be an integer such that $m_0\cdot d + e < 0$. For every $q\not \in \cU_1$, let $\lambda_{\min}\colon\Theta\to \cX$ be a morphism which maps $1\mapsto q$ and such that the normalized Hilbert-Mumford function in $q$ reaches its minimum. So
\begin{align*}
    \frac{\mu^{\cG_m}(q,\lambda_{\min})}{|\lambda_{\min}|} = \frac{\mu^{\cL^{\otimes m}}(q,\lambda_{\min)}}{|\lambda_{\min}|} + \frac{\mu^{\cM}(q,\lambda)}{|\lambda|} \le md + e <0
\end{align*}
for every $m\geq m_0$. This shows that $q\notin \cV_m$.
From now on we assume $m\geq m_0$, and we use $\cV$ to denote $\cV_m$.

    \noindent
    \textbf{Step 3.} We show that if $q\in \cU_1\smallsetminus\cU_2$ then $q\notin \cV$.

    Indeed, as $\cU_1\to \cU_1\to U_1$ is cohomologically affine and the line bundle is a character of a group, we can check if a point in $\cU_1$ lies in $\cU_2$ using the Hilbert-Mumford criterion. In particular, if $q\in \cU_1\smallsetminus\cU_2$ there is a morphism $\lambda:\Theta\to \cU_1$ sending $1\mapsto q$ such that $\mu^{\cM}(q,\lambda)<0$. Since $\lambda(\Theta)\subseteq \cU_1$, there is an invariant section of $\cL$ not vanishing on $\lambda(\Theta)$. As for $\lambda(\Theta)$ to be in $\cU_1$ the point $q$ cannot be strictly stable, we have $\mu^{{\cL}}(q,\lambda)=0$. Now the desired statement follows by computing the Hilbert-Mumford function: 
    \[\mu^{\cL^{\otimes m}\otimes \cM}(q,\lambda) = m\mu^{\cL}(q,\pi\circ\lambda) + \mu^{\cM}(q,\lambda) = \mu^{\cM}(q,\lambda)<0.\]

    \noindent
    \textbf{Step 4.} We prove that if a point $q\in \cU_2$, then it is also in $\cV$. This concludes the proof.
    
    We need to produce a section of $\cL^{\otimes m}\otimes \cM$ which does not vanish at $q$. Let $p:\cX\to X$ be the good moduli space morphism. 
    Since $q\in \cU_2\subset\cU_1$, up to shrinking the base, there is a section $s$ of $p_*\cL$ which does not vanish at $q$. As $U_1=\Proj(\bigoplus_n p_*(\cL^{\otimes n}))$, there is an open affine subset of the form $U_{1,s}$ such that $q$ maps to $U_{1,s}$. 
    As $\cU_2$ is the relative stable locus for $\cM=\Bbbk_\chi$,
    and since $\cX$ is a global quotient, from \Cref{lemma_relative_GIT_for_gms_can_be_checked_using_HM} there is $f\in \oH^0(\cX_{s},\cM)$ such that $f(q)\neq 0$.
    From \cite{Liu_AG_book}*{\S 5.1.25}, which is stated for schemes but holds for algebraic stacks as well, there is a section of $\cL^{\otimes m}\otimes\cM$ on the whole $\cX$ lifting $f\otimes s^n$. This section will not vanish at $q$, as desired. 
 \end{proof}
 The following lemma is the standard tool for passing from projective GIT to affine GIT. \begin{Lemma}\label{lemma_git_for_ample_linearization_is_contained_in_the_affine_git_for_the_cone}
    Let $W$ be a quasiprojective scheme with an action of $G$, and $L$ a $G$-linearized ample line bundle. For $d>0$, consider the character \[\chi:G\times \Gm\to \Gm, \text{ }(g,t)\mapsto t^{-d}.\]
    This induces a $G\times \Gm$-equivariant line bundle $\Bbbk_\chi$, and up to replacing $L$ with $L^{\otimes m}$ for $m>0$,
    \[
    [W(L)^{ss}/G] \subseteq [\spec(\bigoplus_{i\ge 0} \oH^0(L^{\otimes i}))(\Bbbk_{\chi})^{ss}/G\times \Gm].
    \]
    Moreover, if $W$ is either affine or projective, the inclusion above is an isomorphism, and if $W=\spec(A)$ is affine we also have 
    \[
    [\spec(A)(L)^{ss}/G] \cong [\spec(\bigoplus_{i\in \bZ} \oH^0(L^{\otimes i}))(\Bbbk_{\chi})^{ss}/G\times \Gm].
    \]
\end{Lemma}
\begin{proof} As the semistable locus of $\Bbbk_\chi$ agrees with the one of $\Bbbk_\chi^{\otimes k}$ for every $k\ge 1$, we assume $d=1$.
    Without loss of generality, we can assume that $L$ is generated in degree one and very ample, and also the ring $R_L=\oplus_{n\geq 0} H^0(W,L^{\otimes n})^G$ is generated in degree one.
    
    Let $\pi:P_L\to W$ be the $\Gm$-torsor associated to the $G$-linearized line bundle $L$ on $W$, so that there is a $G$-action on $P$ for which $\pi$ is $G$-equivariant. Let $\Bbbk_\chi$ be the $\Gm$-linearized line bundle on $P_L$ associated to the character $\chi$. Then $\pi^*L=\Bbbk_\chi$, and we can regard the latter as a $G\times\Gm$-linearized line bundle. 
    
    We claim that $P_L(\Bbbk_\chi)^{ss}=\pi^{-1}(W(L)^{ss})$, where the first semistable subset is taken with respect to the $G\times \Gm$-action, and the second semistable subset with respect to the $G$-action. This is equivalent to saying that $[W(L)^{ss}/G]\simeq [P_L(\Bbbk_\chi)/G\times \Gm]$.

    The claim follows from the equality $\pi^*\colon H^0(W,L)^{G} \simeq H^0(P_L,\Bbbk_\chi)^{G\times \Gm}$, which follows from the fact that $\pi^*\colon H^0(W,L) \simeq H^0(P_L,\Bbbk_\chi)^{\Gm}$ is an isomorphism of $G$-representations.

    The same argument applies to $\Proj(R_L)$ with the line bundle $\cO(1)$, where
    \[ R_L:= \oplus_{n\geq 0} H^0(W,L^{\otimes n}).\]
    In this particular case the $\Gm$-torsor over $\Proj(R_L)$ is isomorphic to $\spec(R_L)\smallsetminus\{0\}$. From \cite{stacks-project}*{\href{https://stacks.math.columbia.edu/tag/01PZ}{Tag 01PZ}} and \cite{stacks-project}*{\href{https://stacks.math.columbia.edu/tag/01PZ}{Tag 01Q1}}, the induced morphism $\iota:W\to \Proj(R_L)$ is an open embedding with $\iota^*\cO(1)\simeq L$ and we have a cartesian diagram
    \[
    \begin{tikzcd}
        P_L \ar[r, "j"] \ar[d, "\pi"] & \spec(R_L)\smallsetminus\{0\} \ar[d, "\rho"] \\
        W \ar[r, "\iota"] & \Proj(R_L)
    \end{tikzcd}
    \]
    with $\rho^*\cO(1)\simeq \Bbbk_\chi$. By cartesianity, also $j$ is an open embedding.

    We claim that any point $x$ in $P_L(\Bbbk_\chi)^{ss}$ is also $\Bbbk_\chi$-semistable in $\spec(R_L)$. For this, we have to study the restriction homomorphism
    \[ j^*\colon H^0(\spec(R_L),\Bbbk_\chi)^{G\times\Gm} \longrightarrow H^0(P_L,\Bbbk_\chi)^{G\times \Gm}.\]
    Observe that
    \[
        H^0(\spec(R_L),\cO)\simeq R_L = \oplus_{n\geq 0} H^0(W,L^{\otimes n}),
    \]
    hence the composition
    \begin{align*}
        H^0(W,L) \hookrightarrow H^0(\spec(R_L)\smallsetminus\{0\},\Bbbk_\chi)^{G\times \Gm} \overset{j^*}{\longrightarrow} H^0(P_L,\Bbbk_\chi)^{G\times\Gm} \simeq H^0(W,L)^G
    \end{align*}
    is necessarily surjective. From this it follows that if there is a $G\times \Gm$-invariant section of $\Bbbk_\chi$ over $P_L$ which does not vanish in $x$, then there is also a $\Gm\times G$-invariant section of $\Bbbk_\chi$ over $\spec(R_L)\smallsetminus\{0\}$ which does not vanish over $x$. This proves that
    \[
    [W(L)^{ss} /G] \simeq [P_L(\Bbbk_\chi)^{ss} / G\times\Gm ] \overset{\text{open}}{\hookrightarrow} [\spec(R_L)(\Bbbk_\chi)^{ss}/G\times\Gm],
    \]
    since the origin is never semistable for this linearization.

    If $W$ is projective or affine, the open embedding $\iota$ is an isomorphism, which implies that $P_L\simeq \spec(R_L)\smallsetminus\{0\}$, hence
    \[ [W(L)^{ss}/G] \simeq [\spec(R_L)(\Bbbk_\chi)^{ss}/G\times\Gm].\]
    If $W$ is affine, we have that also $P_L$ is affine, hence $P_L\simeq \spec(\cO_{P_L}(P_L))$. By definition $P_L$ is the relative spectrum of the $\cO_W$-algebra $\oplus_{n\in \bZ} L^{\otimes n}$, hence
    \[ H^0(P_L,\cO_{P_L}) \simeq H^0(W, \oplus_{n\in \bZ} L^{\otimes n}) \simeq \oplus_{n\in \bZ} H^0(W,L^{\otimes n}),\]
    from which we deduce $[P_L(\Bbbk_\chi)^{ss}/G]\simeq [\spec(\oplus_{n\in \bZ} H^0(W,L^{\otimes n})(\Bbbk_\chi)^{ss})/G\times\Gm]$.
\end{proof}
We end this subsection giving a few conditions under which an open substack of an algebraic stack is the semistable locus of a line bundle.
\begin{Lemma}\label{lemma_condition_when_X_dm_is_the_ss_locus_for_a_vb_when_X_is_normal}
    Let $\cX$ be an algebraic stack with a good moduli space $\cX\to X$, and let $\cX'\subseteq \cX$ be a dense open substack that admits a good moduli space $X'$ projective over $X$. Let $\cL$ be a line bundle on $\cX$ which, once restricted to $\cX'$, descends to a line bundle $L_{X'}$ on $X'$, that is ample over $X$. Assume finally that the complement of $\cX'$ in $\cX$ has codimension at least 2, and that $\cX$ is normal. Then $\cX'=\cX(\cL)^{ss}_X$.
\end{Lemma}
\begin{proof}The statement is \'etale local over $X$ so we can assume that $\cX=[\spec(A)/G]$ and $X=\spec(A^G)$, so $L'$ is ample.
We will use the following chain of equalities: \begin{equation}\label{equality_sections}
    \oH^0(\cX,\cL^{\otimes m}) = \oH^0(\cX',\cL^{\otimes m}|_{\cX'}) = \oH^0(X',L_{X'}^{\otimes m})
\end{equation} where the first equality follows since $\cX$ is normal and the complement of $\cX'$ has codimension at least 2.

    The inclusion $\cX'\subseteq \cX(\cL)^{ss}_X$ can be proved as follows. For every $p\in \cX'$ there is an $m\gg 0$ and a section of $L_{X'}^{\otimes m}$ which does not vanish at $p$. On the other hand from \Cref{equality_sections}, such a section extends to a section $s\in\oH^0(\cX,\cL^{\otimes m})$.
    The locus where $s$ does not vanish is of the form $[\spec(A)\smallsetminus V(\pi^*s)/G]$ where $\pi:\spec(A)\to \cX$ is the quotient map. Since $\spec(A)\smallsetminus V(\pi^*s)$ is affine, also $[\spec(A)\smallsetminus V(\pi^*s)/G]$ is cohomologically affine. This proves $\cX_\dm\subseteq \cX(\cL)^{ss}_X$.

    From \cite{Alp}*{Theorem 11.5}, there is a morphism $\xi:\cX(\cL)^{ss}_X\to \Proj(\bigoplus \oH^0(\cX,\cL^{\otimes m}))$, and from \Cref{equality_sections} and since $L_{X'}$ is ample, $\Proj(\bigoplus \oH^0(\cX,\cL^{\otimes m}))= X'$. In particular, as $\cX'\to X'$ is surjective, also $\cX(\cL)^{ss}_X\to X'$ is surjective, so $\xi$ is a good moduli space from \cite{Alp}*{Theorem 11.5}. Both $\cX'\to X'$ and $\cX(\cL)^{ss}_X\to X'$ are cohomologically affine, the stack $\cX(\cL)^{ss}$ is normal and the complement of $\cX'$ in $\cX(\cL)^{ss}_X$ is at least 2. So $\cX'=\cX(\cL)^{ss}_X$ as desired.
\end{proof}
\begin{Cor}\label{cor_when_X_dm_is_ss_locus_for_a_lb}
    Let $\cX$ and $\cX'$ be as in \Cref{lemma_condition_when_X_dm_is_the_ss_locus_for_a_vb_when_X_is_normal}, except that $\cX'$ is assumed to be Deligne-Mumford and $\cX$ is not assumed to be normal. If $\cX(\cL)^{ss}_X$ is dense, we have $\cX'=\cX(\cL)^{ss}_X$.
\end{Cor}
Observe that if the properly stable locus $\{x\in X:\pi^{-1}(x)$ is Deligne-Mumford$\}$ is dense then $\cX(\cL)^{ss}_X$ is dense, where $\pi:\cX\to X$ is the good moduli space map.
\begin{proof}The statement is  \'etale local over $X$ so we can assume that $\cX=[\spec(A)/G]$ and $X=\spec(A^G)$, so $L'$ is ample. Observe that $\cX(\cL)^{ss}_X\subseteq \cX'$. Indeed,
consider $\nu:\cX^n\to \cX$ be the normalization. Then $\nu^{-1}(\cX(\cL)^{ss}_X)\subseteq \cX^n(\nu^*\cL)^{ss}_X$ since $\nu $ is affine, and $\cX^n(\nu^*\cL)^{ss}_X = (\cX')^n = \nu^{-1}(\cX')$ where the first equality follows from \Cref{lemma_condition_when_X_dm_is_the_ss_locus_for_a_vb_when_X_is_normal} (the assumpions of \Cref{lemma_condition_when_X_dm_is_the_ss_locus_for_a_vb_when_X_is_normal} are still satisfied as $\nu$ is the normalization). Then $\cX(\cL)^{ss}_X =\nu(\nu^{-1}(\cX(\cL)^{ss}_X))\subseteq \nu(\nu^{-1}(\cX'))=\cX'$. So $\cX(\cL)^{ss}_X$ is an open substack of $\cX'$, so it is Deligne-Mumford, and from \Cref{lemma_git_for_ample_linearization_is_contained_in_the_affine_git_for_the_cone} the coarse moduli space of $\cX(\cL)^{ss}_X$ is projective, hence $\cX(\cL)^{ss}_X=\cX'$.
\end{proof}

 \subsection{Luna slice}\label{subsection_luna_slice}
 For the convenience of the reader, we briefly recall some classical results on GIT, which are well-known.
    Let $X$ be a scheme of finite type with an action of a reductive group $G$, and let $U\subseteq X$ be the semistable locus with respect to a $G$-equivariant line bundle $\cL$ on $X$. If $x\in U$ is a closed point with stabilizer $G_x$, we have a cartesian diagram
    \[
    \xymatrix{
    [\spec(A)/G_x]\ar[r]\ar[d] & [U/G]\ar[d] \\ 
    \spec(A^{G_x})\ar[r]^-\alpha  & U/\!\!/_{\cL} G,
    }
    \]
    where $\alpha$ is an \'{e}tale neighborhood of the image of $x$ in the good moduli space \cite{Alp21}*{Corollary 6.7.3}. If $x$ is a smooth point with tangent space $T_x$, there is an action of $G_x$ on $T_x$. 
    Assume that the orbit of $x$ is closed: then from Matsushima’s Theorem the group $G_x$ is reductive, so we can decompose the $G_x$-representation \[T_x=T_x(G\cdot x)\oplus N_x\] where $T_x(G\cdot x)$ is the tangent space at $x$ of the orbit $Gx$ of $x$, and $N_x$ is $G_x$-invariant. Then it easily follows from  \cite{Alp21}*{Theorem 6.7.5, "Luna’s \'{E}tale Slice Theorem"} applied to $[\spec(A)/G_x]$ that we have a cartesian diagram
    \[
    \xymatrix{[N_x/G_x]\ar[d] & [\spec(A)/G_x]\ar[l]\ar[d] \\
    N_x/\!\!/G_x & \spec(A^{G_x}) \ar[l]_-\beta ,}
    \]
    where $\beta$ is an \'{e}tale neighborhood of $0\in N_x$
\begin{Def}
    With the notations above, we define $[N_x/G_x]$ to be a \textit{Luna slice} for $[U/G]$ at $x$.
\end{Def}
In other terms, the map $[N_x/G_n]\to N_x/\!\!/G_x$ gives a local description of $\pi:[U/G]\to U/\!\!/_\cL G$ in a neighbourhood of $\pi(x)$.

\section{Stable quasimaps}\label{sec:qmaps}
The goal of this section is to extend the existing theory of orbifold quasimaps to algebraic stacks $\cX$ with a projective good moduli space, and containing a proper Deligne-Mumford stack $\cX_\dm\subseteq \cX$. 
This section is organized as follows.

In \S \ref{subsection_def_stable_qmaps} we extend the definition of stable quasimap $\sC\to \cX$, given in \cite{ciocan2014stable}, to arbitrary stacks $\cX$ with a projective good moduli space, and containing a proper Deligne-Mumford stack. Our definition of stable quasimap is the combination of the definition of quasimap given in \cite{ciocan2014stable} with the definition of twisted map of \cite{dli1}.
We show that the moduli stack of stable quasimaps $\cQ_{g,n}(\cX,\cX_\dm)$ is locally of finite type. 

In \S \ref{section_qmaps_satisfy_the_val_crit_for_prop} we show that $\cQ_{g,n}(\cX,\cX_\dm)$ satisfies the valuative criterion for properness. 

In \S \ref{section_boundedness_qmaps} we show that if one fixes the class of a quasimap, defined as in \cite{ciocan2014stable}, the locus of quasimaps of fixed class $\beta$ in $\cQ_{g,n}(\cX,\cX_\dm)$ is bounded.
For this section, we need the additional \Cref{assumptions:extension of line bundle}, namely that $\cX$ is a global quotient and that $\cX_\dm$ is the semistable locus of a line bundle on $\cX$ (observe however that the last condition is often automatic, see \Cref{cor_when_X_dm_is_ss_locus_for_a_lb}).

In \S \ref{subsection_obstruction_theory} we show that $\cQ_{g,n}(\cX,\cX_\dm,\beta)$ carries a perfect obstruction theory for $\cX_\dm$ smooth and $\cX$ lci.

\subsection{Quasimaps and stable quasimaps}\label{subsection_def_stable_qmaps}
Throughout this subsection we will adopt the following 
\begin{Notation}\label{notation_cX_and_cX_dm}
    We will denote by $\cX$ an algebraic stack with a quasi-projective good moduli space $\cX\to X$. Further, we will assume that $\cX$ contains an open dense substack $\cX_\dm\subseteq \cX$ such that $\cX_\dm\to X$ is proper and Deligne-Mumford, with coarse moduli space $\cX_\dm\to X_\dm$ which is projective over $X$. We will denote by $L_X$ an ample line bundle on $X$ and $L_{X_\dm}$ an ample line bundle on $X_\dm$.  
\end{Notation}
The following definition already appeared \cite{orbifold_qmap_theory}, albeit only for quotient stacks $\cX=[W/G]$ with $W$ affine.
\begin{Def}\label{def:quasimap}
    An $n$-marked \emph{quasimap} of genus $g$ over $S$ is $(\phi:\sC\to \cX, \Sigma_1,\ldots, \Sigma_n) $ such that
    \begin{enumerate}
        \item the $\Sigma_i \subset \sC \to S$ are $n$ distinct gerbes over $S$ for $i=1,\ldots,n$, and $(\sC \to S,\Sigma_1,\ldots,\Sigma_n)$ is a twisted, $n$-marked curve of genus $g$ as in \cite{AV_compactifying};
        \item the morphism $\phi:\sC\to \cX$ is representable and maps the generic points, the marked gerbes and the nodal gerbes of $\sC$ to $\cX_\dm$.
    \end{enumerate}
    If $S=\spec(k)$, we will just say an $n$-marked quasimap of genus $g$.
\end{Def}   
We will denote the coarse moduli space of a twisted curve $\sC$ by $C$, and the image of the $\Sigma_i$ in $C$ by $p_i$.
The following definition is the intersection of \cite{dli1}*{Definition 2.2} and \Cref{def:quasimap}:

\begin{Def}\label{def_stable_qmap}
    Let $(\phi:\sC\to \cX, \Sigma_1,\ldots, \Sigma_n) $ be $n$-marked quasimap of genus $g$, over $\spec(\ell)$, with $\ell$ an algebraically closed field. It is \emph{stable} if:
    \begin{enumerate}
        \item the line bundle $\omega_C(\sum p_i)\otimes f^*L_X^{\otimes 3}$ is nef, where $f:C\to X$ is the morphism induced by $\phi$ on good moduli spaces, and
        \item if $\sD\subseteq \sC$ is an irreducible component with coarse space $D$ such that $\omega_C(\sum p_i)\otimes f^*L_X^{\otimes 3}$ has degree zero on $D$, then then $D\cong \bP^1$ and there are two geometric points $p_1,p_2\in \sD$ such that $\phi(p_1)$ and $\phi(p_2)$ are not isomorphic in $\cX$.
    \end{enumerate}
\end{Def}
\begin{Remark}
    {For a quasimap $\phi:\sC\to \cX$, condition (2) in \Cref{def_stable_qmap} is equivalent to $\phi$ being a \textit{twisted map} (see \cite{dli1}*{Definition 2.2}). In other terms, if $\sD$ is obtained by performing a root stack of $\bP^1$ at 0 with order $m_0$ and $\infty$ with order $m_\infty$ (where we allow $m_0$ and $m_\infty$ to be 1), and $\phi:\sD\to \cX$ is a quasimap, then for every pair of points $p_1,p_2\in \sD$ we have that $\phi(p_1)$ and $\phi(p_2)$ map to the same point of $\cX$ if and only if there is $d\in \bZ_{>0}$ such that the map $\sD\to\cX$ factors as $\sD\to \cB\bmu_d\to\cX$. One direction is clear, the other follows from \cite{dli1}*{Lemma 2.5} as in the proof of \cite{dli1}*{Proposition 2.6}.}
\end{Remark}

\begin{Remark}\label{remark_passing_to_a_subcurve_preserves_stability} Observe that being stable is preserved by taking marked subcurves. More precisely, let $(\phi:\sC\to \cX, \Sigma_1,\ldots, \Sigma_n) $ be an $n$-marked quasimap, and let $\sD\subseteq \sC$ be the closed substack given by the union of some irreducible components of $\sC$, attached along the nodes $n_1,\ldots, n_k\subseteq \sD$ which could be gerbes. If $\phi$ is stable, then $(\phi|_\sD:\sD\to \cX, (\Sigma_1)|_\sD,\ldots, (\Sigma_n)|_\sD, n_1,\ldots,n_k)$ is also stable.
\end{Remark}
If $\cX=[\spec(A)/G]$ as in \cites{ciocan2014stable, orbifold_qmap_theory} and $g\neq 1$, this recovers the usual definition of orbifold stable quasimap.
\begin{Lemma}\label{lemma_equivalent_conditions_for_stability}
    Let $(\phi:\sC\to \cX, \Sigma_1,\ldots, \Sigma_n) $ be a quasimap of genus $g\neq 1$, satisfying condition (1) in \Cref{def_stable_qmap}. The following conditions are equivalent:
    \begin{enumerate}
        \item $\phi$ is stable,
        \item for every irreducible component $\sD\subseteq \sC$, with coarse space $D$, such that $\omega_C(\sum p_i)\otimes f^*L_X^{\otimes 3}$ has degree zero on $D$, either $\phi(\sD)\not \subset \cX_{\dm}$, or $L_{X_\dm}$ is ample when pulled back to $D$.    \end{enumerate} 
        Moreover, assume that $\cX=[X/G]$ is a quotient stack and there is $\cL$ a line bundle on $[X/G]$ such that $[X(\cL)^{ss}/G]=\cX_\dm$. If we denote by $\widetilde{L}$ the line bundle on $C$ obtained by descending a high enough power of $\phi^*\cL$, condition (2) is equivalent to $\omega_C(\sum p_i)\otimes f^*L_X^{\otimes 3}\otimes \widetilde{L}^{\otimes \epsilon}$ being ample for $0<\epsilon \ll 1$.
\end{Lemma}
In particular, for the moreover part, if $X=\spec(A)$ is affine, then $L_X$ is trivial as $X/\!\!/G=\spec(A^G)$ is affine so $\sC\to\spec(A^G)$ factors as $\sC\to \spec(k)\to \spec(A^G)$. Then condition (2) is equivalent to the stability conditions given in \cites{ciocan2014stable,orbifold_qmap_theory}.
\begin{proof}
First we prove that (1) and (2) are equivalent. Let $\sD\subseteq \sC$ be as in (2), with attaching nodes $n_1,n_2$ which, as usual, could be gerbes, and let's first assume that $\phi(\sD)\subseteq \cX_\dm$. Then, from \cite{dli1}*{Proposition 2.6} the map $(\phi|_\sD\to \cX, n_1,n_2)$ is stable if and only if it is a twisted \textit{stable} map in the sense of \cite{AV_compactifying}, if and only if $L_{X_\dm}$ is ample on $D$. 
If instead $\phi(\sD)\not \subset \cX_\dm$, then not all points of $\sC$ map to isomorphic points on $\cX$: from the quasimap condition the generic locus of $\sD$ maps to $\cX_\dm$, but $\phi(\sD)\not \subset \cX_\dm$ so there must be a point $x\in \sC$ not mapping to $\cX_\dm$.

%We are left to show that if $\phi(\sD)\not \subset \cX_\dm$, then there cannot be a factorization $\sD\to \cB \bmu_d\to \cX$. If $\sD\to \cX$ factored as $\sD\to \cB \bmu_d\to \cX$, then for each pair of $k$-points $p_1,p_2\to \sD$, the compositions $\phi(p_1)$ and $\phi(p_2)$ would be isomorphic, as they are isomorphic in $\cB\bmu_d$. Since $\phi$ is a quasimap, the set of $k$ points of $\sD$ which map to $\cX_\dm$ is not empty, so all the points would map to $\cX_\dm$, which contradicts $\phi(\sD)\not \subset \cX_\dm$. 


We now prove the moreover part, first assuming (2). It suffices to prove that if $\sD$ and $D$ are is as in (2), then $\widetilde{L}$ has positive degree on $D$. 
By hypothesis, $\phi(\sD)\cap [X(\cL)^{ss}/G]$ is non-empty. By definition of semistability there is a section $s$ of $\cL$ whose pullback to $\sD$ is not the zero section. This implies that $\phi^*\cL|_\sD$ has at least a non-zero global section $s$. If $\phi(\sD)\not\subset \cX_\dm$, then by definition of semistability $s$ must vanish in at least one point, so $\phi^*\cL$ has a section which is not zero and not constant, so $\deg(\phi^*\cL|_\sD)>0$. If instead $\phi(\sD)\subset\cX_\dm$, since $L_{X_\dm}$ is ample on $D$, the map $D\to X_\dm$ is finite, so again $\widetilde{L}$ has positive degree on $D$.

The other direction is similar. Indeed, if $\omega_C(\sum p_i)\otimes f^*L_X^{\otimes 3}\otimes \widetilde{L}^{\otimes \epsilon}$ is ample for $0<\epsilon \ll 1$, then $\omega_C(\sum p_i)\otimes f^*L_X^{\otimes 3}$ is nef, and if $D$ is a component on which $\omega_C(\sum p_i)\otimes f^*L_X^{\otimes 3}$ has degree 0, then $D\cong \bP^1$ since $g\neq 1$. Moreover, the line bundle $\widetilde{L}$ has positive degree on $D$, and $\cL|_{\cX_\dm}$ descends to an ample line bundle on $X_\dm$. So either $\phi(\sD)\not \subset \cX_\dm$, or if $\phi(\sD)\subset \cX_\dm$, then $f(D)$ cannot be a point. So $L_{X_\dm}$ is ample on $D$.
\end{proof}
\begin{Assumptions}
    From now on, we assume that $X$, the good moduli space of $\cX$, is a projective variety.
\end{Assumptions}



\begin{Def}
We define the category $\cQ_{g,n}(\cX,\cX_\dm)$ of stable quasimaps to $\cX$ with respect to $\cX_\dm$, as the following category fibred in groupoids over the big \'{e}tale site of schemes over $k$:
\begin{enumerate}
    \item for any scheme $S$, the objects of $\cQ_{g,n}(\cX,\cX_\dm)(S)$ are the $n$-marked quasimaps to $\cX$ of genus $g$ over $S$ whose geometric fibers are stable, and
    \item the morphisms $(\phi\colon\sC\to\cX,\Sigma) \longrightarrow (\phi'\colon\sC'\to\cX,\Sigma')$ over $S$ consist of an isomorphism $\alpha\colon\sC\simeq \sC'$ such that $\alpha^{-1}(\Sigma'_i)=\Sigma_i$ and an isomorphism $\gamma\colon\phi'\circ\alpha \simeq \phi$.
\end{enumerate}
The category $\mathfrak{Q}_{g,n}(\cX,\cX_\dm)$ of quasimaps to $\cX$ is defined in the same way but without requiring the geometric fibres to be stable.
\end{Def}
\begin{Prop}\label{prop_qmap_gives_open_cond}
    The stack $\mathfrak{Q}_{g,n}(\cX,\cX_\dm)$ is an open substack of the Hom-stack 
    \[\operatorname{Hom}_{\mathfrak{M}_{g,n}^{\rm tw}}(\mathfrak{C}, \cX\times \mathfrak{M}_{g,n}^{\rm tw})
    \]
    constructed in \cite{hall2019coherent},
    where $\mathfrak{M}_{g,n}^{\rm tw}$ is the moduli stack of twisted curves of genus $g$ with $n$ marked point as constructed in \cite{olsson2007log}. In particular, $\mathfrak{Q}_{g,n}(\cX,\cX_\dm)$ is locally of finite type.
\end{Prop}
\begin{proof}

The proof is analogous to the one in \cite{orbifold_qmap_theory}, we sketch it here for the convenience of the reader. 
    
    Consider $\cH:=\operatorname{Hom}_{\mathfrak{M}_{g,n}^{\rm tw}}(\mathfrak{C}, \cX\times \mathfrak{M}_{g,n}^{\rm tw})$, the Hom-stack constructed in \cite{hall2019coherent},where $\mathfrak{C}\to \mathfrak{M}_{g,n}^{\rm tw}$ is the universal curve. First we show that the conditions for being a quasimap are representable by open embeddings. Consider then $\cU\to \cH$ an atlas which is a union of affine schemes.
    The morphism $\cU\to \cH$ gives a family of twisted curves $\pi:\sC\to \cU$, with a morphism $\phi:\sC\to \cX$, and it suffices to check that the points in $\cU$ where $\phi$ is a quasimap is open. 
    
    The condition that $\phi$ is representable is open from \cite{olshom}*{Corollary 1.6}.
Moreover, $\phi^{-1}(\cX\smallsetminus \cX_\dm)$ is a closed subset of $\cC$, so from the upper-semicontinuity of the dimension of the fibers, the locus where $\phi^{-1}(\cX\smallsetminus \cX_\dm)$ is either empty or has dimension 0 on $\cC$ is open. Similarly, the locus in $\cU$ where the closed $\Delta:=\phi^{-1}(\cX\smallsetminus \cX_\dm)\cap (\sC^{\operatorname{sing}}\cup \bigcup \Sigma_i)$ is empty is again open. Indeed, $\Delta$ is closed in $\sC$ as intersection of closed is closed, and since $\pi:\cC\to \cU$ is proper, also $\pi(\Delta)$, so $\cU\smallsetminus\pi(\Delta)$ is open as desired. 
\end{proof}
\begin{Prop}\label{prop_stable_qmap_is_an_open_cond}
    The inclusion $\cQ_{g,n}(\cX,\cX_\dm)\to \mathfrak{Q}_{g,n}(\cX,\cX_\dm)$ is an open embedding. In particular, $\cQ_{g,n}(\cX,\cX_\dm)$ is algebraic and locally of finite type over $\mathfrak{M}_{g,n}^{\rm tw}$.
\end{Prop}
\begin{proof}
    As $\mathfrak{Q}_{g,n}(\cX,\cX_\dm)\to \mathfrak{M}_{g,n}^{\rm tw}$ is locally of finite presentation over $\mathfrak{M}_{g,n}^{\rm tw}$ and $\mathfrak{M}_{g,n}^{\rm tw}$ is locally of finite type over $\spec(k)$, also $\mathfrak{Q}_{g,n}(\cX,\cX_\dm)$ is locally of finite type. In particular, there is a smooth cover $\cU\to \mathfrak{Q}_{g,n}(\cX,\cX_\dm)$ with $\cU$ a disjoint union of affine schemes of finite type over $\spec(k)$, with the two universal maps $\pi:\cC\to \cU$ and $\phi:\cC\to \cX$. Let $f:C\to X$ be the morphism induced by $\phi$ on good moduli spaces.
    
    From \cite{stacks-project}*{\href{https://stacks.math.columbia.edu/tag/0903}{Tag 0903}}, to check that being stable is an open condition, it suffices to check
     that the line bundle $\omega_\pi(p_1+\ldots+p_n)\otimes L^{\otimes 3}$ being nef and condition (2) in \Cref{lemma_equivalent_conditions_for_stability} satisfy these two conditions:
     \begin{enumerate}
         \item[(Const)] they hold for the generic point $\xi_W$ of a closed irreducible subscheme $W\subseteq \cU$ if and only if they hold for an open and dense subset of $W$, and
         \item[(Gen)] they are stable under generalization.
     \end{enumerate}
     It is now standard to show that being a quasimap satisfies (Const): the topological type of $\cC\to \cU$ is constructible, so it suffices to check that stability is constructible for families of curves with constant topological type, which is standard. So now we check (Gen): we can assume that $\cU=\spec(R)$ is the spectrum of a DVR with generic point $\eta$, and we need to prove that if the special fiber is stable, the generic one is stable. Observe that for doing so, we can replace $\spec(R)$ with a cover of it, possibly ramified.

     To show that $\omega_\pi(p_1+\ldots+p_n)\otimes L^{\otimes 3}$ is nef on the generic fiber, we can normalize $C$ to get $C^n\to C$ and work one connected component of $C^n$ at the time. So in particular, we can assume that $C_\eta$ is smooth, where $\eta$ is the generic point of $\cU$. Now from Riemann-Roch for nodal curves, and since the Euler characteristic of a line bundle is constant on flat and proper families, the degree of $\omega_\pi(p_1+\ldots+p_n)\otimes L^{\otimes 3}$ is constant. As on the special fiber is is non-negative, on the generic one will be non-negative as well.

     Assume then by contradiction that there is $\cD_\eta\subseteq \cC_\eta$ an irreducible component of the generic fiber, such that $\phi(\cD_\eta)\subseteq \cX_\dm$, but $f_{\dm,\eta}^*L_{X_\dm}$ is not ample, where $f_{\dm,\eta}:D_\eta\to X_\dm$ is the morphism induced by $\phi$ on coarse moduli spaces.
     Consider then $\cD$ the closure of $\cD_\eta$ in $\cC$, and $D$ its coarse moduli space. Then, up to possibly replacing $\spec(R)$ with a cover of it, the morphism $D\to \spec(R)$ is a family of rational curves, where the generic fiber is isomorphic to $\bP^1$, and the special fiber is a chain of rational curves.
     Moreover, $\cD_\eta$ has two closed substacks $n_{1,\eta}, n_{2,\eta}\to \cD_\eta$ which are gerbes over $\eta$, correspondings to the two nodes to which $\cD_\eta$ is attached.
     Similarly, the morphism $f:\cD\to \cX\to X$ maps each fiber of $\cD\to \spec(R)$ to a point. In other terms, the curves in $D$ are contracted via $f$.
     
     First observe that $\phi(\cD)$ cannot be contained in $\cX_\dm$, otherwise proceeding as before we would have that $f_{\dm,\eta}^*L_{X_\dm}$ has positive degree on the generic fiber, as it would have positive degree on the special one. Then, since $\phi$ is a quasimap, the locus $\phi^{-1}(\cX\smallsetminus \cX_\dm)$ is a disjoint union of closed smooth points on the special fiber. Pick then $\cD_p$ an
     irreducible component of the special fiber, with coarse moduli space $D_p$, and let $x\in \cD_p$ be a point on it which does not map to $\cX_\dm$.
     
     \textbf{Step 1.} We reduce to the case when the special fiber of $D\to \spec(R)$ is a single rational curve. 
     
     Consider $\widetilde{\Delta}$ the surface obtained by contracting all the irreducible components of the special fiber of $D$, with the exception of $D_p$. Then $\widetilde{\Delta}$ is a family of $\bP^1$s, so it is smooth as each fiber is smooth.
     Recall that we had two closed substacks $n_{i,\eta}\subseteq \cD_\eta$ which are gerbes, so $\cD_\eta\to D_\eta$ is the root-stack along the two nodal points of $C_\eta$ on $D_\eta$, with indices $d_1$ and $d_2$ respectively. Let $\widetilde{n}_{i,\eta}$ be the coarse moduli space of of ${n}_{i,\eta}$,
     and let $\widetilde{n}_{i}$ be its closure in $\Delta$. From how $\widetilde{\Delta}$ is constructed, one can check that $\widetilde{n}_{1}$ and $\widetilde{n}_{2}$ are smooth and don't intersect. Let finally $\Delta$ be the root-stack of $\widetilde{\Delta}$ along $\widetilde{n}_{1}$ and $\widetilde{n}_{2}$ with indices $d_1$ and $d_2$ respectively, and let $n_i$ be the closure of $n_{i,\eta}$. We summarize the notations:
     \begin{enumerate}
         \item $\widetilde{\Delta}$ is a family of $\bP^1$s, and is the coarse moduli space of $\Delta$.
         \item $\Delta$ is obtained from $\widetilde{\Delta}$ by taking root stacks along $\widetilde{n_i}$, in a way such that $\Delta_\eta$ agrees with the original $\cD_\eta$,
         \item $n_i=\widetilde{n_i}\times_{\widetilde{\Delta}}\Delta$ is the closure in $\Delta$ of the nodes on of $\cC_\eta$ that lie on $\cD_\eta$.
     \end{enumerate}In particular:
     \begin{enumerate}
         \item $\Delta\to\spec(R)$ is a family of smooth twisted curves, with the generic fiber that has a map to $\alpha: \Delta_\eta\to \cX_\dm$ that factors via $\Delta_\eta\to \cB\bmu_d\to \cX_\dm$,
         \item as $\Delta$ and $\widetilde{\Delta}$ agree away from $n_{i}$, there is a map $\alpha:\Delta \smallsetminus (\widetilde{n}_{1,p}\cup \widetilde{n}_{2,p})\to \cX$, where $\widetilde{n}_{i,p}$ is the intersection of $\widetilde{n}_{i}$ with the special fiber, and which extends $\alpha_\eta$,
         \item the codimension one locus of $\Delta$ maps to $\cX_\dm$ via $\alpha$, but there is a closed point $x$ of $\Delta$ that does not map to $\cX_\dm$, and
         \item the composition $\Delta \smallsetminus (\widetilde{n}_{1,p}\cup \widetilde{n}_{2,p})\to \cX\to X$ extends to $\Delta\to X$, and contracts each fiber of $\Delta\to\spec(R)$.
     \end{enumerate}
     Then from \Cref{lemma_generalization_purity_tsm1}, the map $\alpha$ extends to a morphism $\Delta\to \cX$. We still denote such an extension by $\alpha.$ 

     \textbf{Step 2.} The morphism $\alpha|_{\Delta^\circ}:\Delta^\circ\to \cX_\dm$ defined on the codimension one locus $\Delta^\circ\subseteq \Delta$ that maps to $\cX_\dm$ extends to $\alpha':\Delta\to \cX_\dm$.

     Indeed, the morphism on the generic fiber $\Delta_\eta\to \cX_\dm\to X_\dm$, by assumption, factors via $\eta\to X_\dm$ as follows, up to potentially taking an extension of $\spec(R)$:
     \[
     \xymatrix{\Delta_\eta\ar[r]\ar[d] & \cX_\dm\ar[r] & X_\dm\\\eta\ar[rru] & & }
     \]
     As by assumption $X_\dm$ is proper, there is an extension $\eta\to X_\dm$ to $\spec(R)\to X_\dm$ giving a map \[\widetilde{\alpha}':\Delta\to \spec(R)\to X_\dm.\] Since $X_\dm$ is separated and $\Delta$ is smooth, the two maps $\widetilde{\alpha}|_{\Delta^\circ}$ and $\widetilde{\alpha}'$ agree on codimension one, so up to shrinking $\Delta^\circ$ we can assume that they agree. In other terms, we can extend $\Delta^\circ\to \cX_\dm\to X_\dm$ to a morphism $\widetilde{\alpha}':\Delta\to X_\dm$. Then from \Cref{lemma_generalization_purity_tsm1}, the map $\alpha'$ extends to a morphism $\Delta\to \cX_\dm$. Moreover $\alpha$ and $\alpha'$ agree in codimension one. 

     \textbf{Contradiction.} From \Cref{lemma_two_maps_that_agree_in_codim_2_are_the_same}, the two maps $\alpha$ and $\alpha'$ agree. However there was a point $x\in \Delta$ such that $\alpha(x)\not \in \cX_\dm$, which is the desired contradiction. \end{proof}

\subsection{Valuative criterion for properness}\label{section_qmaps_satisfy_the_val_crit_for_prop}
In this subsection we will prove that $\cQ_{g,n}(\cX,\cX_\dm)$ satisfies the valuative criterion for properness.
\begin{Teo}
    We will adopt \Cref{notation_cX_and_cX_dm}, moreover let $R$ be a DVR, let $\eta$ be the generic point of $\spec(R)$ and $p$ the closed one. Let $(\phi_\eta:\sC_\eta\to \cX; \Sigma_{1},\ldots,\Sigma_{n})$ be an $n$-marked stable quasimap to $(\cX_{\dm},\cX)$, over $\eta$.
    Then, up to replacing $\spec(R)$ with a possibly ramified cover of it, there is a unique $n$-marked stable quasimap to $(\cX_{\dm},\cX)$ over $\spec(R)$ extending $(\phi_\eta:\sC_\eta\to \cX; p_{1},\ldots,p_{n})$. 
\end{Teo}
\begin{proof}Our argument adapts \cite{orbifold_qmap_theory}*{2.4.4} using some results from \Cref{section_background}.    

    \textbf{Step 1.} Let $q_1,\ldots ,q_r$ be the points of $\sC_\eta$ which do not map to $\cX_{\dm}$. Then there is a root-stack $\cC_\eta\to \sC_\eta$ supported at the points $q_1,...,q_r$ such that the rational map $\phi_\eta|_{\sC_\eta\smallsetminus \{q_1,\ldots q_r\}}$ extends to $\psi_\eta:\cC_\eta\to \cX_{\dm}$.
    This follows for example from \cite{bresciani2022arithmetic} or \cite{orbifold_qmap_theory}*{Lemma 2.5}. We denote by $\Xi_i$ the gerbe over the point $q_i$. Then $(\psi_\eta:\cC_\eta\to \cX_{\dm}, \Sigma_1,\ldots,\Sigma_n,\Xi_1,\ldots,\Xi_r)$ is a pointed twisted stable map as in \cite{AV_compactifying}.

    \textbf{Step 2.} Since the moduli space of twisted stable maps constructed in \cite{AV_compactifying} is proper, up to replacing $R$ with a possibly ramified cover of it, there is a twisted stable map $(\psi_R:\cC_R\to \cX_{\dm}, \overline{\Sigma}_1,\ldots,\overline{\Sigma}_n,\overline{\Xi}_1,\ldots,\overline{\Xi}_r)$ extending the twisted stable map of the previous point.

    \textbf{Step 3.} Let $\cC_R\to C_R'$ be the coarse moduli space. Since the coarse moduli space of $\cX_{\dm}$ maps to the good moduli space $X$ of $\cX$, there is a morphism $C_R'\to X$. Let then $\Delta_1,\ldots,\Delta_m$ be the rational tails on the central fiber of $C_R'\to \spec(R)$ which are contracted by $C_R'\to X$ and which have no markings among $\overline{p}_1,\ldots, \overline{p}_n$, where $\overline{p}_i$ is the coarse moduli space of $\overline{\Sigma}_i$. 
    Observe that $\Delta_1\cup\ldots\cup \Delta_m$ do not contain the whole central fiber. Indeed, if that was the case, the central fiber of $C'_R$ would map to a point on $X$, it would intersect no marking $\overline{p}_i$, and would be a chain of $\bP^1$s. But then also the generic fiber would map to a point, have no markings and be a chain of $\bP^1$s, which is impossible since $\sC_\eta\to \cX$ is stable. Therefore let $C_R$ be the curve obtained from $C'_R$ by contracting $\Delta_1,\ldots,\Delta_m$. So to summarize:
    \begin{enumerate}
        \item[(i)] $\cC_R\to C'_R$ is the coarse moduli space,
        \item[(ii)] $C'_R\to C_R$ is the contraction of the rational tails on the central fiber, which do not intersect the sections $\overline{p}_i$ and which map to a point on $X$, and
        \item[(iii)] there is a morphism $C_R\to X$.
    \end{enumerate}Consider then the twisted curve $\sC_R\to \spec(R)$ constructed as follows:
    \begin{enumerate}
        \item it has $C_R$ as coarse moduli space,
        \item is isomorphic to $\cC_R$ away from $\overline{\xi}_1,\ldots \overline{\xi}_r$, where $\overline{\xi}_i$ is the coarse moduli space of $\overline{\Xi}_i$, and from the locus where $\cC_R\to C_R$ is not quasi-finite, and 
        \item it is isomorphic to $C_R$ elsewhere.
    \end{enumerate}
    In particular, the generic fiber of $\sC_R\to \spec(R)$ is isomorphic to $\sC_\eta$, it has the same twisted nodes as $\cC_R$, there is a morphism $\cC_R\to \sC_R$ which is an isomorphism away from both $\{\overline{q}_i\}_i$ and the locus where $\cC_R\to C_R$ is not quasi-finite. We summarize the situation in the following diagram, where the vertical maps are coarse moduli spaces, both $\cC_R$ and $\sC_R$ are twisted curves, and the map $\alpha$ contracts some rational tails:
    \[
    \xymatrix{\cC_R\ar[d] \ar[r] & \sC_R\ar[d] \\ C'_R\ar[r]^{\alpha} & C_R}
    \]

    \textbf{Step 4.} We have a rational morphism $\sC_R\dashrightarrow \cX$ which on the generic fiber is given by $f_\eta$, and on the locus where $\cC_R\to \sC_R$ is an isomorphism it is given by $\psi_R$. In particular, the indeterminacy locus of $\sC_R\to \cX$ consists of finitely many smooth points on the central fiber of $\sC_R$. Then we can extend it using \Cref{lemma_generalization_purity_tsm1}.

    \textbf{Step 5.} We prove that the resulting map $(\sC_R\to \cX, \overline{\Sigma}_1,\ldots ,\overline{\Sigma}_n)$ is stable. It is stable along the generic fiber, so it suffices to prove that it is stable along the central fiber. We denote by $C_0$ be the central fiber of $C_R\to \spec(R)$, with the two maps $f_0:C_0\to X$ and $\psi_0:C_0\to X_\dm$, and by $\overline{p}_{0,i}$ the one of $\overline{p}_i$ for every $i$. 

    Since to construct $C_R$ we contracted all the rational tails with no markings and which map to a point in $X$, we have that $\omega_{C_0}( \overline{p}_{0,1},\ldots,\overline{p}_{0,n})\otimes f_0^*L^{\otimes 3}$ is nef. So let's assume that $D\subseteq C_0$ is an irreducible component such that $\omega_{C_0}( \overline{p}_{0,1},\ldots,\overline{p}_{0,n})\otimes f_0^*L^{\otimes 3}$ has degree 0 over it. Then it is either an elliptic curve mapping to a point and with no markings, or a curve isomorphic to $\bP^1$ with two markings. If it does not map to $\cX_\dm$, then it is stable by definition of stability.
    If it does map to $\cX_\dm$, then by construction of $\cC$ it is isomorphic to an irreducible component of the special fiber of $\cC$, which does not intersect $\overline{\Xi}_i$. Since $(\cC,\overline{\Sigma}_1,\ldots,\overline{\Sigma}_n, \overline{\Xi}_1,\ldots,\overline{\Xi}_r)\to \cX_\dm$ is stable, $L_{X_\dm}$ is ample when pulled back to $D$. 

    \textbf{Step 6.} We now focus on the uniqueness part of the valuative criterion for properness. First we prove that if there are two extensions $\phi:\sC\to \cX$ and $\phi':\sC'\to \cX$ such that $\phi$ and $\phi'$ are quasimaps, a priori not stable, but with $\sC$ and $ \sC'$ having isomorphic coarse moduli spaces, then $\sC\cong \sC'$ and $\phi\cong \phi'$. 

    Let $U$ be the locus on $\sC_R$ where $\sC_R\cong \sC'_R$. The $U$ contains the generic fiber and the codimension one locus of the special fiber. The maps $\phi|_U, \phi|_U':U\to \cX$ agree in codimension one, as their codimension one points either in $\sC_\eta$, or on the locus where $\phi|_U$ and $\phi'|_U$ map to $\cX_\dm$. For points in $\sC_\eta$ we have that $\phi|_{\sC_\eta} = \phi'|_{\sC'_\eta}$ by assumption. Instead, each point $\xi$ of codimension one not in $\sC_\eta$ maps to $\cX_\dm$, and $\cO_{\sC_R,\xi}$ is a DVR, with generic point contained in $\sC_\eta$. So two extensions of $\phi_\eta$ to $\xi$ must agree from the valuative criterion for separatedness applied to $\cX_\dm$.
    
    We can now use \Cref{lemma_two_maps_that_agree_in_codim_2_are_the_same} to argue that $\phi|_U\cong \phi'|_U$. Moreover, from the uniqueness in \Cref{lemma_generalization_purity_tsm1}, the twisted curves $\sC_R$ and $\sC'_R$ agree on the locus where $C_R$ and $C'_R$ have local quotient singularities on the special fiber and $\phi\cong \phi'$ on this locus. Observe that this locus consists of the complement of the nodes of the special fiber which do not smooth. But then $\sC_R\cong \sC'_R$ as they are twisted curves, and the two maps $\phi\cong \phi'$ from \cite{AV_compactifying}*{Lemma 2.4.1}, as by assumption $\phi$ and $\phi'$ are quasimaps, so map the nodal locus in $\cX_\dm$.

     \textbf{Step 7.} Let $\phi:\sC_R\to \cX$ and $\phi':\sC_R'\to \cX$ two extensions, with coarse moduli spaces $C_R$ and $C'_R$ respectively. We show that, if $\widetilde{C}_R$ is a minimal resolution of the rational map $C_R \dashrightarrow C'_R$ with two maps $a:\widetilde{C}_R\to C_R$ and $a':\widetilde{C}_R\to C'_R$, then up to replacing $\spec(R)$ with a possibly ramified cover of it, there is a twisted curve $\widetilde{\sC}_R$ with coarse moduli space $\widetilde{C}_R$, and two maps $\alpha:\widetilde{\sC}_R\to \sC_R$ and $\alpha':\widetilde{\sC}_R\to \sC_R'$ restricting to $a$ and $a'$.

     It suffices to lift $\widetilde{C}_R\to C_R\times C'_R$ to $\widetilde{C}_R\to \sC_R\times \sC'_R$. Observing that $\sC_R\times \sC'_R\to C_R\times C'_R$ is the coarse moduli space map, we proceed as in \cite{AV_compactifying}. We can lift $a\times a'$ to $\alpha\times \alpha'$ in codimension one, up to performing a ramified cover of $\spec(R)$, using Abihankar's lemma and descent as in \cite{AV_compactifying}*{Proposition 6.0.4, Step 2}. We can also lift $a$ to get $\alpha$ away from the nodes of the special fiber of $\widetilde{C}_R$ which smooth. As those are locally finite quotient singularities, can finally use \Cref{lemma_generalization_purity_tsm1} to obtain $\alpha\times \alpha'$. One can check that the resulting $\widetilde{\sC}_R\to \spec(R)$ is a family of twisted curves.

     \textbf{Step 8.} End of the argument.

     We now have two maps $\widetilde{\sC}_R\xrightarrow{\alpha} \sC_R\xrightarrow{\phi} \cX$ and $\widetilde{\sC}_R\xrightarrow{\alpha'} \sC_R'\xrightarrow{\phi'} \cX$. As each component of $\widetilde{C}_R$ which is contracted with either $a$ or $a'$ maps to a node, and since $\phi$ and $\phi'$ are quasimaps, both $\phi\circ\alpha$ and $\phi'\circ\alpha'$ are quasimaps. So from Step 6 they agree. Assume by contradiction that there is $\sE\subseteq \widetilde{\sC}_R$ which is contracted via $\alpha'$ but not via $\alpha$. Then by the stability condition (2) in \Cref{lemma_equivalent_conditions_for_stability}, there are two points $p,q\in \sE$ such that $\phi(\alpha(p))$ and $\phi(\alpha(q))$ are not isomorphic. But $\alpha'(p)$ and $\alpha'(q)$ are isomorphic instead, as they map to a node in $\sC'_R$. So $\phi'(\alpha'(p))\cong \phi'(\alpha'(q))$, which is the desired contradiction.

     It is straightforward to check that $\Sigma_i$ extend uniquely once $\sC_R$ is fixed.
\end{proof}
\subsection{Boundedness}\label{section_boundedness_qmaps}
We now add some numeric invariants to stable quasimaps, to prove that if one fixes those the resulting moduli space is bounded.
We begin by generalizing the definition of \textit{class} of a quasimap to our setting:
\begin{Def}\label{def_beta}
    Given a quasimap $(\phi\colon \sC \to \cX,\Sigma_1,\ldots,\Sigma_n)$ over a geometric point, the \emph{class} of the quasimap is the homomorphism
    \[\phi_*[\sC] \colon \Pic(\cX) \longrightarrow \bQ,\quad L \longmapsto \deg(\phi^*L).\]
\end{Def}
In the definition above, the degree of a line bundle $M$ on a twisted curve $\sC$ is computed as follows: pick $e$ such that $M^{\otimes e}=\pi^*N$, where $\pi:\sC\to C$ is the coarse moduli space. Then $\deg(M)=\frac{\deg(N)}{e}$.
\begin{Remark}
    Given a morphism $g:\cX \to \cY$, we can define $g_*(\phi_*[\sC])$ as the homomorphism $L\longmapsto \phi_*[\sC](g^*L)$. This pushforward is functorial with respect to composition of morphisms, and if $\rho:\sC\to C$ is the coarse moduli space, we have $\rho_*({\rm id}_*[\sC])={\rm id}_*[C]$.
\end{Remark}

Throughout the remaining part of this section, we will make the following additional
\begin{Assumptions}\label{assumptions:extension of line bundle}
We assume that:
    \begin{enumerate}
        \item there exists a line bundle $\cL$ on $\cX$ such that $\cX_\dm = \cX(\cL)^{ss}_X$, where the latter denotes the relative $\cL$-semistable locus of $\cX$ over the good moduli space $X$. In particular, there exists $m>0$ such that
    \[ \cL|_{\cX_\dm}^{\otimes m} \simeq \pi^*L_{X_\dm}, \]
    where $\pi\colon \cX_{\dm} \to X_\dm $ is the coarse moduli space morphism,
    \item the stack $\cX$ is a global quotient stack, i.e. $\cX=[W/G]$, with $G$ reductive.
    \end{enumerate}
\end{Assumptions}
\begin{Remark}\label{rmk_if_W_is_Qfactorial_then_any_open_is_the_ss_locus_of_a_linebundle}
    Condition (1) is automatically satisfied if $\cX$ is smooth by \cite{Alp}*{Theorem 11.14.}, and the same proof goes through if $\cX=[W/G]$ and $W$ is $\bQ$-factorial. See also \Cref{cor_when_X_dm_is_ss_locus_for_a_lb} for other cases when Condition (1) holds.
\end{Remark}


    \begin{Def}
        We define  $\cQ_{g,n}(\cX,\cX_\dm,\beta)$ as the locus in  $\cQ_{g,n}(\cX,\cX_\dm)$ parametrizing stabe quasimaps with class $\beta$. 
    \end{Def}
    \begin{Remark}
        Observe that $\cQ_{g,n}(\cX,\cX_\dm,\beta)$ are open and close in $\cQ_{g,n}(\cX,\cX_\dm)$.
    \end{Remark}
    The main goal of this subsection is to prove the following.
    \begin{Teo}\label{thm_boundedness}
    Given $\cX$ an algebraic stack satisfying \Cref{assumptions:extension of line bundle}, the algebraic stack
        $\cQ_{g,n}(\cX,\cX_\dm,\beta)$ is of finite type.
    \end{Teo}
    We begin with the following

\begin{Lemma}\label{lemma:type is bounded}
        Let $(\phi:\sC\to \cX,\Sigma_1,\ldots,\Sigma_n)$ be a stable quasimap of class $\beta$. Then the topological type of $C$, the number of stacky points of $\sC$ and the automorphism groups of the stacky points are bounded.
\end{Lemma}
In this argument we will use (1) of \Cref{assumptions:extension of line bundle}. The proof is a mild adaptation of \cite{ciocan2010moduli}*{Corollary 3.1.5}.
\begin{proof}
    Assuming that the topological type of $C$ is bounded, the number of stacky points is then also bounded, because non-trivial stacky structures are only allowed at the nodes and at the markings. The automorphism groups must be cyclic groups, so we only need to bound the orders. As $\phi$ is representable and the stacky points are mapped to $\cX_\dm$ (which is quasicompact), we deduce that the orders are bounded.

    To bound the topological type, let $C\to C^{st}$ be the contraction induced by the nef line bundle $\omega_C(\sum p_i)\otimes f^*L_X^{\otimes 3}$. The topological type of $C^{st}$ is bounded, because if it is not a point then $f^{st}\colon C^{st}\to X$ is Kontsevich-stable of class $\pi_*(\beta)$, where $\pi\colon\cX\to X$ is the good moduli space homomorphism. 
    
    We are left with bounding the topological type of the exceptional locus of the contraction, which includes the case in which $C^{st}$ is a point. As genus one smooth curves are bounded, we only need to show that the number of rational curves which are contracted by $C\to C^{st}$ is also bounded.

    For this, it is enough to observe that, from the moreover part of \Cref{lemma_equivalent_conditions_for_stability}, for every such component $\sD$ we have $\deg(\phi^*\cL|_\sD)>0$. The desired conclusion follows as $\sum_{\sD\subset \sC} \deg(\phi^*\cL|_\sD) \leq \beta(\cL)$. 
\end{proof}

\subsubsection{Boundedness of quasimaps $C\to \cX$ from a fixed curve}
To bound the moduli space of quasimaps with a given class $\beta$, we need to bound the stack of quasimaps from a fixed curve $C$ to $\cX$. For doing so, we plan to use the following.
\begin{Teo}[\cite{ciocan2014stable}*{Theorem 3.2.5}]
    Let $\beta \in \Hom(\chi(G),\bZ)$  and a smooth projective
curve $C$ be fixed. Let $V$ be a vector space with an action of $G$ via a
representation $G\to \GL(V)$ with finite kernel and let $\chi\in \mathbf{X}(G)$ be
a character such that $V^s(G,\chi)\neq\emptyset$. Then the family of principal $G$-bundles $P$ on $C$ of degree $\beta$ such that the vector bundle $P \times_G V$
admits a section $u$ which sends the generic point of $C$ to $V^s(G,\chi)$ is
bounded.
\end{Teo}
For this, we will prove the following.
\begin{Teo}\label{thm_boundedness_embedding_in_vector_space}
    Assuming \Cref{assumptions:extension of line bundle}, there is a group $G'$, a character $\chi:G'\to \Gm$ and an action on $\bA^n$, such that there is a locally closed embedding $\iota:\cX\to [\bA^n/G']$ satisfying $\iota^{-1}([\bA^n(\Bbbk_\chi)^{ss}/G'])=\cX_\dm$.
\end{Teo}

We begin with the following lemma, which will allow to pass from relative GIT for a line bundle, to relative GIT for a character, which is better behaved, as we can use the affine Hilbert-Mumford criterion, see \Cref{lemma_relative_GIT_for_gms_can_be_checked_using_HM}.

\begin{Lemma}\label{lemma_cone_construction_over_X}
    Let $\cX=[W/G]$ be a quotient stack by a reductive group $G$, admitting a good moduli space $\cX\to X$ and let $\cL$ be a line bundle on $\cX$. Let $\chi:G\times \Gm\to \Gm$ the character $(g,t)\mapsto t^{-1}$.
    Then there is a scheme $C$ with an action of $G\times \Gm$ such that
\begin{itemize}
    \item there is an isomorphism $[C(\Bbbk_\chi)^{ss}_X/G\times \Gm]\cong \cX(\cL)^{ss}_X$, 
    \item  there is an isomorphism $[C/G\times \Gm]\cong \cX$, and 
    \item the inclusion $C(\Bbbk_\chi)^{ss}_X\subseteq C$ induces $\cX(\cL)^{ss}_X\to \cX$.
\end{itemize}
    In particular, for every Zariski cover $\spec(A)\to X$, 
        \[C(\Bbbk_\chi)^{ss}_X\times_X\spec(A)=(C\times_X\spec(A))(\Bbbk_\chi)^{ss}.\]
\end{Lemma}
\begin{proof}

    With the same notation of the proof of \Cref{lemma_git_for_ample_linearization_is_contained_in_the_affine_git_for_the_cone}, we can pick $C=P_\cL$ the $\Gm$-torsor associated to the line bundle $\cL$, so that $[C/G\times\Gm]\simeq\cX$. 

    For every Zariski open cover $\spec(A) \to X$, we have that $W\times_X \spec(A)$ is affine, because $\pi:W\to X$ is affine as $W\to [W/G]$ is a $G$-torsor and $[W/G]\to X$ is cohomologically affine. In particular, this implies that the pullback of $\cL$ to $W\times_X\spec(A)$ is ample from \cite{stacks-project}*{\href{https://stacks.math.columbia.edu/tag/0890}{Tag 0890}} and the fact that, over an affine scheme, each coherent module is globally generated. 

    Furthermore, we can assume that $\cX(\cL)^{ss}_{X}\times_X \spec(A)$ is isomorphic to $[(W\times_X \spec(A))(L)^{ss} /G]$, where $L$ denotes the $G$-equivariant line bundle on $W\times_X \spec(A)$ which descends to the pullback of $\cL$: indeed, by definition of relative semistable locus, for every point $p\in \cX(\cL)^{ss}_{X}\times_X \spec(A)$, up to shrinking the base we can always assume that there exists a section of the pullback of $\cL$ that does not vanish on $p$.
    
    By \Cref{lemma_git_for_ample_linearization_is_contained_in_the_affine_git_for_the_cone} we get that $[(W\times_X \spec(A))(L)^{ss} /G]$ is isomorphic to $[(C\times_X \spec(A))(\Bbbk_\chi)^{ss}/G\times\Gm]$.
    On the other hand, by definition, the latter is equal to $[C(\Bbbk_\chi)^{ss}_X \times_X \spec(A) / G\times\Gm]$ as desired.
\end{proof}
\begin{Lemma}\label{lemma_L_otimes_epsilon_k_chi}
   Let $\cX=[W/G]$ be a quotient stack by a reductive group $G$, admitting a projective good moduli space $\cX\to X$ and let $\chi:G\to \Gm$ be a character.
    Then there are $G\times\Gm$-equivariant schemes $\spec(A)$ and $V$ together with a $G\times\Gm$-equivariant open embedding $V\subseteq \spec(A)$ and a character $\rho:G\times \Gm\to \Gm$ such that:
    \begin{itemize}
    \item there is an isomorphism $[\spec(A)(\Bbbk_\rho)^{ss}/G\times \Gm]\cong \cX(\Bbbk_\chi)^{ss}_X$, 
    \item  there is an isomorphism $[V/G\times \Gm]\cong \cX$, and 
    \item there is an inclusion $\spec(A)(\Bbbk_\rho)^{ss}\subseteq V$ inducing $\cX(\Bbbk_\chi)^{ss}_X\to \cX$.
\end{itemize}
\end{Lemma}
%\begin{Def} Let $\cX\to \cY$ be a morphism of algebraic stacks with a line bundle $\cL$ on $\cX$.We call the Hilbert-Mumford stable locus of $\cX$ with respect to $\cY$, the set of points $u\in \cX$ such that for every morphism $\phi:\Theta\to \cX$ sending $1\mapsto u$ and $\cB \Gm\mapsto \overline{u}$, and such that $\Theta\to \cX\to \cY$ factors via $\Theta\to \spec(k)\to \cY$, the weight of $\Gm$ on $\phi^*\cL_2|_{\cB\Gm}$ is non-positive. We denote this locus by $\cX(\cL)^{\operatorname{HM}}_\cY$.\end{Def}

%\begin{Prop}
 %Let $\cX_1\to \cX_2\to \cS$ be a morphism of algebraic stacks, with line bundles $\cL_i$ on $\cX_i$. Assume that $\cX_2(\cL_2)^{\operatorname{HM}}_\cY$ is open non-empty, ad assume that also  $\cU:=(\cX_2(\cL_1)^{\operatorname{HM}}_\cY\times_{\cX_2}\cX_1)(\cL_2)^{\operatorname{HM}}_{\cX_2(\cL_1)^{\operatorname{HM}}_\cY}$ is open non-empty. Then $\cU=\cX_1(\pi^*\cL_2^{\otimes m}\otimes \cL_1)^{\operatorname{HM}}_\cY$ for $m\gg1$.\end{Prop}

\begin{proof}
    From \cite{teleman2000quantization}*{Lemma 6.1} and its proof, there is a $G$-equivariant compactification $W\subseteq \overline{W}$, with $\overline{W}$ a projective variety with ample $G$-linearization $L$, and such that $W=\overline{W}(L)^{ss}$. We can now apply \Cref{prop_stable_locus_of_pi*L^n_otimes_G}
    with $\cX_1=\cX_2=[\spec(\bigoplus_{n\ge 0}\oH^0(\overline{W}, L^{\otimes n}))/G\times \Gm]$, the linearization of \Cref{prop_stable_locus_of_pi*L^n_otimes_G} denoted by $\cL$ is given by $\theta:G\times \Gm\to \Gm$, $(g,t)\mapsto t^{-1}$ and the other linearlization in \Cref{prop_stable_locus_of_pi*L^n_otimes_G} is $\chi$.
    Up to replacing $L$ with a higher power of it, there is an inclusion $\overline{W}\subseteq \cX_1$ and from \Cref{lemma_git_for_ample_linearization_is_contained_in_the_affine_git_for_the_cone}, up to replacing $L$ with a higher power, the linearization $\cL$ has $W$ as semistable locus over $\spec(k)$. Observe that, if we denote by $X_1$ the good moduli space of $\cX_1$, then $X_1\to \spec(k)$ is finite: indeed, we have $X_1=\spec(H^0(\overline{W},\cO_{\overline{W}})^G)$,  and the latter is a finite dimensional vector space as $\overline{W}$ is projective. Then by definition of relative semistable locus, $\cX_i(\cG)^{ss}_X = \cX_i(\cG)^{ss}_{\spec(k)}$ for every $i$ and for every line bundle $\cG$ on $\cX_i$.
    
    So from \Cref{prop_stable_locus_of_pi*L^n_otimes_G} the semistable locus of $\cX_2$ over $X_1$ (hence also over $\spec(k)$) for the line bundle $\cL^{\otimes m}\otimes\Bbbk_\chi$ is $\cX(\Bbbk_\chi)^{ss}_X$. In other terms,
    \[
    \cX(\Bbbk_\chi)^{ss}_X = [\spec(\bigoplus_{n\ge 0}\oH^0(\overline{W}, L^{\otimes n}))/G\times \Gm](\Bbbk_{\theta}^m\otimes \Bbbk_\chi)^{ss}.
    \]
We can take $[V/G\times \Gm]=[\spec(\bigoplus_{n\ge 0}\oH^0(\overline{W}, L^{\otimes n}))/G\times \Gm](\Bbbk_{\theta})^{ss}=[W/G]$ and $\rho=\Bbbk_{\theta}^m\otimes \Bbbk_\chi$, and so \[[\spec(\bigoplus_{n\ge 0}\oH^0(\overline{W}, L^{\otimes n}))/G\times \Gm](\Bbbk_{\theta}^m\otimes \Bbbk_\chi)^{ss} = \cX(\Bbbk_\chi)^{ss}_X\subseteq \cX =[\spec(\bigoplus_{n\ge 0}\oH^0(\overline{W}, L^{\otimes n}))/G\times \Gm](\Bbbk_{\theta})^{ss}\] as desired.
\end{proof}
\begin{Lemma}\label{cor_embed_into_spec(A)/G}
    Let $\cX=[W/G]$ be an algebraic stack with a projective good moduli space $\cX\to X$, and let $\cL$ be a line bundle on $\cX$. Then there are schemes $\spec(A)$ and $V$, both endowed with an action of $G\times \Gm^2$, an equivariant embedding $V\subset\spec(A)$ and a character $\rho:G\times \Gm^2\to \Gm$ such that:
    \begin{itemize}
    \item there is an isomorphism $[\spec(A)(\Bbbk_\rho)^{ss}/G\times \Gm]\cong \cX(\cL)^{ss}_X$, 
    \item  there is an isomorphism $[V/G\times \Gm]\cong \cX$, and 
    \item there is an inclusion $\spec(A)(\Bbbk_\rho)^{ss}\subseteq V$ inducing $\cX(\cL)^{ss}_X\to \cX$.
\end{itemize}
\end{Lemma}
\begin{proof}
From \Cref{lemma_cone_construction_over_X}, up to replacing $G$ with $G\times \Gm$, we can assume that $\cX(\cL)^{ss}_X$ is given by a character, and the case in which $\cL$ is a character is treated in \Cref{lemma_L_otimes_epsilon_k_chi}.
\end{proof}
\begin{proof}[Proof of \Cref{thm_boundedness_embedding_in_vector_space}] This now follows from \Cref{cor_embed_into_spec(A)/G} since we can find an equivariant closed embedding $\spec(A)\to \bA^n$ which maps the semistable locus in the semistable locus.
\end{proof}


    \subsubsection{Boundedness of $\cQ_{g,n}(\cX,\cX_\dm,\beta)$}
    
    We are finally ready to prove \Cref{thm_boundedness}. Our strategy follows \cite{orbifold_qmap_theory}*{\S 2.4.3}, and can be understood as follows.
There is a morphism $\pi:\cQ_{g,n}(\cX,\cX_\dm,\beta)\to \mathfrak{M}_{g,n}^{\operatorname{tw}}$, its image is a substack of finite type, so it suffices to show that $\pi$ is of finite type. 

    In \cite{orbifold_qmap_theory}*{\S 2.4.3} it is proved that the fibers of $\pi$ are bounded. The proof in \emph{loc. cit.} actually goes through in families, up to stratifying the base of the family.
        We now recall the steps needed to bound quasimaps from a fixed twisted curve $\cC$, as in \cite{orbifold_qmap_theory}*{\S 2.4.3}. In \cite{orbifold_qmap_theory}*{\S 2.4.3} the authors first observe that $\cC$ fits in a push-out diagram wit $\cN\sqcup \cN\to \cC^n$ and $\cN\sqcup \cN\to \cN$, where $\cN\to \cC$ is the nodal locus and $\cC^n\to \cC$ the normalization.
        So it suffices to bound the quasimaps $\cC^n\to \cX$ and the gluing data on $\cN$. This the content of Step 2 of the proof of \Cref{thm_boundedness}. Bounding the gluing data on $\cN$ is instead achieved in Step 1 instead . To bound the maps $\cC^n\to \cX$, in \cite{orbifold_qmap_theory}*{\S 2.4.3} the authors take a finite flat cover $C\to \cC^n$, and show that it suffices to bound quasimaps $C\to \cX$ together with an isomorphism $C\times_\cC C\to \cX$ which satisfies the cocycle condition. In other terms, there is a closed embedding \[\Hom(\cC^n,\cX)\to \Hom(C,\cX)\times_{\Hom(C\times_{\cC^n}C,\cX)}\Hom(C,\cX)\]
    given by the morphisms which satisfy the cocycle condition. This is the content of Step 3. Finally one can use \Cref{thm_boundedness_embedding_in_vector_space} to reduce to \cite{ciocan2014stable}*{Theorem 3.2.5}.

    Observe also that, while \cite{ciocan2014stable}*{Theorem 3.2.5} is stated for a fixed curve, their proof goes through in families of smooth curves, up to possibly stratifying the base. Indeed, their argument is as follows, at least for the case $G=\GL_n$. First they recall that each vector bundle on a curve can be obtained as a consecutive extensions of line bundles.
    This is proved in \cite{ciocan2014stable}*{Lemma 3.2.6}, \cite{seshadri1972quotient}*{Theorem 1.9}, but in the case of $\GL_n$ the proof in \cite{beauville1996complex}*{Lemma III.11.1} generalizes for $n\ge 2$. They then show that all the vector bundles on $C$ induced by $C\xrightarrow{\phi} [\spec(A)/\GL_n]\to \cB \GL_n$ where $\phi$ has degree $\beta$, are consecutive extensions by line bundles \textit{with bounded degree} (this is a key step, achieved in \cite{ciocan2014stable}*{Lemma 3.2.8}). As line bundles with bounded degree are bounded (as they are parametrized by a relative Jacobian), and extensions of a bounded family of
    pairs consisting of a smooth curve $C$ and a vector bundle on $C$, by a fixed degree line bundle, are bounded (as parametrized by the push-forward, via a relative Jacobian, of a certain $\cE xt^1$), the possible vector bundles which appear are bounded.
    This is the argument of \cite{ciocan2014stable}*{Lemma 3.2.7},  using \cite{holla2001generalisation}*{Proposition 3.1 and Lemma 3.3} and \cite{grothendieck1957fondements}*{Proposition 1.2, page 221}. \footnote{Technically, the version of boundedness proved in \cite{grothendieck1957fondements}*{Proposition 1.2, page 221} and \cite{holla2001generalisation}*{Proposition 3.1 and Lemma 3.3} is weaker than what we need. Indeed, in \textit{loc. cit.} by saying that a class of objects $\mathfrak{C}$ is bounded the author means that there is a morphism of finite type $f:X\to B$ where each member of $\mathfrak{C}$ is a fiber of $f$. However, in our setting, the same proof gives a stronger boundedness: there is a morphism of finite type $f:X\to B$ where each member of $\mathfrak{C}$ is a fiber of $f$, and each fiber of $f$ is an element of $\mathfrak{C}$.} 

    
    \begin{proof}[Proof of \Cref{thm_boundedness}]
        By construction there is a morphism $\cQ_{g,n}(\cX, \cX_\dm,\beta)\to \mathfrak{M}_{g,n}^{\operatorname{tw}}$. It follows from \Cref{lemma:type is bounded},
        there is a finite type substack $\sS\subseteq \mathfrak{M}_{g,n}^{\operatorname{tw}}$ and a factorization $\cQ_{g,n}(\cX, \cX_\dm,\beta)\to \sS\to  \mathfrak{M}_{g,n}^{\operatorname{tw}}$.
        There is a surjective morphism $S\to \sS$ with $S$
         an affine scheme, so it suffices to check that $\cQ_{g,n}(\cX,\cX_\dm,\beta)\times_\sS S$ is of finite type. For doing so, we can replace $S$ with any surjective morphism $S'\to S$ of finite type.         
         
         
         In what follows, we will often replace $S$ with a surjective and locally closed stratification of it, to simplify the behaviour of $\sC\to S$, the the twisted curve given by the morphism $S\to \mathfrak{M}_{g,n}^{\operatorname{tw}}$. These are always obtained in the same way: recursively, starting from a generic point $\eta$ of $S$, and spreading out. We explain what we mean in the following paragraph, and we will omit some of the details for all the other stratifications as the procedure will be very similar.
         
         Up to replacing $S$ with a sequence of surjective and locally closed embeddings $S'\to S$, we can assume that $\sC\to S$ is such that the topological type of the fibers of $\sC\to S$ is constant on connected components of $S$. Such a stratification is obtained as follows. Let $\eta\in S$ be a generic point of $S$.
         The fiber $\sC_\eta\to \eta$ has a specific topological type, which is constant in an open neighbourhood of $\eta$, as acquiring nodal singularities is a closed condition. In particular, there is $U\subseteq S$ an open neighbourhood of $\eta$ where the topological type is constant. We can then replace $S$ with $(S\smallsetminus U) \sqcup U$, and proceed inductively on $S\smallsetminus U$. This process will stop as $S$ is noetherian, so we will have replaced $S$
         with a sequence $S_1\sqcup\ldots\sqcup S_n$ of locally closed subschemes where $\cC\times_S (S_1\sqcup\ldots\sqcup S_n)\to (S_1\sqcup\ldots\sqcup S_n)$ has fibers whose topological type is locally constant.
         
         Up to replacing $S$ with a resolution of singularities $S'\to S$, we can assume that $S$ is smooth, and up to possibly replacing $S$ with a locally closed stratification followed by a finite cover $S'\to S$, the normalization $\sC^n\to\sC$ is a fiberwise normalization, i.e. for every $s\in S$ we have that $(\sC^n)_s= (\sC_s)^n$. Now, from \cite{LMB}*{Theorem 16.6} there is a finite and generically
         \'etale cover $C\to \sC^n$, and up to replacing $S$ with a stratification $S'\to S$ we can assume that $C\to \sC^n$ is generically
         \'etale along each fiber, and $C\to S$ is still a family of smooth curves. Up to further stratifying $S$ we can assume that $C\to \sC^n$ is an fppf morphism. Up to replacing $S$ with the Stein factorization of $C\to S$, we can assume that for every $s\in S$ belonging to a connected component $S_k\subseteq S$ there is a bijection between the connected components of $C_s$ and those of $C_{S_k}$. Finally, let $\sN\subseteq \sC$ be the nodes in $\sC$; up to replacing $S$ with an \'etale cover of it we can assume that $\sN\to S$ are \'etale, and they are disjoint unions of trivial gerbes. We finished our preparatory steps, we are ready to prove our main result.

         \textbf{Step 1.} The morphism $\Hom_S(\sN,\cX)\to \Hom_S(\sN,\cX\times \cX)$ is of finite type, where the morphism is induced by the diagonal $\cX\to \cX\times \cX$.

         It suffices to prove that for every $\alpha\colon B\to \Hom_S(\sN,\cX\times \cX)$, the fiber product $F$ is of finite type over $B$. A morphism $\alpha$ as above induces $\sN_B\to \cX_B\times \cX_B$, and one can check that $F$ represents the same functor as $\sN_B\times_{\cX_B\times \cX_B}\cX_B$. As $\sN$ and $\cX$ are of finite type over $S$, the morphism $\sF\to B$ is of finite type; this concludes Step 1.

         There are finitely many $g_i$ such that each irreducible component of $\sC^n$ has genus $g_i$, and if $\sC\to \cX$ has class $\beta$ there are finitely many $\beta_i$ such that an irreducible component of $\sC^n$ are such that their morphisms to $\cX$ have class $\beta_i$. 

         \textbf{Step 2.} If we know that the locus in $\mathfrak{Q}_{g_i, n_i}(\cX,\cX_\dm,\beta_j)$ where the domain curve is smooth but possibly an orbifold is of finite type for every $i,j$, then $\cQ_{g,n}(\cX,\cX_\dm,\beta)$ is of finite type.

         
         From how we stratified $S$, there are two closed embeddings $i_1,i_2:\sN\to \sC^n$ whose compositions correspond to the inclusion of the nodal locus in $\sN\to\sC$. This induces a morphism $\Hom(\sC^n,\cX)\to \Hom(\sN,\cX\times \cX)$, and we can form the fiber product \[\sF:=\Hom(\sC^n,\cX)\times_{\Hom(\sN,\cX\times \cX)}\Hom(\sN,\cX).\]
         A morphism $B\to \sF$ corresponds to a morphism
         $f:\sC^n_B\to \cX_B$ and an isomorphism $f\circ (i_1)|_B\to f\circ (i_2)|_B$. Since $\sC$ fits in a pushout diagram with $\sN\sqcup\sN\xrightarrow{i_1\sqcup i_2}\sC^n$ and $\sN\sqcup\sN\to \sN$ from \cite{alper2024artin}*{Theorem 1.8}, we have $\Hom(\sC,\cX)\cong \sF$. Then it suffices to prove that $\sF$ is of finite type. 
         
         From the previous step, we realized $\Hom(\sC,\cX)$ as a stack (namely $\sF$) which is of finite type over $\Hom(\sC^n,\cX)$. Recall that we are not interested in \textit{every} homomorphism $\sC\to \cX$, but only in those which satisfy the quasimap condition, which from \Cref{prop_qmap_gives_open_cond} and \Cref{prop_stable_qmap_is_an_open_cond} is open. Moreover,  $\sC\to \cX$ is a quasimap if and only if $\sC^n\to \cX$ is a quasimap.
         So if we know that the locus in $\Hom(\sC^n,\cX)$ where the quasimap condition is satisfied with the given invariants is of finite type, also $\cQ_{g,n}(\cX,\cX_\dm,\beta)$ will be of finite type. This concludes Step 2.

         As before, from how we stratified $\sS$, there are finitely many $g_i$ such that each irreducible component of $C$ has genus $g_i$, and if $C\to \cX$ has class $\beta$ there are finitely many $\beta_i$ such that the irreducible components of $C$ are such that their morphism to $\cX$ have class $\beta_i$.

         \textbf{Step 3.} It suffices to show that the locus in $\mathfrak{Q}_{g_i, n_i}(\cX,\cX_\dm,\beta_j)$ where the domain curve is a smooth schematic curve is of finite type for every $i,j$.

         This follows since from fppf descent and the cocycle condition there is a closed embedding \[\Hom(\sC^n,\cX)\to\Hom(C^n,\cX)\times_{\Hom(C^n\times_{\sC^n} C^n,\cX)}\Hom(C^n,\cX).\]
         Moreover, the map $\sC^n\to \cX$ is a (pointed) quasimap if and only if the corresponding morphism $C\to \cX$ is a (pointed) quasimap. 

         \textbf{End of the argument.}
         Combining \Cref{assumptions:extension of line bundle} and \Cref{thm_boundedness_embedding_in_vector_space}, we obtain that there exists an affine scheme $V=\spec(A)$ endowed with an action of a reductive group $H$ such that $j:\cX \overset{\text{open}}{\hookrightarrow} [V/H]$ and $\cX_\dm \simeq [V(\Bbbk_\chi)^{ss}/H]$ for a character $\chi$.

    Therefore, we have a well-defined functor
    \[ \mathfrak{Q}_{g,n}(\cX,\cX_\dm,\beta) \longrightarrow \mathfrak{Q}_{g,n}([V/H],[V(\Bbbk_\chi)^{ss}/H],j_*\beta).\]
    We claim that the functor above an open embedding. This would conclude the proof, as from \cite{ciocan2014stable}*{Rmk 4.2.1} the locus in $\mathfrak{Q}_{g,n}([V/H],[V(\Bbbk_\chi)^{ss}/H],j_*\beta)$ parametrizing smooth schematic curves is of finite type. 
    
    Let $(\phi:\sC\to [V/H], \Sigma)$ be a quasimap of class $j_*\beta$ over a scheme $S$. This is equivalent to giving an $H$-torsor $P\to \sC$ and a section $u:\sC\to P\times_H V$ to the projection morphism $p:P\times_H V \to \sC$, satisfying the usual hypotheses. Let $U\subset V$ be the $H$-invariant open subset such that $\cX\simeq [U/H]$, and consider the subset $U'=(P\times_H U)\times_{P\times_H V} u(\sC)$, which is open in $u(\sC)$. We deduce that the complement $Z$ of $p(U')$ is closed in $\sC$, and as the latter is proper over $S$, also the image of $Z$ in $S$ is closed. The complement of this closed subscheme is precisely the set of points $s$ such that $\sC_s$ is contained in $\cX$, which is open in $S$.
    \end{proof}
\begin{Remark}\label{remark_weight_vector_case}
    Observe that, in the case in which the $n$ marked points are schematic, one could have given a slightly different definition of stable quasimap, where the points are allowed to collide based on a weight vector $\vec{a}=(a_1,...,a_n)\in (0,1]^n$ as in \cite{Hassett_w}. Specifically, one could allow the points $p_{i_1},\ldots ,p_{i_k}$ to collide if $\sum_{j=1}^ka_{i_j}\le 1$, and instead of condition (1) in \Cref{def_stable_qmap}, require the line bundle $\omega_C(\sum a_ip_i)\otimes f^*L^{\otimes 3}$ to be nef. If we denote by $\cQ_{g,\vec{a}}(\cX,\cX_\dm,\beta)$ the resulting stack, we believe that minor modifications of the arguments given in the above sections would prove that $\cQ_{g,\vec{a}}(\cX,\cX_\dm,\beta)$ is Deligne-Mumford and proper. \end{Remark}
    \subsection{Obstruction theory}\label{subsection_obstruction_theory}
The existence of a perfect obstruction theory on $\cQ_{g,n}(\cX,\cX_\dm,\beta)$ follows from \Cref{thm_boundedness_embedding_in_vector_space} and the same arguments in \cite{ciocan2014stable}*{\S 4.5} or \cite{orbifold_qmap_theory}*{\S 2.4.5}:
\begin{Teo}
    Assume that $\cX$ satisfies \Cref{assumptions:extension of line bundle}. If $\cX$ has lci singularities and $\cX_\dm$ is smooth, the moduli stack $\cQ_{g,n}(\cX,\cX_\dm,\beta)$ carries a perfect obstruction theory.
\end{Teo}
\begin{proof}
From \Cref{thm_boundedness_embedding_in_vector_space} we can find an open embedding $\cX\hookrightarrow [\spec(A)/G]$ for a reductive group $G$. As in \cite{ciocan2014stable}*{\S 4.5} or \cite{orbifold_qmap_theory}*{\S 2.4.5}, if $\pi:\cC\to \cQ_{g,n}(\cX,\cX_\dm,\beta)$ is the universal family, with universal morphism $[u]:\cC\to \cX\to [\spec(A)/G]$, there is a principal $G$-torsor $\cP\to \cC$ and we can consider the fiber bundle $\rho:\cP\times_G\spec(A)\to \cC$ with a section $u:\cC\to \cP\times_G\spec(A)$. While the singularities of $\spec(A)$ could be arbitrary, by construction the section $u$ is such that:
\begin{enumerate}
    \item away from finitely many points it lands in $\cX_\dm$ which is smooth, and
    \item it is contained in $\cX$ which is lci.
\end{enumerate}
In particular, the complex $[u]^*\bL_{[\spec(A)/G]}$ is a complex of vector bundles concentrated in degrees $[-1,0]$. Then, proceeding exactly as in \cite{orbifold_qmap_theory}, one can prove that the complex \[
(R\pi_*[u]^*R\cH om(\bL_{[\spec(A)/G]},\cO_{[\spec(A)/G]}))^\vee
\]
 gives a perfect obstruction theory for $\cQ_{g,n}(\cX,\cX_\dm,\beta)$ relative to $\mathfrak{M}^{\operatorname{tw}}_{g,m}$, and since $\mathfrak{M}^{\operatorname{tw}}_{g,m}$ is smooth, the stack $\cQ_{g,n}(\cX,\cX_\dm,\beta)$ has an absolute perfect obstruction theory.
 \end{proof}
 The four previous subsections complete the proof of \Cref{teo_intro_cX_contains_open_proper_dm}.

\section{Applications of \Cref{teo_intro_cX_contains_open_proper_dm}}\label{section_examples_qmaps}
In this section we report two examples when one can use \Cref{teo_intro_cX_contains_open_proper_dm} to compactify the space of maps from curves to certain algebraic stacks. We first show how to use \Cref{teo_intro_cX_contains_open_proper_dm} to construct a compact moduli of fibered Calabi-Yau pairs. Then we show that if $\cX$ is a quotient of a separated Deligne-Mumford stack by a torus, and the set of properly stable points is dense in $\cX$, then $\cX$ always contains an open substack which is a proper Deligne-Mumford stack.

\subsection{Moduli of fibered Calabi-Yau pairs}
We begin by recalling that:
\begin{enumerate}
    \item there is an algebraic stack $\cX:=\cD\cP_m^{\CY}$ with a projective good moduli space parametrizing pairs $(S,cD)$ where $\dim(S)\le 2$, $D\subseteq S$ is a divisor such that the singularities of $(S,cD)$ are semi-log-canonical, $K_S+cD\sim_\mathbb{Q}0$ the sheaf $(\omega_S^{\otimes m})^{*}$ is an ample Cartier divisor, and $S$ admits a $\bQ$-Gorenstein smoothing; see \cite{blum2024good}*{Theorems 1.1, 1.2 and 3.6}.
    \item When $c<1$ and $m\gg 0$, there is an open, dense and proper Deligne-Mumford stack $\cX_\dm:=\cD\cP_m^{\operatorname{KSBA}}\subseteq \cX$ parametrizing pairs $(S,cD)$ such that $(S,(c+\epsilon)D)$ is still semi-log-canonical for some $0<\epsilon \ll 1$ \cite{blum2024good}*{Theorems 1.1 and 1.2}
    \item  For $c<1$ and $m\gg1$, if $p\colon(\sS,c\sD)\to \cX$ is the universal family, then from cohomology and base change $p_*\cO_\cS(d K_{\sS/\cX} + d(c+\epsilon)\sD)$ is a vector bundle, and $\det\left(p_*\cO_\sS(d K_{\cS/\cX} + d(c+\epsilon)\sD)\right)$ restricted to $\cX_\dm$ descends to an ample line bundle on $X_\dm$ \cite{KP17}.
    \item $\cX$ is a global quotient stack \cite{ascher2023moduli}.
    \item For $m\gg 1$, the algebraic stack $\cX$ admits a projective good moduli space $q:\cX\to X$ \cite{blum2024good}*{\S 6.4}.
    \item The map $q$ identifies two pairs $(S_1,cD_1)$ and $(S_2,cD_2)$ if and only if they are $S$-equivalent \cite{blum2024good}. 
\end{enumerate}
{We aim at proving \Cref{teo_intro_kodairadim1}, whose content we recall below for the convenience of the reader.
\begin{Teo}
    Set $\cX:=\cD\cP^{\CY}_m$ and $\cX_\dm=\cD\cP^{\operatorname{KSBA}}_m$. Then the assumptions of \Cref{teo_intro_cX_contains_open_proper_dm} apply for the inclusion $\cX_\dm\subseteq \cX$. In particular:
    \begin{enumerate}
        \item the stack $\cQ_g(\cX,\cX_\dm,\beta)$ compactifies the space of maps $\pi\colon(Y,cD)\to C$ with fibers in $\cX$ such that:
        \begin{itemize}
            \item[(Q)] the curve $C$ is smooth and the generic fiber of $\pi$ has klt singularities,
            \item[(S)] either $\omega_C$ is ample, or not all the fibers of $\pi$ are $S$-equivalent,
            \item[(N)] the family $\pi$ comes from a map $C\to \cX$ of class $\beta$ from a curve of genus $g$.
        \end{itemize}
        \item the boundary of $\cQ_g(\cX,\cX_\dm,\beta)$ parametrizes families $\pi:(\sY,c\sD)\to \sC$ of pairs in $\cX$ with fibers in $\cX$, fibered over a twisted curve $\sC$, such that:
        \begin{itemize}
            \item[(Q)] the set $\Delta:=\{p\in \sC:(\sY_p,(c+\epsilon)\sD_p)$ does \underline{not} have semi-log-canonical singularities for any $0<\epsilon \ll 1\}$ is a finite union of smooth points $\sC$,
            \item[(S)] if $\sR\subseteq \sC$ is an irreducible component such that $\deg(\omega_\sC |_\sR)< 0$, then not all the fibers of $\pi|_\sR\colon(\sY|_\sR,c\sD|_\sR)\to \sR$ are $S$-equivalent; whereas if $\deg(\omega_\sC |_\sR)=0$, then not all fibers of $\pi|_\sR$ are isomorphic, and
            \item[(N)] the family $\pi$ comes from a map $\sC\to \cX$ of class $\beta$ from a twisted curve of genus $g$.
        \end{itemize}
    \end{enumerate}
\end{Teo}}
\begin{proof}[Proof of \Cref{teo_intro_kodairadim1}]{To prove that $\cX_\dm$ is the semistable locus relative to $X$ for a line bundle on $\cX$ we can use \Cref{cor_when_X_dm_is_ss_locus_for_a_lb} with $\det\left(p_*\cO_\sS(d K_{\cS/\cX} + d(c+\epsilon)\sD)\right)$. The only other statement one needs to check is that if $f\colon\sC\to \sX$ is stable, then it satisfies condition (2) part (S) in \Cref{teo_intro_kodairadim1}. This follows from \Cref{lemma_equivalent_conditions_for_stability}: since the generic points of $\sC$ map to $\sX_\dm$, if $f$ is stable then any irreducible component $\sR\subseteq \sC$ such that $\deg((\omega_\sC)|_\sR)<0$ must map quasi-finitely to $X$, and so from point (6) above there must be two fibers of $\pi|_\cR$ not $S$-equivalent. Assume $\sR$ is such that $\deg((\omega_\sC)|_\sR)=0$. Since $f$ is a quasimap, the generic point of $\sR$ maps to a point $x\in\cX_\dm$, and $f$ is stable unless all the points of $\sR$ map to $x$. }    
\end{proof}
\subsection{Toric quotients}\label{section_toric_quotients}
We prove that if $\cX$ is a quotient of a Deligne-Mumford stack by a split torus, and there is an open dense $U\subseteq X$ such that $\cX\times_XU$ is a Deligne-Mumford stack, then $\cX$ always contains an open substack which is a proper Deligne-Mumford stack. So for such stacks, the enlargement $\cX\subseteq \widetilde{\cX}$ of \Cref{thm_intro_simplified} is not needed{: if $\cX$ is a global quotient, \Cref{teo_intro_cX_contains_open_proper_dm} applies}.
\begin{Teo}\label{teo_toric_stack_has_open_dense_dm}
    Let $\cX$ be an algebraic stack with a good moduli space $\cX\to X$ which is separated. Assume that:
    \begin{itemize}
        \item there is a morphism $\cX\to \cB \Gm^n$,
    representable in Deligne-Mumford stacks, and
    \item there is a dense open $U\subseteq  X$ such that $\cX\times_XU$ is Deligne-Mumford.
    \end{itemize} Then there are finitely many hyperplanes $H_1,...,H_m$ in the space $\mathbf{X}(\Gm^n)_\bQ$ of rational characters of $\Gm^n$ such that, if $\mu \notin \bigcup H_i$, then $\cX(\Bbbk_\mu)^{ss}_X$ is a separated Deligne-Mumford stack, with coarse moduli space projective over $X$ and such that $\cX\times_XU\subseteq \cX(\Bbbk_\mu)^{ss}_X$.
\end{Teo}

\begin{proof}
We will denote by $G$ a group which is a central extension of $\Gm^r$ by a finite group $F$. 

\textbf{Affine case.} We will assume that $\cX=[\spec(A)/G]$, and that $\spec(A)$ has a $G$-fixed point. In particular, our assumptions guarantee that there is a morphism $\psi:G\to \Gm^n$ with finite kernel.  

    \textbf{Step 1.} We can assume $G=\Gm^r\times F$.
    
    From \cite{brion2015extensions}, there is a finite subgroup $F<G$ and a surjective morphism $\Gm^r \rtimes F \to G$ with finite kernel. As $\Gm^r$ is contained in the center of $G$, the product is direct. In particular, there is a morphism $[\spec(A)/\Gm^r\times F]\to [\spec(A)/G]$ which is separated and a gerbe, so if $\mu$ is a character of $\Gm^n$, then $[\spec(A)/\Gm^r\times F](\Bbbk_\mu)^{ss}$ satisfies the condition of the theorem if and only if $[\spec(A)/G](\Bbbk_\mu)^{ss}$ does. In particular, we can assume that $G=\Gm^r\times F$.

    \textbf{Step 2.} We can assume that $F=\{1\}$.

    Indeed, up to post-composing $\cX\to \cB\Gm^n$ with the map $\cB \Gm^n\to \cB\Gm^n$ induced by $(t_1,\ldots ,t_n)\mapsto (t_1^N,\ldots ,t_n^N)$ for $N\gg 1$, we can assume that the morphism $\psi:\Gm^r\times F\to \Gm^n$ is trivial if restricted to $F$.  In particular, we can factor $[\spec(A)/\Gm^r\times F]\to \cB\Gm^n\times \spec(A^{\Gm\times F})$ as
    \[[\spec(A)/\Gm^r\times F]\to [\spec(A^F)/\Gm^r]\to \cB\Gm^n\times \spec(A^{\Gm\times F})\]
    and again, if we can find such a $\mu$ for $[\spec(A^F)/\Gm^r]$, the same $\mu$ will also work for $[\spec(A)/\Gm^r\times F]$. 

    \textbf{Step 3.} We can assume that $\spec(A)=\bA^N$.

    Indeed, there is an equivariant closed embedding $i:\spec(A)\hookrightarrow \bA^N$. From the affine Hilbert-Mumford criterion, since $i$ is a closed embedding, we have that $i^{-1}(\bA^N(\Bbbk_\mu)^{ss}) =\spec(A)(\Bbbk_\mu)^{ss}$. Since $\Gm^r$ is reductive, there is a closed embedding $\spec(A^{\Gm^r})\hookrightarrow \bA^N/\!\!/\Gm^r$ where we denoted by $\bA^N/\!\!/\Gm^r$ the good moduli space of $[\bA^N/\Gm^r]$. So if we find such a $\mu$ for $[\bA^N/\Gm^r]$, the same $\mu$ will also work for $[\spec(A)/\Gm^r]$.

    \textbf{The $\bA^N$ case.} Up to performing a change of coordinates, we can assume that the action is diagonal, given by $N$ characters $\chi_1,\ldots,\chi_N$.
     
    Given a character $\widetilde{\mu}$, recall that:
    \begin{itemize}
        \item from \cite{Hos}*{Proposition 2.5} a point $p\in \bA^N$ is $\widetilde{\mu}$-stable
    if, for every one-parameter subgroup $\lambda:\Gm\to \Gm^r$ such that $\lim_{t\to 0}\lambda(t)p$ exists, we have that $<\lambda, \widetilde{\mu}>\le 0$, and
    \item given $\lambda\in \mathbf{X}(\Gm^r)^*$, if a point $p$ is such that $\lim_{t\to 0}\lambda(t)p$ exists and if $<\lambda,\chi_j><0$, then the $j$-th coordinate of $p$ is 0.
    \end{itemize} 
    
    Since the generic stabilizer is finite, the morphism $\Xi:\Gm^r\to \Gm^N$ that sends $g\mapsto (\chi_1(g),\ldots, \chi_N(g))$ has finite kernel. In particular, the point $(1,\ldots,1)$ is stable, as there is no one-parameter subgroup $\lambda$ such that $\lim_{t\to 0}\lambda(t)p$ exists. Moreover,
    if we choose $\widetilde{\mu}$ such that the hyperplane in
    $\mathbf{X}(\Gm^r)^*$
    given by $<\cdot, \widetilde{\mu}>=0$ is different from $<\cdot, \chi_i>=0$ for every $i$, then the stable locus and the semistable locus agree. Indeed, if $p$ is a point such that,
    for a given $\lambda$, we have that $<\lambda,\widetilde{\mu}>=0$ and $\lim_{t\to 0}\lambda(t)p$ exists, then $\lambda$ is away from the hyperplanes $<\cdot, \chi_i> = 0$. In particular, we can slightly perturb $\lambda$ in a way such that $\lim_{t\to 0}\lambda(t)p$ still exists, but $<\lambda,\widetilde{\mu}> \gg 0$, so $p$ was not semistable.
    
    Finally, since $\psi:G\to\Gm^n$ has finite kernel, the morphism
    $\Psi:\mathbf{X}(\Gm^n)_{\mathbb{Q}}\to \mathbf{X}(\Gm^r)_{\mathbb{Q}}$ induced by $\psi:G\to \Gm^n$ is surjective.
    So as long as we take $\mu$ away from $\bigcup\Psi^{-1}(\chi_i)$,
    the character $\widetilde{\mu}=\Psi(\mu)$ will work. 

    \textbf{General case.} Since $\cX\to \cB\Gm^n$ is representable in Deligne-Mumford stacks, from \cite{di2024effective}*{Corollary 4.7 \& Theorem 2.19} the stabilizers $G$ of the geometric points are central extensions of a split torus $\Gm^r$ by a finite group $F$.
    In particular, for every $p\in X$ there is an \'etale neighbourhood of $p$ of the form $\spec(A^G)\to X$ such that the following diagram is cartesian:
    \[
    \xymatrix{[\spec(A)/G]\ar[r] \ar[d] & \cX\ar[d] \\ \spec(A^G)\ar[r] & X}
    \]
    and there is a homomorphism $\psi:G\to \Gm^n$ with finite kernel. So for each neighbourhood as above we can find a finite set of hyperplanes of $\mathbf{X}(\Gm^n)_{\mathbb{Q}}$ as in the statement of the theorem.
    Now the desired result follows since $X$ is quasicompact, so we can cover $X$ with finitely many \'etale charts as above. 
\end{proof}
\begin{example}
We give an application of \Cref{teo_toric_stack_has_open_dense_dm} to study stable maps to a specific quotient. Consider the action of $\Gm$ on $\bA^2$ with weights 1 and $-1$. The good moduli space $[\bA^2/\Gm]$ is $\bA^1$, and the quotient map $[\bA^2/\Gm]\to \bA^1$ is an isomorphism away from a point (which we set to be 0). We can consider the stack $\sP^1$ obtained by replacing the chart in $\bP^1$ given by $x_0\neq 0$ with $[\bA^2/\Gm]$.
More specifically, this is the gluing of $\bA^1_t$ and $[\bA^2_{x,y}/\Gm]$ along $\bA^1\smallsetminus\{0\}$ and $[\bA^2_{x,y}\smallsetminus\{xy=0\}/\Gm]$ via the map that sends $t\mapsto \frac{1}{xy}$. The stack $\sP^1$ is an algebraic stack with $\bP^1$ with a good moduli space. It is smooth and it has a Zariski open cover where it is a quotient of a separated scheme by a torus.
Then from the main result of \cite{di2024effective}, it is globally a quotient of a separated Deligne-Mumford stack by a torus. Indeed in \textit{loc. cit.} we give a criterion for when a smooth algebraic stack with a good moduli space admits such a presentation, and the criterion is Zariski local on the good moduli space. Therefore \Cref{teo_toric_stack_has_open_dense_dm} applies, and one can check explicitly that $\sP^1$ contains $\bP^1$ as an open substack, since the quotient map $[\bA^2/\Gm]\to \bA^1$ admits a section given by $[\bA^2\smallsetminus\{x=0\}/\Gm]\to \bA^1$. This extends to a section $\bP^1\to \sP^1$. 

Consider the map $\bA^1_t\to \bA^2_{x,y}$, $x\mapsto t-a$ and $y\mapsto t-b$ for $a,b\in k$. This gives a map $\bA^1\to [\bA^2/\Gm]$ which, if composed to $[\bA^2/\Gm]\to \bA^1\hookrightarrow\bP^1$, compactifies to a map $\phi:\bP^1\to\sP^1$ whose composition with the good quotient of $\sP^1$ gives $\bP^1\to \sP^1\to \bP^1$ of degree 2, and not ramified at 0 (the point on $\bP^1$ corresponding to a polystable and not stable point of $\sP^1$).  If we take as open proper Deligne-Mumford of $\sP^1$ the locus in $[\bA^2/\Gm]$ where $x\neq0$, and we denote by $C$ the domain of $\phi$, then we have a unique point in $C$ which does not map to $\sP^1_\dm$. If one takes the closure of the locus of maps $\bP^1\to \sP^1$ constructed as above, it is straightforward to check that the nodal curves that one finds on the boundary are either $D$ the nodal union of two $\bP^1$s, with the composition $D\to \sP^1\to \bP^1$ being unramified over $0$, or the nodal union $D$ of three $\bP^1$s, with the two $\bP^1$s which have a single node mapping finitely to the good quotient of $\sP^1$, and the $\bP^1$ with two nodes mapping to $0$. The map $\phi:D\to \sP^1$ will still be a quasimap, so $\phi^{-1}(\sP^1\smallsetminus\sP^1_\dm)$ will still be a unique smooth point.
\end{example}

\section{Extended weighted blow-ups}\label{section_extended_blowups}
In this section we first introduce extended weighted blow-ups, which are a mild generalization of weighted blow-ups, and are more amenable for compactifying moduli spaces of maps to algebraic stacks. Then we prove \Cref{thm_you_can_embed_a_Stack_in_one_with_open_dm_by_taking_deformations_to_nc}, which allows to embed any algebraic stack $\cX$ with a good moduli space $\cX\to X$ and a properly stable point in an algebraic stack $\widetilde{\cX}$ with the same good moduli space $\widetilde{\cX}\to X$, and such that $\widetilde{\cX}\to X$ contains an open substack which is proper and Deligne-Mumford.
\begin{Teo}\label{thm_you_can_embed_a_Stack_in_one_with_open_dm_by_taking_deformations_to_nc}
    Let $\cX$ be an algebraic stack with a good moduli space $p:\cX\to X$, and with an open dense $U\subseteq X$ such that $\cX\times_XU$ is Deligne-Mumford. Then there is an algebraic stack $\widetilde{\cX}$ with an open embedding $i:\cX\hookrightarrow \widetilde{\cX}$ such that:
    \begin{enumerate}
        \item $\widetilde{\cX}$ has a good moduli space which is isomorphic to $X$, and $i$ induces an isomorphism on good moduli spaces, 
        \item there is a morphism $\pi:\widetilde{\cX}\to \cX$ which is an isomorphism over the open subset of $\cX$ given by $(\pi\circ p)^{-1}(U)$,
        \item the morphism $\pi\circ i$ is isomorphic to the identity,
        \item there is a line bundle $\cL_{DM}$ on $\widetilde{\cX}$ such that $\widetilde{\cX}(\cL_{DM})^{ss}_X$ is  Deligne-Mumford and proper over $X$,
        \item if $\cX$ is a global quotient by a reductive group, then $\widetilde{\cX}$ can be chosen to be a global quotient by a reductive group.
    \end{enumerate}
    In particular, the stack $\widetilde{\cX}$ satisfies \Cref{assumptions:extension of line bundle} if $\cX=[W/G]$ for $G$ reductive and $X$ is projective.
\end{Teo}
In particular, we can combine \Cref{thm_you_can_embed_a_Stack_in_one_with_open_dm_by_taking_deformations_to_nc} with the results of \Cref{sec:qmaps} to get the following
\begin{Cor}
    Let $\cX$ be an algebraic stack with a good moduli space $p:\cX\to X$, and with an open dense $U\subseteq X$ such that $\cX\times_XU$ is Deligne-Mumford. Assume that $\cX=[W/G]$, where $G$ is a reductive group. Then there is a schematically dense open embedding $\cX\subseteq \widetilde{\cX}$ as in \Cref{thm_you_can_embed_a_Stack_in_one_with_open_dm_by_taking_deformations_to_nc}, and there is an algebraic stack $\cQ_{g,n}(\widetilde{\cX}, \widetilde{\cX}(\cL_\dm)^{ss}_X,\beta)$ parametrizing stable quasimaps $\sC\to \widetilde{\cX}$ of genus $g$ and class $\beta$. 
    \end{Cor}
\begin{Remark}
    Recall locus $\cS:=\{x\in X:p^{-1}(x)$ is a Deligne-Mumford gerbe$\}$ is open in $X$ from \cite{ER21}*{Proposition 2.6}. So if $X$ is irreducible, if there is a point in $\cS$ then $\cS$ is a dense open subscheme of $X$.
\end{Remark}
\subsection{Extended weighted blow-ups}
In this subsection we define \textit{extended weighted blow-ups}.
On a first approximation, an extended blow-up of an algebraic stack $\cX$ along an ideal $\cI$ is the $\Gm$-quotient of the deformation of $\cX$ to the weighted normal cone determined by $\cI$.
More specifically, if we denote by $\cZ\subset\cX$ the closed substack defined by $\cI$, the deformation to the weighted normal cone of $\cX$ is (a subset of) the weighted blow-up of $\cX\times\bA^1$ along $\cZ\times\{0\}$, where the parameter $T$ of $\bA^1$ is defined of degree $-1$. The induced $\Gm$-action on $\cX\times\bA^1$ extends then to
the deformation of the normal cone, and the extended weighted blow-up is the $\Gm$-quotient of the latter.

More formally, we begin with the following

\begin{Def}[\cite{QR}*{Definition 3.1.1}]\label{def_weighted_embedding}
   Let $\cX$ be an Artin stack. A weighted embedding $\cY_{\bullet} \hookrightarrow \cX$ is defined by a sequence of closed embeddings $\{ \cY_n=V(\cI_n) \hookrightarrow \cX \}_{n \ge 0}$ such that: \begin{itemize}
        \item $\cI_0 \supset \cI_1 \supset \dots \supset \cI_n \supset \dots$
        \item $\cI_n \cI_m \subset \cI_{n+m}$
        \item Locally in the smooth topology on $X$, there exists a sufficiently large positive integer $d$ such that for all integers $n \ge 1$, \[\cI_n=\left(\cI_1^{l_1}\cI_2^{l_2} \cdots \cI_d^{l_d} \; : \; l_i \in \bN, \ \sum_{i=1}^d il_i=n\right)\] in which case, we say $\cI_\bullet$ is generated in degrees $\le d$. 
    \end{itemize} 
    The last condition can be understood as a needed condition for $\bigoplus_n \cI_n$ to be of finite type.
    Furthermore, we set $\cI_n=\cO_{\cX}$ for $n\leq 0$ and we call the sequence of ideals $\{\cI_n\}_{n\in \bZ}$ a \textit{weighted ideal sequence}. 
\end{Def} 

Consider then the graded $\cO_\cX$-algebra $\cA:=\oplus_{n\in\bZ} \cI_n$. We will use an auxiliary variable $T$ to denote the degree, so in particular we will write an element of $\cA$ as $\sum_{j=-k}^ka_jT^j$, where $a_j$ is an element of $\cI_j$ in degree $j$. Observe that:
\begin{enumerate}
    \item there is a morphism $\pi^\#:\cO_\cX\to \cA$, which sends $\cO_\cX$ to the degree zero component;
    \item the previous morphism has a left inverse $i^\#:\cA \to \cO_\cX$ defined by $\sum_{j=-k}^ka_jT^j\mapsto \sum_{j=-k}^ka_j$.
\end{enumerate}
Moreover, we have the following.
\begin{Lemma}\label{lm:action}
    There is a strict $\Gm$-action on $\spec_{\cO_\cX}(\cA)$ over $\cX$ in the sense of \cite{Rom05}*{Definition 1.3}. In particular, there is a quotient $[\spec_{\cO_\cX}(\cA)/\Gm] \to \cX$ whose formation commutes with base change.
\end{Lemma}
\begin{proof}
    The grading of $\cA$ induces a co-action
    \[ \cA \longrightarrow \cA[u^{\pm 1}],\quad aT^j\longmapsto aT^ju^j.\]
    We can then define an action $\mu\colon\Gm \times \spec_{\cO_\cX}(\cA) \to \spec_{\cO_\cX}(\cA)$ as follows: \begin{center} the object $(\lambda \in \cO_S^*(S), \xi\colon S\to\cX, f\colon \xi^*\cA \to \cO_S)$ is mapped to $(\xi, \xi^*\cA \to \xi^*\cA[u^{\pm 1}] \overset{{u=\lambda}}{\longrightarrow} \xi^*\cA \overset{f}{\to} \cO_S).$\end{center}
    Moreover, as the automorphism group of $(\lambda,\xi,f)$ coincides with the one of $(\xi,f)$ and $\xi$, the morphism between automorphism groups is just the identity; in this way, we get a morphism of algebraic stacks.

    To check that the morphism above defines a strict action, we need to check that the identity (regarded as a $2$-morphism) makes the two following diagrams of algebraic stacks $2$-commute:
    \[
    \begin{tikzcd}
        \Gm \times \Gm \times \spec_{\cO_\cX}(\cA) \ar[r, "m\times \id"] \ar[d, "\id\times\mu"] & \Gm \times \spec_{\cO_\cX}(\cA) \ar[d, "\mu"] \\
        \Gm \times \spec_{\cO_\cX}(\cA) \ar[r, "\mu"] & \spec_{\cO_\cX}(\cA)
    \end{tikzcd} \quad
    \begin{tikzcd}
        \Gm \times \spec_{\cO_\cX}(\cA) \ar[r, "\mu"] & \spec_{\cO_\cX}(\cA) \\
        \spec_{\cO_\cX}(\cA). \ar[u, "1\times\id"] \ar[ur, "\id"]
    \end{tikzcd}
    \]
    This is pretty straightforward, and it follows from the fact that the action morphism is a morphism over $\cX$, where the $2$-commutativity is strict (meaning that the $2$-morphism making the diagram $2$-commute is actually the identity); we omit the details. Then from \cite{Rom05}*{Theorem 4.1} it follows that there exists an algebraic stack $[\spec_{\cO_\cX}(\cA)/\Gm]\to \cX$ whose formation commute with base change.
\end{proof}
Given \Cref{lm:action}, we observe also the following:
\begin{enumerate}
    \setcounter{enumi}{2}
    \item if we denote by $\pi$ the morphism induced by $\pi^\#$ on sections, then $\pi$ is $\Gm$-equivariant and the induced morphism $[\spec_{\cO_\cX}(\cA)/\Gm]\to X$ is a good moduli space as the $\Gm$-invariant functions are those on degree zero;
    \item if we denote by $i$ the morphism $\cX\to [\spec_{\cO_\cX}(\cA)/\Gm]$ induced by $i^\#$, this is an open embedding \cite{QR}*{Proposition 4.3.4};
    \item there is a morphism $ [\spec_{\cO_\cX}(\cA)/\Gm]\to \cX\times[\bA^1/\Gm]$, given by the inclusion $\bigoplus_{n\le 0}\cO_\cX[T^{-1}]\to \cA$, which induces an isomorphism of $i(X)$ with the preimage of $X\times \{1\}$ and of its complement with the complement of $\{T^{-1}\}=0$; in particular, we have that $[\spec_{\cO_\cX}(\cA)/\Gm]\smallsetminus i(X)$ is Cartier.
    \item from \cite{QR}*{Remark 3.2.5}, we have that $[\spec_{\cO_\cX}(\cA)\smallsetminus V(\cA^+)/\Gm]$ is the weighted blow-up of $\cX$ along $\cI_\bullet$.
\end{enumerate}
\begin{Def}\label{def_ext_w_blowup} Given a weighted ideal sequence $\{\cI_n\}$ as above, we define \[
\EB_{\cI_\bullet}\cX:=[\spec_{\cO_\cX}(\cA)/\Gm]
\]
     the \textit{extended weighted blow-up} of $\cX$ along $\cI_\bullet$. When $\cI_n = \cI_1^n$ for $n\ge 1$, we call $\EB_{\cI_\bullet}\cX$ simply an \textit{extended blow-up} and we denote it by $\EB_{\cI}\cX$.
\end{Def}
\begin{Remark}
    An extended weighted blow-up $\pi:\EB_{\cI^\bullet}\cX\to \cX$ is such that $\pi_*\cO_{\EB_{\cI^\bullet}\cX}=\cO_{\cX}$. Indeed, one can check this smooth locally over $\cX$ so we can assume that $\cX$ is a scheme, in which case it follows from point (3) above.
\end{Remark}

\begin{example}[Extended blow-up of $0\in \bA^2$]\label{example_ext_blowup_of_origin_in_A2} In this example we work out explicitly the extended blow-up of $0\in \bA^2$.
    The deformation to the normal cone of the origin $V(x,y)\subseteq \bA^2$ can be identified with $\spec(k[x,y,T,X,Y]/(XT-x, YT-y)$ where:
    \begin{itemize}
        \item $X,Y$ are the generators of the maximal ideal in degree 1,
        \item $T$ has degree $-1$,
        \item $x,y$ have degree 0,
        \item the map $X\to \spec(I^\bullet)$ is given by sending $T\mapsto 1$, and 
        \item the map $\spec(I^\bullet)\to X$ is given by the inclusion of the degree 0 component.
    \end{itemize}
    Therefore we can identify it with $\bA^3=\spec(k[X,Y,T])$, which has a $\Gm$-action with weights $(1,1,-1)$.
    The map $\bA^3\to \bA^2$ given by $(a,b,u)\mapsto (au,bu)$
    induces the good moduli space morphism \[\pi:[\bA^3/\Gm]\to \bA^2\] which is the extended blow-up of $0\in \bA^2$. The section $i$ of point (4) is given by the locus where $T=1$ in $[\bA^3/\Gm]$. As for point (6), we can identify $[\{(a,b,t):(a,b)\neq (0,0)\}/\Gm]\hookrightarrow[\bA^3/\Gm]$ with the blow-up of the origin in $\bA^2$ as in \cite{Inc}*{\S 2.2}.
    
    Similarly, when a group $G$ acts on $\bA^2$, we can extend its action to $\bA^3$,
    by acting trivially on the $T$ component. So as before one has two morphism $[\bA^2/G]\to[\bA^3/G\times \Gm] $ and $[\bA^3/G\times \Gm]\to [\bA^2/G]$ whose composition is the identity, where the first is the complement of a Cartier divisor, and $[\bA^3/G\times \Gm]$ contains as an open the blow-up of $\cB G$ in $[\bA^2/G]$. 
\end{example}
\begin{Def}\label{def_gen_blowup_comes_with_a_cartier_divisor} Given a stack $\cX$ and a weighted ideal sequence $\cI^\bullet$, there is an effective Cartier divisor $\cE\subseteq \EB_{\cI^\bullet}X$ which restrict to the exceptional divisor of the blow-up, and which in the examples given above it agrees with the vanishing locus of $T^{-1}$. We call such an effective Cartier divisor the \textit{exceptional divisor}. Observe that the exceptional divisor of an extended weighted blow-up $\EB_{\cI^\bullet}X$ is the scheme-theoretic closure of the (usual) exceptional divisor of the weighted blow-up contained (as an open) in $\EB_{\cI^\bullet}X$.
\end{Def}
\begin{Remark}
    One can check that if $\cZ\subseteq \cX$ is a smooth closed substack of $\cX$ with ideal $\cI$, and $\cI^\bullet = \{\cI^n\}$ then the extended blowup $\EB_{\cI^\bullet}\cX$ is smooth.
\end{Remark}
\subsection{Proof of \Cref{thm_you_can_embed_a_Stack_in_one_with_open_dm_by_taking_deformations_to_nc}}
We begin with the following.
\begin{Lemma}\label{lemma_ext_blowup_is_relative_ss_locus}
    Let $\cX$ be an algebraic stack and $\cI_\bullet:=\{\cI_n\}_{n\in \bZ}$ a weighted ideal sequence of $\cO_\cX$. Let $\theta:\Gm\to \Gm$ be the character $t\mapsto t^{-1}$.
    Then $(\EB_{\cI_\bullet} \cX)(\Bbbk_{\theta} )^{ss}_{\cX}$
    is the weighted blow-up of $\cX$ along $\cI_\bullet$. Moreover, if $\pi:\cX\to X$ is a good moduli space and $\cI_n=\cI_1^n$, then $(\EB_{\cI_\bullet} \cX)(\Bbbk_{\theta} )^{ss}_{X}\cong \Bl^\pi_\cI\cX$, where the latter is a saturated blow-up (see \cite{ER21}*{Definitions 3.1 and 3.2}).
\end{Lemma}
\begin{proof}If we denote by $p:\EB_{\cI_\bullet}\cX\to \cX$ the projection, observe that \[
    \pi_*\Bbbk_{\theta}^{\otimes n} = \cI^n.
    \]
    This can be checked smooth locally over $\cX$, where it suffices to observe that the sections of $\spec_{\cO_\cX}(\bigoplus_{n\in \bZ}\cI^n)$ which are $\Bbbk_{\theta}^{\otimes n}$-semiinvariant are homogeneous of degree $n$.
    The first part now follows since $(\EB_{\cI_\bullet}\cX)(\Bbbk_{\theta}  )^{ss}_{\cX} = [\spec(\bigoplus_{n\in \bZ} \cI_n)\smallsetminus V(<\bigoplus_{n>0}\cI_n>)/\Gm]$ and from \cite{QR}*{Remark 3.2.4 and \S 1.1} the latter is the desired weighted blowup.

    
     For the moreover part, observe that $(\EB_{\cI} \cX)(\Bbbk_{\theta} )^{ss}_{X}$ is the locus of $\EB_{\cI_\bullet} \cX$ given by the complement of \[V((\pi\circ p)^{-1}(\pi\circ p)_*(\bigoplus_{n\in \bZ} \Bbbk_{\theta}^{\otimes n})) = V((\pi\circ p)^{-1}\pi_*(\bigoplus_{n\in \bZ} \cI^n))=V(p^{-1}\pi^{-1}\pi_*(\bigoplus_{n\in \bZ} \cI^n))=:\cZ.\]
Since $\pi^{-1}\pi^*\cI_n\subseteq \cI^n$, we have that $V(\bigoplus_{n\ge 1}\cI_n)\subseteq \cZ$ so the latter is the locus in the weighted blow-up of $\cI_\bullet$ given by the complement of $V(\pi^{-1}\pi_*(\bigoplus_{n\in \bZ} \cI_n))$, namely $\Bl^\pi_\cI\cX$.
\end{proof}
\begin{proof}[Proof of \Cref{thm_you_can_embed_a_Stack_in_one_with_open_dm_by_taking_deformations_to_nc}] By the main theorem in \cite{ER21} there is a sequence of saturated blow-ups $\cX_n\to \cX_{n-1}\to \ldots \to \cX_1\to \cX$ such that $\cX_n$ is Deligne-Mumford. As taking an extended blow-up commutes with open embeddings, it suffices to prove by induction on $n$ that there is an algebraic stack $\widetilde{\cX}_j$ such that:
\begin{enumerate}
    \item there is an open embedding $i_j:\cX\to\widetilde{\cX}_j$ with a projection $\pi_j:\widetilde{\cX}_j\to \cX$ inducing isomorphisms on good moduli spaces and such that $\pi_j\circ i_j$ is isomorphic to the identity,
    \item if $U\subseteq X$ consists of the open subset of properly stable points in $X$, then $\pi_j$ is an isomorphism on the preimage of $U$,
    \item there is a morphism $\xi_j:\widetilde{\cX}_j\to \cB\Gm^{\oplus j}$ and a character $\theta_j:\Gm^{\oplus j}\to \Gm$ such that if we denote by $\cL_j:=\xi_j^*\Bbbk_{\theta_j}$, the semistable locus for $\cL_j$ over $X$ is isomorphic to $\cX_j$.
\end{enumerate}
For the case $n=1$ we can take the extended blow-up corresponding to the saturated blow-up $\cX_1\to \cX$, and apply \Cref{lemma_ext_blowup_is_relative_ss_locus}: all the points above follow from the properties of extended blow-ups and from \Cref{lemma_ext_blowup_is_relative_ss_locus}.

Assuming that the desired result holds for $j$, we show it holds for $j+1$. Let $\cZ_j\subseteq \cX_j$ the closed locus we would blow-up to obtain $\cX_{j+1}$, and let $\widetilde{\cZ}_j$ be its closure in $\widetilde{\cX}_j$. As $\cX_j$ is open in $\widetilde{\cX}_j$, the extended blow-up $\tau:\EB_{\widetilde{\cZ}_j}\widetilde{\cX}_j\to \widetilde{\cX}_j$ restricts to $\EB_{\cZ_j}\cX_j$. From \Cref{lemma_ext_blowup_is_relative_ss_locus} there is a morphism $\EB_{\cZ_j}\cX_j\to \cB \Gm$ and a character $\theta_{j+1}$ such that the semistable locus over the good moduli space $X_j$ of $\cX_j$ for $\Bbbk_{\theta_{j+1}}$ is $\cX_{j+1}$ . The morphism $\EB_{\cZ_j}\cX_j\to \cB \Gm$ comes from the $\Gm$-quotient presentation of the definition of extended blow-up, so it extends to a morphism $\EB_{\widetilde{\cZ}_j}\widetilde{\cX}_j\to \cB \Gm$. There is a morphism $\iota:\widetilde{\cX}_j\to \EB_{\widetilde{\cZ}_j}\widetilde{\cX}_j$ from the properties of extended blow-ups, and it is straightforward to check that points (1) and (2) follow by taking $i_{j+1}:=\iota\circ i_i$ and $\pi_{j+1} = \pi_j\circ \tau$. Point (3) instead follows from \Cref{prop_stable_locus_of_pi*L^n_otimes_G}, we explain how. By induction we have a line bundle $\cL$ on $\widetilde{\cX}_j$ whose semistable locus over $X$ is $\cX_j\subseteq \widetilde{\cX}_j$, the $j$-th step of the algorithm of Edidin and Rydh.
We know that the $(j+1)$-th
step $\cX_{j+1}\to \cX_j$ is a saturated blow-up, which by \Cref{lemma_ext_blowup_is_relative_ss_locus} is the relative semistable locus, over the good moduli space $X_j$ of $\cX_j$, for a line bundle $\cM$ of the extended blow-up $\EB_{\cZ_{j+1}}\cX_{j+1}$. Then \Cref{prop_stable_locus_of_pi*L^n_otimes_G} allows us to take both semistable loci simultaneously, so $\cX_{j+1}$ will be the $\pi^*\cL^{\otimes m}\otimes \cM$-semistable locus of $\widetilde{\cX}_{j+1}$ over $X$, which proves (3). 

Finally, as the weighted blow-up $\EB_\cI \cX$ admits a representable (in fact, affine) morphism $\EB_\cI\to \cX\times \cB \Gm$, if $\cX$ is a global quotient then also $\widetilde{\cX}_i$ is a global quotient for every $i$.
\end{proof}




\begin{comment}
\subsubsection{Deformation to the normal cone}\label{subsubsection_def_to_normal_cone}
Recall that, given a closed subscheme $Z\subseteq X$ with ideal sheaf $I$, one can construct the algebra $I^\bullet:=\bigoplus_{n\in \bZ}I^n$, where we fix the convention that $I^n=\cO_X$ for $n\ge 0$. We will use an auxiliary variable $T$ to denote the degree, so in particular we will write an element of $I^\bullet$ as $\sum_{j=-k}^ka_jT^j$, where $a_j$ is an element of $I^j$ in degree $j$. Observe now that:
\begin{enumerate}
    \item there is a morphism $\pi^\#:\cO_X\to I^\bullet$, which sends $\cO_X$ to the degree 0 component,
    \item the previous morphism has a left inverse $i^\#:I^\bullet \to \cO_X$ defined by $\sum_{j=-k}^ka_jT^j\mapsto \sum_{j=-k}^ka_j$,
    \item the grading induces an action of $\Gm$ on $\spec_{\cO_X}(I^\bullet)$, and from \cite{QR}*{Rmk. 3.2.5}, we have that $[\spec_{\cO_X}(I^\bullet)\smallsetminus V(I^\bullet_+)/\Gm]$ is the blow-up of $X$ along $I$,
    \item if we denote by $\pi$ the morphism induced by $\pi^\#$ on sections, then $\pi$ is $\Gm$-equivariant and the induced morphism $[\spec_{\cO_X}(I^\bullet)/\Gm]\to X$ is a good moduli space as the $\Gm$-invariant functions are those on degree zero, and
    \item if we denote by $i$ the morphism $X\to [\spec_{\cO_X}(I^\bullet)/\Gm]$ induced by $i^\#$, this is an open embedding \cite{QR}*{Prop. 4.3.4},
    \item there is a morphism $ [\spec_{\cO_X}(I^\bullet)/\Gm]\to [X\times \bA^1/\Gm]$, given by the inclusion $\bigoplus_{n\le 0}\cO_X[T^{-1}]\to I^\bullet$, which $i(X)$ with the preimage of $X\times \{1\}$ and its complement with the complement of $T^{-1}=0$, so in particular $[\spec_{\cO_X}(I^\bullet)/\Gm]\smallsetminus i(X)$ is Cartier.
\end{enumerate}
So to summarize, given $X$ with a closed subscheme $Z\subseteq X$, we can construct an algebraic stack $\widetilde{X}$ which has $X$ as a good moduli space, the good moduli space map $\widetilde{X}\to X$ has a section $i:X\to \widetilde{X}$ which is an open embedding, the complement of $i(X)$ is a Cartier divisor, and $\cX$ contains as an open substack the blow-up of $X$ at $Z$.

One can perform the construction above equivariantly. More explicitly, if a group $G$ acts on a scheme $X$, and $Z\subseteq X$ is a closed subscheme which is $G$-invariant with ideal sheaf $I$, then $G$ acts on $\spec_{\cO_X}(I^\bullet)$ leaving invariant the locus where $T^{-1}=0$ (namely, $V(I^\bullet_+)$) invariant. In this case, one can obtain again an algebraic stack $[\spec_{\cO_X}(I^\bullet)/G\times \Gm]$ which has $[X/G]$ as an open substack (the complement of a Cartier divisor) and which contains the blow-up of $[X/G]$ at $[Z/G]$
\begin{example}
    The deformation to the normal cone of the origin $V(x,y)\subseteq \bA^2$ is by definition $\spec(I^\bullet)$, which can be identified with $\spec(k[x,y,T,X,Y]/(XT-x, YT-y)$ where:
    \begin{itemize}
        \item $X,Y$ are the generators of the maximal ideal in degree 1,
        \item $T$ has degree $-1$,
        \item $x,y$ have degree 0,
        \item the map $X\to \spec(I^\bullet)$ is given by sending $T\mapsto 1$, and 
        \item the map $\spec(I^\bullet)\to X$ is given by the inclusion of the degree 0 component.
    \end{itemize}
    Therefore we can identify it with $\spec(k[X,Y,T])$ with the action of $\Gm$ with weights $(1,1,-1)$, the map $\bA^3\to \bA^2$ given by $(a,b,u)\mapsto (au,bu)$
    induces the good moduli space morphism $[\bA^3/\Gm]\to \bA^2$ and the section is the locus where $T=1$ in $[\bA^3/\Gm]$. We can identify $[\{(a,b,t):(a,b)\neq (0,0)\}/\Gm]\hookrightarrow[\bA^3/\Gm]$ with the blow-up of the origin in $\bA^2$ as in \cite{Inc}*{\S 2.2}.
    
    If a group $G$ acts on $\bA^2$, we can extend its action to $\bA^3$,
    by acting trivially on the $T$ component. So as before one has two morphism $[\bA^2/G]\to[\bA^3/G\times \Gm] $ and $[\bA^3/G\times \Gm]\to [\bA^2/G]$ whose composition is the identity, where the first is the complement of a Cartier divisor, and $[\bA^3/G\times \Gm]$ contains as an open the blow-up of $\cB G$ in $[\bA^2/G]$. 
\end{example}
\begin{Lemma}\label{lemma_gen_blowup_comes_with_a_cartier_divisor}
    Let $I$ be an ideal sheaf on $[X/G]$, and consider $\pi:\cB_I[X/G]\to [X/G]$ the extended blow-up of $[X/G]$ along $I$. Then there is a line bundle $\cL$ on $\cB_I[X/G]$ with a section $s\in \oH^0(\cB_I[X/G],\cL)$ whose zero locus is contained in $\pi^-1(V(I))$, which is the scheme-theoretic closure of the ideal sheaf of the exceptional divisor of $B_I[X/G]\subseteq \cB_I[X/G]$. 
\end{Lemma}
\begin{Def}
    We define the line bundle of \Cref{lemma_gen_blowup_comes_with_a_cartier_divisor} the $\cO(-1)$ \textit{of} $\cB_I[X/G]$, and the zero-locus of the section of \Cref{lemma_gen_blowup_comes_with_a_cartier_divisor} the \textit{exceptional locus}.
\end{Def}
\begin{proof}
    \textcolor{red}{to be added}
\end{proof}
\begin{proof}[Proof of \Cref{thm_you_can_embed_a_Stack_in_one_with_open_dm_by_taking_deformations_to_nc}] From \cite{ER21}, there is a sequence of blow-ups $\cX_n\to \cX_{n-1}\to ... \to \cX_1 = \cX$ such that $\cX_n$ contains an open substack $\cX_n^{\dm}\subseteq \cX_n$ which is Deligne-Mumford,
and it admits a coarse moduli space $\cX^\dm_n\to X^\dm_n$ which is projective over $X$. We can realize each of these blow-ups as open in an appropriate deformation to the normal cone as in \Cref{subsubsection_def_to_normal_cone}. We are left with proving that we can realize $\cX^\dm$ as the semistable locus of a line bundle \textcolor{red}{to be finished}
\end{proof}
\end{comment}
\section{An application of \Cref{thm_intro_simplified} and  \Cref{thm_intro_you_can_embed_a_Stack_in_one_with_open_dm_by_taking_deformations_to_nc}}\label{section_examples_extended}
\Cref{thm_intro_simplified} now follows from the results in \Cref{sec:qmaps} and \Cref{section_extended_blowups}: from \Cref{section_extended_blowups} given an algebraic stack $\cX$ with a dense properly stable locus, one compactify the space of maps to $\cX$ as follows. Either the assumptions of \Cref{teo_intro_cX_contains_open_proper_dm} apply, or one can enlarge $\cX\subseteq \widetilde{\cX}$ using \Cref{thm_intro_you_can_embed_a_Stack_in_one_with_open_dm_by_taking_deformations_to_nc}, so that the assumptions of \Cref{teo_intro_cX_contains_open_proper_dm} apply to $\widetilde{\cX}$. One can use the inclusion $\cX\subseteq \widetilde{\cX}$ and \Cref{teo_intro_cX_contains_open_proper_dm} to compactify the space of maps to $\cX$, and the projection morphism $\widetilde{\cX}\to \cX$ to give a modular interpretation to maps to $\widetilde{\cX}$ in terms of maps to $\cX$.

In this section we give an example of how to use the \Cref{thm_intro_you_can_embed_a_Stack_in_one_with_open_dm_by_taking_deformations_to_nc} and \Cref{teo_intro_cX_contains_open_proper_dm} to compactify the space of maps to the GIT moduli space of plane cubics, which we denote by $\sP^1$. Recall that this GIT moduli space is defined as $\sP^1:= [\bP(\oH^0(\cO_{\bP^2}(3)))^{ss}/\PGL_3]$, where $\bP(\oH^0(\cO_{\bP^2}(3)))^{ss}$ is the open in $\bP(\oH^0(\cO_{\bP^2}(3)))$ parametrizing plane cubics with at most nodal singularities. This algebraic stack both does not have an open substack which is proper and Deligne-Mumford, and it does not have (as far as we know) a modular enlargement containing a proper and Deligne-Mumford open substack.
We show that also in this case, from the properties of extended blow-ups, one can still carry a detailed analysis of the objects on the boundary of $\cQ_g(\widetilde{\sP}^1,\widetilde{\sP}^1_\dm,\beta)$ for an extended blow-up $\widetilde{\sP}^1\to \sP^1$, by post-composing a map $\cC\to\widetilde{\sP}^1$ with $\cC\to\widetilde{\sP}^1\to \sP^1$. In \Cref{prop_conditions_induced_by_phi_tilda_on_phi_and_f} we list some of the properties of maps $C\to \widetilde{\sP}^1$ appearing as degenerations in the quasimaps moduli space of maps $\bP^1\to \bP(\oH^0(\cO_{\bP^2}(3)))^{ss}/\PGL_3]$ coming from pencils of cubics.


\subsection{GIT moduli space of plane cubics}We begin by reporting a few classical results, which one might find in \cites{MFK, ascher2023moduli}
\begin{enumerate}
    \item each smooth cubic $C$ is GIT-stable,
    \item each nodal cubic can be degenerated, along $[\bA^1/\Gm]$, to $x_0x_1x_2$
    \item each smooth plane cubic can be written as $x_0^3+ x_1^3 + x_2^3 + \lambda x_0x_1x_2$, \cites{nakamura2004planar, hesse1897ludwig}, so each plane cubic contains $\bmu_3\rtimes S_3$ as automorphism group,
    \item the stabilizer of $x_0x_1x_2$ is $\Gm^3/{\Gm}_{, \Delta}\rtimes S_3$ where ${\Gm}_{, \Delta}$ is the diagonal subgroup of $\Gm^3$ and the $S_3$ action permutes the coordinates $x_0, x_1, x_2$. This is straightforward once we realize that any automorphism of $\PGL_3$ which fixes $x_0x_1x_2=0$ fixes the three nodes $[1,0,0]$, $[0,1,0]$ and $[0,0,1]$,
    \item the GIT quotient $\bP(\oH^0(\cO_{\bP^2}(3)))/\!\!/\PGL_3$ is isomorphic to $\bP^1$.
\end{enumerate}
In particular, we introduce the following
\begin{Notation}
    We denote by $G:=\Gm^3/{\Gm}_{, \Delta}\rtimes S_3$, the stabilizer of $x_0x_1x_2$, and by $\cG:=\Gm^3\rtimes S_3$ the group which surjects to $G$ in the obvious way. We further denote $\sP^1:= [\bP(\oH^0(\cO_{\bP^2}(3)))^{ss}/\PGL_3]$, and if $\cB G\hookrightarrow \sP^1$ is the inclusion of the point corresponding to $x_0x_1x_2$, we assume that it corresponds to the point at infinity in $\bP^1$.
\end{Notation}

\begin{Lemma}\label{lemma_luna_slice_curlyP1}
    There is a Luna slice for $\sP^1$ at $x_0x_1x_2$ isomorphic to $[\bA^3/G]$ where the action of $G$ on $\bA^3$ comes from an action of $\cG$ defined as follows:
    \begin{align*}
    ((\lambda,\mu,\nu),\sigma)\in \Gm^3\rtimes S_3\text{ acts as }((\lambda,\mu,\nu),\sigma)*(x_1,x_2,x_3) = \left(\frac{\lambda^2}{\mu\nu}x_{\sigma^{-1}(1)}, \frac{\mu^2}{\lambda\nu}x_{\sigma^{-1}(2)}, \frac{\nu^2}{\lambda\mu}x_{\sigma^{-1}(3)}\right).
    \end{align*}
\end{Lemma}
\begin{proof} We denote by $x$ the point $x_0x_1x_2$.
It suffices to study the action of $\cG$ on $T_x$, the tangent space of $x\in \bP(\oH^0(\cO_{\bP^2}(3)))$, and the image of the map on tangent spaces $T_{\Id}\PGL_3\to T_x$ given by the inclusion of the orbit of $x$, and 
where $\Id$ is the identity in $\PGL_3$. Recall that we can identify the tangent space of the identity of $\PGL_3$ as the matrices as follows
\[
\begin{bmatrix}1+\epsilon_1&\epsilon_4&\epsilon_7\\\epsilon_2&1+\epsilon_5&\epsilon_8\\\epsilon_3&\epsilon_6&1+\epsilon_9
\end{bmatrix}\text{ such that }\epsilon_1+\epsilon_5+\epsilon_9=0\text{ and }\epsilon_i\epsilon_j=0\text{ for every }i,j.\]
Then the image of $T_{\Id}\PGL_3\to T_xX$ can be identified with the following forms \[\{((1+\epsilon_1)x_0+\epsilon_2x_1+\epsilon_3x_2)(\epsilon_4x_0+(1+\epsilon_5)x_1+\epsilon_6x_2)(\epsilon_7x_0+\epsilon_8x_1+(1+\epsilon_9)x_2):\epsilon_i\epsilon_j=0, \epsilon_1+ \epsilon_5+\epsilon_9=0\}.\]
The above expression simplifies to
\[x_0x_1x_2 + \epsilon_7x_0^2x_1 +\epsilon_4x_0^2x_2+\epsilon_8x_0x_1^2 +\epsilon_2x_1^2x_2 + \epsilon_6x_0x_2^2 + \epsilon_3x_1x_2^2.\]
In particular, in the affine chart of $\bP(\oH^0(\cO_{\bP^2}(3)))$ given by $x_0x_1x_2\neq 0$, we can take as coordinates for $T_xX$ the following ones:
\[\frac{x_0^3}{x_0x_1x_2}, \frac{x_1^3}{x_0x_1x_2}, \frac{x_2^3}{x_0x_1x_2}, \frac{x_0^2x_1}{x_0x_1x_2},
\frac{x_0^2x_2}{x_0x_1x_2},
\frac{x_1^2x_0}{x_0x_1x_2},
\frac{x_1^2x_2}{x_0x_1x_2},
\frac{x_2^2x_0}{x_0x_1x_2},
\frac{x_2^2x_1}{x_0x_1x_2}.
\] Then the image of $T_{\Id}\PGL_3\to T_xX$ can be identified with the linear subspace generated by \[
<\frac{x_0^2x_1}{x_0x_1x_2},
\frac{x_0^2x_2}{x_0x_1x_2},
\frac{x_1^2x_0}{x_0x_1x_2},
\frac{x_1^2x_2}{x_0x_1x_2},
\frac{x_2^2x_0}{x_0x_1x_2},
\frac{x_2^2x_1}{x_0x_1x_2}>.
\] Since the linear subspace $<\frac{x_0^3}{x_0x_1x_2}, \frac{x_1^3}{x_0x_1x_2}, \frac{x_2^3}{x_0x_1x_2}>$ is $G$-invariant, we can take it as our $W$. It is straightforward to check now that the action of $G$ on $W$ is the desired one.  
\end{proof}
The advantage of our local description is that we can explicitly determine how many extended blow-ups are needed to apply \Cref{thm_you_can_embed_a_Stack_in_one_with_open_dm_by_taking_deformations_to_nc}. In particular:
\begin{Lemma}
Consider $\bA^3$ with the action of $G$ as in \Cref{lemma_luna_slice_curlyP1}.
The extended blow-up of the origin $\EB_0\bA^3$ contains a Deligne-Mumford stack $(\EB_0\bA^3)_\dm$ which is proper over the good moduli space of $\EB_0\bA^3\to \bA^1$, and the stabilizer of the unique closed point of $(\EB_0\bA^3)_\dm$ over $0\in \bA^1$ is $\bmu_3^{\otimes 3}/\bmu_{3,\Delta}\rtimes S_3$.
\end{Lemma}
\begin{proof}Recall that we can identify $\EB_0\bA^3$ with $[\bA^4/G\times \Gm]$ where the action of $G$ on the first three components is the original one and trivial on the last component. Similarly, the action of $\Gm$ on the first three components is with weight $1$ and on the last component is with weight $-1$. It is now straightforward to check that the point $(1,1,1,0)$ has $\bmu_3^{\otimes 3}/\bmu_{3,\Delta}\rtimes S_3$ as stabilizers. 
\end{proof}
Since the formation of extended blow-ups commutes with \'etale base change, we have
\begin{Cor}
    The extended blow-up of $\cB G \subseteq \sP^1$, which we denote by $\widetilde{\sP}^1$, contains a
    proper Deligne-Mumford stack, which we denote by $\widetilde{\sP}^1_\dm$. The coarse moduli space of $\widetilde{\sP}^1_\dm$ is $\bP^1$. 
\end{Cor}
\subsection{Space of maps to $\sP^1$}
We now focus on the problem of compactifying maps to $\sP^1$, in a specific example. Consider a moduli space $\cU$ of pencils of plane cubics, i.e. of maps $\bP^1\to \bP(\oH^0(\cO_{\bP^2}(3)))$, up to $\PGL_3$ (namely, up to change of coordinates) and such that each map $\bP^1\to \bP(\oH^0(\cO_{\bP^2}(3)))$ intersects the discriminant locus transversally (this is a $\PGL_3$-invariant condition). In particular:
\begin{enumerate}
    \item there is a fibration $\cC\to \bP^1$ where each fiber is a plane cubic, the generic fiber is smooth and the singular fibers have a single node,
    \item there are exactly 12 nodal singular fibers, as the discriminant is a hypersurface of degree 12 in $\bP(\oH^0(\cO_{\bP^2}(3)))$,
    \item if we denote by $\phi:\bP^1\to \sP^1$ the induced morphism, then $\phi^{-1}\cB G=\emptyset$ as the fibers of $\cC\to\bP^1$ have a single node,
    \item if we denote by $f:\bP^1\to \sP^1\to\bP^1$ the composition of $\phi$ with the good moduli space map, then $f$ has degree 12 and it is unramified over $\infty$ (the point corresponding to $x_0x_1x_2$).
\end{enumerate}
Then there is a map $\cU\to \cQ_{0}(\widetilde{\sP}^1,\widetilde{\sP}^1_\dm)$. Assume it is a monomorphism. 
\begin{Notation}\label{notation_sections}
    We denote by $\overline{\cU}$ the closure of $\cU$ inside $\cQ_{0}(\widetilde{\sP}^1,\widetilde{\sP}^1_\dm)$. Given a map $\psi:C\to\sP^1\to \widetilde{\sP}^1$ corresponding to a point in $\cU$, there are 12 points in $C$ such that $\psi(p)\not \in \widetilde{\sP}^1_\dm$. We denote by \[\widetilde{\phi}:(C\xrightarrow{\pi} B;x_1,\ldots,x_{12})\to \sP^1\]
    the objects $\widetilde{\phi}:C\to\sP^1$ of $\overline{\cU}(B)$, such that there are 12 \textit{sections} $x_1,\ldots,x_{12}$ of $\pi$ that satisfy \[\widetilde{\phi}^{-1}(\widetilde{\sP}^1\smallsetminus\widetilde{\sP}^1_\dm) = \{x_1,\ldots,x_{12}\}.\] 
\end{Notation}
The following remark follows from the definition of stable quasimap (\Cref{def_stable_qmap}), and since the coarse moduli space of $\widetilde{\sP}^1_\dm$ is $\bP^1$.
\begin{Remark}
    Let $\widetilde{\phi}:\cC\to \widetilde{\sP}^1$ a morphism in $\overline{\cU}$ inducing $f:C\to \bP^1$ on good moduli spaces. Then if $\{x_1,\ldots,x_m\}=\widetilde{\phi}^{-1}(\widetilde{\sP}^1\smallsetminus\widetilde{\sP}^1_\dm)$, for every $0<\epsilon $ the line bundle $\omega_C(\epsilon\sum x_i)\otimes f^*\cO_{\bP^1}(3)$ is ample on $C$.
\end{Remark}

\begin{Lemma}\label{lemma_irred_cp_which_map_to_infty_map_to_exc_locus}
     Let $R$ be a DVR with generic point $\eta$ and closed point $p$ and let $(C\to \spec(R);x_1,...,x_{12})\xrightarrow{\widetilde{\phi}}\widetilde{\sP}^1$ correspond to a morphism $\spec(R)\to \overline{\cU}$ witn $12$ sections as in \Cref{notation_sections}. Assume that $\eta\to\overline{\cU}$ maps into $\cU$, and let $D\subseteq C_p$ be an irreducible component of the central fiber of $C\to \spec(R)$. Then:
     \begin{enumerate}
         \item if $f:D\to C\to \bP^1$ is finite and $x_i\in D$ is a marked point, $\widetilde{\phi}(x_i)$ is not contained in the exceptional divisor of $\widetilde{\sP}^1$,
         \item if $D\to C\to \bP^1$ maps to a point $z$, then $D$ is contained in the exceptional locus of $\widetilde{\sP}^1$ if and only if $z=\infty$, and
         \item if $f:D\to C\to \bP^1$ is finite $x\in D$ is a smooth point of $D$, then the multiplicity of $f^{-1}(\infty)$ at $x$ is $\#\{i:x_i=x\}$.         
     \end{enumerate}
\end{Lemma}
\begin{proof}
    Let $\cO(-1)$ be the line bundle on $\widetilde{\sP}^1$ given by the exceptional divisor, let $\cL=\widetilde{\phi}^*\cO(-1)$ and let $t$ be the section of $\cL$ given by the pull-back of the exceptional locus.
    By construction of the morphism $\eta\to \cU$, the section $t$ does not vanish along the generic fiber of $C$. Then $V(t)$ is, set theoretically, a union of irreducible components of the special fiber $C_p\subseteq C$ that don't map finitely to $\bP^1$, so point (1) follows since the markings $x_i$ are smooth on the special fiber.
    
    For (2), since the exceptional locus is contained in $\pi^{-1}\cB G$ from \Cref{def_gen_blowup_comes_with_a_cartier_divisor}, if $D\subseteq C_p$ maps to a point different from $\infty$ via $f$, then $D$ is not contained in $V(t)$. We are left with showing that if $f(D)=\infty$ then $t$ vanishes on $D$. By assumption, the generic point of $\widetilde{\phi}(D)$ maps to the Deligne-Mumford locus in $\widetilde{\sP}^1$, and there is a unique point of $\widetilde{\sP}^1_\dm$ over $\infty$. Such a point is in the exceptional locus, so $V(t)$ vanishes over it.

    Finally for (3) observe that the Cartier divisor $f^{-1}(\infty)$ is of the form $\cO_C(a_1x_1+\ldots +a_{12}x_{12})\otimes \cO_C(\Delta)$, where $\Delta$ is effective and supported on the irreducible components of the special fiber mapping to $\infty$ and $a_i\ge 1$. So it suffices to prove that $a_i=1$ for every $i$, which can be checked on the generic fiber. This follows by assumption, as the maps in $\cU$ intersect the discriminant locus transversally.
\end{proof}
We will need the next auxiliary lemma
\begin{Lemma}\label{lemma_when_a_map_from_a_dvr_to_An_lifts_to_the_blow_up_without_the_proper_transform_of_the_axis}
    Let $R$ be a DVR with generic point $\eta$ and closed point $p$, let $\Gm^r$ be a torus acting on $\bA^n$ diagonally, let $\pi:B_0[\bA^n/\Gm^r]\to [\bA^n/\Gm^r]$ be the blow-up of $[0/\Gm^r]\hookrightarrow [\bA^n/\Gm^r]$.  Let $\phi:\spec(R)\to [\bA^n/\Gm^r]$ be a morphism sending the generic point $\eta\in\spec(R)$ to a point in $[\bA^n/\Gm^r]$ which does not lie in $[\{x_i=0\}/\Gm^r]\subseteq [\bA^n/\Gm^r]$ for any $i$, where $x_i$ are the coordinates on $\bA^n$, and assume that $\phi(p)$ maps instead to $[0/\Gm^r]$. Then $\phi$ lifts uniquely to a morphism $\widetilde{\phi}:\spec(R)\to B_0[\bA^n/\Gm^r]$. Moreover, if we denote by $x_1,...,x_n$ the coordinates in $\bA^n$, we have that $\widetilde{\phi}(p)$ lies along the closed substack of $B_0[\bA^n/\Gm^r]$ given by the proper transform of the axis in $[\bA^n/\Gm^r]$ if and only if the there are $i,j$ such that $\phi^*x_i$ and $\phi^*x_j$ vanish at $p$ with different orders.
\end{Lemma}
\begin{proof}
    The existence of $\widetilde{\phi}$ follows from the valuative criterion for properness since $B_0[\bA^n/\Gm^r]\to [\bA^n/\Gm^r]$ is representable and proper, and an isomorphism away from $[0/\Gm^r]$. For the moreover part, one can identify the extended blow-up $\EB_0[\bA^n/\Gm^r]\to [\bA^n/\Gm^r]$ with $[\bA^{n+1}/\Gm^{r+1}]$ as in \Cref{example_ext_blowup_of_origin_in_A2} for the case $n=2$, and if $\bA^{n+1}$ has coordinates $X_1,...,X_n,T$, the map $[\bA^{n+1}/\Gm^{r+1}]\to [\bA^{n}/\Gm^{r}]$ sends $x_i\mapsto X_iT$. A morphism $\spec(R)\to [\bA^n/\Gm^r]$ corresponds to $r$ line bundles $\cL_1,\ldots,\cL_r$, together with a section $\phi^*x_i\in \oH^0(\spec(R),\cG_i)$ for $1\le i\le n$, where each $\cG_i$ is a linear combination of $\cL_1,\ldots,\cL_r$.
    If we denote by $m_i$ is the order of vanishing of $\phi^*x_i$ for every $i$ and $m=\min\{m_i\}$, the sections $\phi^*x_i$ lift uniquely to sections of $\widetilde{x}_i\in \oH^0(\spec(R),\cG_i\otimes \cO_{\spec(R)}(-mp))$. One can check that the lift $\widetilde{\phi}$ comes from the line bundles \[\{\cG_i\otimes \cO_{\spec(R)}(-mp)\}_i\cup\{\cO_{\spec(R)}(mp)\}\]
    and the sections being the sections $\widetilde{x}_i$ and the one coming from dualizing the inclusion of the ideal sheaf $\cO_{\spec(R)}(-mp)\to \cO_{\spec(R)}$. Now the desired statement follows observing that $m<m_j$ is equivalent to the vanishing of $\widetilde{x}_j$ at $p$.
\end{proof}
\begin{Prop}\label{prop_conditions_induced_by_phi_tilda_on_phi_and_f}
    Let $(C\to B;x_1,...,x_{12})\xrightarrow{\widetilde{\phi}}\widetilde{\sP}^1$ be a $B$-point in $\overline{\cU}$ as in \Cref{notation_sections}, let $\phi:C\to \sP^1$ be the composition of $\widetilde{\phi}$ and the extended weighted blowup $\widetilde{\sP}^1\to \sP^1$, and let $f:C\to \bP^1$ the composition to the good moduli space map $\sP^1\to \bP^1$. Then:
    \begin{enumerate}
        \item there is a line bundle $\cL$ on $C$ with a section $t\in \oH^0(C,\cL)$ such that, for every $b\in B$ we have an equality in sets $V(t)_b = \{D\subseteq C_b$ irreducible component such that $f_b(D)=\infty\}$,
        \item if $D\subseteq C_b$ is an irreducible component such that $f_b(D)=\infty$, then the map $D\to \sP^1$ factors via $D\to \cB G\to \sP^1$,
        \item $f^*\cO_{\bP^1}(1)\cong \cO_\cC(x_0+ \ldots+ x_{12})\otimes \cL^{\otimes 3}$ and $f^{-1}(\infty)=V(x_0\cdot \ldots\cdot x_{12}\cdot t^3)$, where the latter is a section of $\cO_C(x_0+ \ldots+ x_{12})\otimes \cL^{\otimes 3}$,
        \item let $D\subseteq C_b$ be an irreducible component such that $f(D)$ is not a point and let $n\in D$ be a node which maps to $\infty$. Let $V\to D$ be an \'etale cover of $n$ which is a scheme, let $v\in V$ be a point mapping to $n$, and let $\cC\to V$ be the resulting family of plane cubics. Then the singularities of $\cC$ around each node of $\cC_v$ are isomorphic.
    \end{enumerate}
\end{Prop}
\begin{proof}
    If we denote by $\cL=\widetilde{\phi}^*\cO(1)$ and $t$ the pull-back of a section that vanishes along exceptional locus, then point (1) follows from \Cref{lemma_irred_cp_which_map_to_infty_map_to_exc_locus}. Point (2) instead follows from \Cref{def_gen_blowup_comes_with_a_cartier_divisor}. Point (3) is true away from the vanishing locus of $t$ from \Cref{lemma_irred_cp_which_map_to_infty_map_to_exc_locus}, so it suffices to prove it in an \'etale neighbourhood of $t=0$. Similarly, point (4) can be checked on an \'etale neighbourhood of $t=0$, so we choose a neighbourhood where we will check both, as follows. We already computed a Luna slice around the polystable point of $\sP^1$ in \Cref{lemma_luna_slice_curlyP1},  so from \Cref{subsection_luna_slice} there are diagrams as follows  \[
    \xymatrix{[\widetilde{W}/G\times \Gm]\ar[d] \ar[r]^\gamma &[\bA^4/G\times \Gm]\ar[d]\\
    [W/G]\ar[d]\ar[r]^\beta  &[\bA^3/G]\ar[d]\\
    W/\!\!/G\ar[r]^\alpha& \bA^1}\text{ }\text{ }\text{ }\text{ }\text{ }\text{ }\xymatrix{[\widetilde{W}/G\times \Gm]\ar[d] \ar[r] & \widetilde{\sP}^1\ar[d]\\
    [W/G]\ar[d]^\pi\ar[r]^\iota  &\sP^1\ar[d]\\
    W/\!\!/G\ar[r]^i& \bP^1}
    \]
    with $i$ and $\alpha$ \'etale, with all squares cartesian, and with a point $w\in W/\!\!/G$ mapping to $0\in \bA^1$ and $\infty \in \bP^1$. Recall that $[\bA^4/G\times \Gm]\to [\bA^3/G]$ is the extended blow-up of the origin in $[\bA^3/\Gm]$.
    We can then pull-back the morphism $\phi:C\to \sP^1$ along $\iota$, this will give us an \'etale neighbourhood where we will check conditions (3) and (4). We denote by $U:=C\times_{\sP^1}[W/G]$ and by \[
        \phi_U:U\to [W/G]\to [\bA^3/G]\text{ and }\widetilde{\phi}_U:U\to [\widetilde{W}/G\times\Gm]\to [\bA^4/G\times \Gm]\] the composition of the two second projections with the maps $\beta$ and $\gamma$.

    To check condition (3), observe that if $x_1,x_2,x_3$ are the coordinates of $\bA^3$, then the ring of $G$-invariants consists of $\spec(k[x_1x_2x_3])$. Moreover if $X_1,X_2,X_3,T$ are the coordinates on $\bA^4$, the blowdown morphism $[\bA^4/G\times\Gm]\to [\bA^3/G]$ is induced by $x_i\mapsto X_iT$. Then the composition sends $x_1x_2x_3\to X_1X_2X_3T^3$, so if we denote by $\widetilde{\pi}:[\bA^4/G\times \Gm]\to \bA^1$ is the composition of the extended blow-down $[\bA^4/G\times \Gm]\to [\bA^3/G]$ and the good moduli space map $[\bA^3/G]\to \bA^1$, then \begin{equation}\label{equation}\widetilde{\pi}^{-1}(0)=[V(X_1X_2X_3T^3)/G\times \Gm]. \end{equation} In particular, as $V(T)$ is the exceptional divisor, the Cartier divisor $V(t)\subseteq U$ agrees with $V(\widetilde{\phi}_U^*T)$. Similarly, from \Cref{lemma_irred_cp_which_map_to_infty_map_to_exc_locus}, the Cartier divisor $V(\widetilde{\phi}_U^*(X_1X_2X_3))$ agrees with $V(p_1\cdot\ldots\cdot p_{12})$, where we denoted by $p_i$ the pull-backs of the markings on $C$ to $U$. Now point (3) follows from \Cref{equation}.

    To prove (4), we can take a further \'etale cover of $[\bA^3/G]$ as follows. Recall that $G=(\Gm^3/\Gmdelta)\rtimes S_3$, so one can take the $S_3$-torsor $[\bA^3/(\Gm^3/\Gmdelta)]\to [\bA^3/G]$. This gives a family of maps $\sC\to [\bA^3/(\Gm^3/\Gmdelta)]$, and over the origin one can check from the explicit description of the Luna slice in \Cref{lemma_luna_slice_curlyP1} that the directions $\{x_i=0\}\subseteq \bA^3$ correspond to degenerations of $V(x_0x_1x_2)\subseteq \bP^2$ along which a node does not smooth. Extended blow-ups commutes with flat base changes, so $[\bA^4/G\times \Gm]\times_{[\bA^3/\Gm]}[\bA^3/(\Gm^3/\Gmdelta)]\cong[\bA^4/(\Gm^3/\Gmdelta)\times \Gm]$. 
    Since $\widetilde{\phi}(n)$ maps to $\widetilde{\sP}^1_\dm$, $\widetilde{\phi}_U(n)$ does not map in $[V(X_1X_2X_3)/G\times \Gm]\subseteq [\bA^4/G\times \Gm]$. In other terms, $\widetilde{\phi}_U(n)$ maps to the unique point of $[\bA^4/G\times \Gm]$ which is in $\widetilde{\sP}^1_\dm\times_{\widetilde{\sP}^1}[\bA^4/G\times \Gm]$ and is over $\infty$, namely the point $(1,1,1,0)\subseteq [\bA^4/G\times \Gm]$. Then from \Cref{lemma_when_a_map_from_a_dvr_to_An_lifts_to_the_blow_up_without_the_proper_transform_of_the_axis} the order of vanishing of $\widetilde{\phi}^*_Ux_i$ does not depend on $i$. The conclusion now follows from \Cref{lemma_when_a_map_from_a_dvr_to_An_lifts_to_the_blow_up_without_the_proper_transform_of_the_axis}.
    \end{proof}
\appendix
\section{Extended blow-ups, rational pointed curves and Hassett's conjecture}\label{appendix}
In this appendix, we give a criterion for when a stack can be obtained as an extended weighted blow-up, and we use it in order to prove \Cref{thm_intro_hassett}.
\subsection{A criterion for being an extended blow-up}\label{subsection_criterion_for_ext_wblowups}
In this subsection we give a few criteria under which a morphism of algebraic stacks $\widetilde{\cX}\to \cX$ is an extended weighted blow-up; this will come in handy later.
\begin{Prop}\label{prop:criterion}
    Let $\pi\colon\widetilde{\cX}\to\cX$ be a morphism of algebraic stacks such that:
    \begin{enumerate}
        \item the stacks $\widetilde{\cX}$ and $\cX$ have the same good moduli space, i.e. there is a good moduli space $\rho:\cX\to X$ and the composition $\rho\circ\pi$ also defines a good moduli space for $\widetilde{\cX}$;
        \item the canonical map $\cO_\cX\to \pi_*\cO_{\widetilde{\cX}}$ is an isomorphism;
        \item there is a morphism $\widetilde{\cX}\to [\bA^1/\Gm]$ such that $\widetilde{\cX}\to \cX\times[\bA^1/\Gm]$ is representable;
        \item if $\widetilde{\cE}\subseteq\widetilde{\cX}$ is the Cartier divisor induced by $\widetilde{\cX}\to\cX\times [\bA^1/\Gm]\to [\bA^1/\Gm]$, we have $(\pi|_{\widetilde{\cE}})_*(\cN_{\widetilde{\cE}}^{\otimes d})=0$ for $d>0$, where $\cN_{\widetilde{\cE}}$ is the normal bundle of $\widetilde{\cE}$.
    \end{enumerate}
    Then $\pi\colon\widetilde{\cX}\to \cX$ is an extended weighted blow-up of $\cX$ along the filtration
    \[ \cO_\cX\supset\pi_*\cO(-\widetilde{\cE}) \supset \pi_*\cO(-2\widetilde{\cE}) \supset \ldots \supset \pi_*\cO(-n\widetilde{\cE}) \supset \ldots\]
\end{Prop}
\begin{proof}
    First, observe that (1) implies that $\widetilde{\cX}\to \cX$ is cohomologically affine \cite{Alp}*{Proposition 3.13}. 
    Let $\cY$ be the total space of the $\Gm$-torsor $g\colon\cY\to\widetilde{\cX}$ associated to $\widetilde{\cE}$. As $\cY\to\widetilde{\cX}$ is affine and $\widetilde{\cX}\to\cX$ is cohomologically affine, also $\cY\to\cX$ is cohomologically affine and representable from (3), hence affine \cite{Alp}*{Proposition 3.3}.

    Write $\cY=\spec_{\cO_\cX}(\cA)$. Observe that, as $g\colon\cY\to\widetilde{\cX}$ is the $\Gm$-torsor associated to $\widetilde{\cE}$, we have that $g_*\cO_{\cY}=\oplus_{n\in\bZ} \cO(-n\widetilde{\cE})$. We deduce that $\cA=\oplus_{n\in\bZ} \pi_*\cO(-n\widetilde{\cE})$.    We prove that $\pi_*\cO(n\widetilde{\cE})=\cO_{\cX}$ for $n\geq 0$.
    
    Observe that by (2) the map $\cO_\cX \to \pi_*\cO_{\widetilde{\cX}}$ is an isomorphism, and consider the exact sequence
    \[ 0 \to \cO_{\widetilde{\cX}}(n\widetilde{\cE}) \to \cO_{\widetilde{\cX}}((n+1)\widetilde{\cE}) \to j_* \cN_{\widetilde{\cE}}^{\otimes (n+1)} \to 0. \]
    Pushing it forward along $\pi$, we obtain an isomorphism $\pi_*\cO_{\widetilde{\cX}}(n\widetilde{\cE})\simeq \pi_*\cO_{\widetilde{\cX}}((n+1)\widetilde{\cE})$, because the third term of the sequence vanishes by (4). We can therefore conclude by induction that $\pi_*\cO(n\widetilde{\cE})\simeq\cO_{\cX}$ for $n\geq 0$.

    Set then $\cI_n:=\pi_*\cO_{\widetilde{\cX}}(-n\widetilde{\cE})$ for $n\in\bZ$. Then by pushing down the filtration
    \[ \cO_{\widetilde{\cX}}\supset \cO_{\widetilde{\cX}}(-\widetilde{\cE}) \supset \ldots \supset \cO_{\widetilde{\cX}}(-n\widetilde{\cE}) \supset \ldots \]
    we get a filtration 
    \[ \cO_{\cX}\supset \cI_1 \supset \cI_2 \supset \ldots \supset \cI_n \supset \ldots. \]
    One can check that it satisfies the conditions of \Cref{def_weighted_embedding},
    hence $\cA=\oplus_{n\in\bZ} \cI_n$ and $\widetilde{\cX}=\EB_{\cI_\bullet}\cX$.
\end{proof}
\begin{Lemma}\label{lm:iso of O}
    Let $\pi:\widetilde{\cX}\to \cX$ be a morphism of algebraic stacks,
    assume that there is a schematically dense open immersion $i:\cX\hookrightarrow \widetilde{\cX}$ satisfying $\pi\circ i = \id_{\cX}$. Then $\pi_*\cO_{\widetilde{\cX}}=\cO_\cX$.
\end{Lemma} 
For example, if $i$ is topologically dense and $\widetilde{\cX}$ is reduced then $i$ is schematically dense.
\begin{proof}
    The morphism $\cO_{\widetilde{\cX}}\to i_*\cO_{\cX}$ is injective, since $i$ is schematically dense. So also $\pi_*\cO_{\widetilde{\cX}}\to \pi_*i_*\cO_{\cX}$ is injective. Since the composition $\cO_{\cX} \longrightarrow \pi_*\cO_{\widetilde{\cX}} \longrightarrow \pi_*i_*\cO_{\widetilde{\cX}}=\cO_{\cX}$
    is surjective, the morphism $\pi_*\cO_{\widetilde{\cX}}\to \pi_*i_*\cO_{\cX}$ is also surjective. This implies that $\cO_\cX \to \pi_*\cO_{\widetilde{\cX}}$ is an isomorphism as claimed.
\end{proof}
\begin{Cor}\label{cor:criterion 2}
    In the setting of \Cref{prop:criterion}, suppose that (1) and (3) hold, together with 
    \begin{itemize}
        \item[(5)] the stack $\cX$ is normal, $\widetilde{\cX}$ is integral, and 
       there is a schematically dense open immersion $i:\cX\hookrightarrow \widetilde{\cX}$ satisfying $\pi\circ i = \id_{\cX}$ such that $\pi(\widetilde{\cE})\subset \cX$ has codimension $\geq 2$.
    \end{itemize}
    Then the same conclusion of \Cref{prop:criterion} holds true.
\end{Cor}
\begin{proof}
    Set $\cZ:=\pi(\widetilde{\cE})$ and observe that by (5) and \Cref{lm:iso of O} there is an isomorphism $j:\cO_{\cX} \overset{\sim}{\to} j_*\cO_{\cX\smallsetminus\cZ}$. Consider the morphism $\alpha\colon\cO_{\cX}\simeq \pi_*\cO_{\widetilde{\cX}} \to \pi_*\cO(n\widetilde{\cE})$ for $n>0$. This is injective since $\cO_{\widetilde{\cX}}\to \cO_{\widetilde{\cX}}(n\widetilde{\cE})$ is injective, as $\widetilde{\cX}$ is integral. We now show that it is surjective.
    As $\widetilde{\cX}$ is integral, we have an injection $ \cO_{\widetilde{\cX}}\to \widetilde{j}_*\cO_{\widetilde{\cX}\smallsetminus \widetilde{\cE}}$ given by the restriction to the complement of $\widetilde{\cE}$. So:
    \[0\longrightarrow\cO(n\widetilde{\cE}) \longrightarrow \cO(n\widetilde{\cE})\otimes \widetilde{j}_*\cO_{\widetilde{\cX}\smallsetminus\widetilde{\cE}} \simeq \widetilde{j}_*\widetilde{j}^*\cO(n\widetilde{\cE}) \simeq \widetilde{j}_*\cO_{\widetilde{\cX}\smallsetminus \widetilde{\cE}}.\] 
    By pushing forward along $\pi$ we get an injection 
    \[0 \longrightarrow \pi_*\cO(n\widetilde{\cE}) \longrightarrow \pi_*\widetilde{j}_*\cO_{\widetilde{\cX}\smallsetminus\widetilde{\cE}} \simeq j_*\cO_{\cX\smallsetminus\cZ} \simeq \cO_{\cX}.\]
    This provides a section to $\alpha$, hence $\alpha$ is surjective. Now the same proof of \Cref{prop:criterion} goes through.
\end{proof}
\subsection{Rational pointed curves}\label{subsection_pts_on_p1}

In this subsection we study a \textit{modular} enlargement $\cC^\git_{2n} \subseteq \widetilde{\cC}^{\CY}_{2n}$ of the GIT moduli stack of genus 0 curves with a divisor of degree $2n$, which we denote by $\cC^\git_{2n}$, and which we show is an extended weighted blow-up.  We will use results from \Cref{section_extended_blowups} to prove \Cref{thm_intro_hassett}. In \Cref{subsubsection_ordered_pts} we study the problem for the case of $2n$ \textit{ordered} points on $\bP^1$.
\subsubsection{Unordered points on $\bP^1$}
For $n\geq 1$, let $\cC_{2n}^{\git}$ be the stack of smooth rational curves endowed with a divisor $D=p_1+\ldots+p_{2n}$ such that if $p_{i_1}=\ldots=p_{i_m}$, then $m\leq n$. This stack is isomorphic to the smooth quotient stack $[\bP H^0(\bP^1,\cO(2n))^{ss}/\PGL_2]$, where the $\PGL_2$-action is linearized via the line bundle $\cO(1)$, and it admits a projective good moduli space \cite{MFK}*{Proposition 4.1}.

The good moduli space $\cC_{2n}^{\git}\to C_{2n}^{\git}$ is properly stable, but it is not a coarse moduli space due to the presence of strictly semistable points in $\bP H^0(\bP^1,\cO(2n))^{ss}$: the point in the stack corresponding to $\bP^1$ and $D=n\cdot 0 + n\cdot \infty$ corresponds to the unique closed orbit of the strictly semistable locus.

Consider the stack $\Ctilde_{2n}^{\CY}$ whose objects are pairs $(P\to S, D)$ where:
\begin{enumerate}
    \item the morphism $P\to S$ is a twisted conic \cite{DLV}*{\S 1.1}, i.e. a proper, flat morphism whose geometric fibers are rational twisted curves having at most one node, and such node have automorphism group isomorphic to $\bmu_2$;
    \item the divisor $D\subset P$ is contained in the schematic locus of $P$, is finite of degree $2n$ over $S$, and on the geometric fibers $D=p_1+
    \ldots +p_{2n}$ has degree at least $n$ on each irreducible component, and if $p_{i_1}=\ldots =p_{i_m}$, then $m\leq n$.
\end{enumerate}
We aim at proving the following.
\begin{Teo}\label{thm:CY extended blow up}
    The following hold true:
    \begin{enumerate}
        \item there is a morphism $\pi\colon \Ctilde^{\CY}_{2n}\to\cC^{\git}_{2n}$ which realizes the first stack as an extended weighted blow-up of the second;
        \item the stack $\Ctilde^{\CY}_{2n}$ contains a proper Deligne-Mumford stack $\overline{C}_{2n}^{\CY}$ admitting a projective coarse moduli space.
    \end{enumerate}
\end{Teo}
In particular, one can compactify the space of maps to $\cC_{2n}^\git$ by considering maps to $\Ctilde^{\CY}_{2n}$ as the enlargement of \Cref{thm_intro_simplified}.
So while $\cC_{2n}^\git$ does not contain an open
substack which is Deligne-Mumford, one can choose an enlargement of $\cC_{2n}^\git$ 
(namely, $\Ctilde^{\CY}_{2n}$) which is itself a moduli space, and which contains an open substack which is proper and Deligne-Mumford.

Recall that there exists a smooth algebraic stack $\cC^{\CY}_{2n}$ whose objects are the same as the ones of $\Ctilde_{2n}^{\CY}$ albeit the node of the singular conics is not twisted \cite{ascher2023moduli}*{Definition 16.7, Lemma 16.8}. This stack has a natural forgetful morphism $\cC^{\CY}_{2n} \to \cC$, where $\cC$ is the stack of rational curves having at most one node.
Let $\cR$ be the stack of twisted conics \cite{DLV}*{Definition 1.1}, which also has a forgetful morphism $\cR\to \cC$. Then the following is straightforward.
\begin{Lemma}\label{lm:cartesian}
    We have $\Ctilde^{\CY}_{2n} \simeq \cC^{\CY}_{2n} \times_{\cC} \cR$.
\end{Lemma}
We will also need the following.
\begin{Lemma}\label{lm:root}
    The stack $\cR$ is an order $2$ root stack of $\cC$ along the smooth divisor of singular curves.
\end{Lemma}
\begin{proof}
    Observe that $\cC$ is isomorphic to the quotient stack $[\bP (\oH^0(\bP^2,\cO(2) \smallsetminus \Delta_2)/\PGL_3]$, where $\Delta_2$ is the locus of quadrics of rank $\leq 1$. Then the substack $\cC_1$ of singular curves is isomorphic to $[(\Delta_1\smallsetminus\Delta_2)/\PGL_3]$, which is a smooth divisor.
    
    Let $\Ctilde$ denote the root stack of $\cC$ along $\cC_1$. In order to define a morphism $\cR\to \Ctilde$, consider a twisted conic $P\to S$ and let $\overline{P}\to S$ be its coarse space, where $S=\spec(A)$ is smooth. By picking an atlas of $\cR$ and up to further shrinking it, we can assume that the locus of points with singular fibers is the vanishing locus of $f\in A$, where $f$ is not a zero divisor.
    
    \'{E}tale locally around the twisted node, the twisted conic is isomorphic to $\spec(A[u,v]/(uv-f)/\bmu_2]$, where the action is given by $u\mapsto (-u)$ and $v\mapsto (-v)$. The coarse space is instead isomorphic to $\spec(A[x,y]/(xy-f^2))$, and the coarse moduli space morphism sends $x\mapsto u^2$ and $v\mapsto v^2$. Let $S_1\subset S$ be the locus of points whose fiber in $\overline{P}\to S$ is singular: then we have just showed that $S_1$ has a root of order $2$, locally defined by the vanishing locus of $f$. In this way, we have defined a morphism $\cR \to \widetilde{\cC}$.

    On the other hand, given a family of rational curves $\overline{P}\to S=\spec(A)$, and $T=V(f)\subset S$ a divisor such that $2T=S_1$, we have that \'etale locally $\overline{P}$ looks like $\spec(A[x,y]/(xy-f^2))$, hence we can replace each of these \'etale local charts with the stack $[\spec(A[u,v]/(uv-f))/\bmu_2]$, and glue them back together, thus obtaining a twisted conic $P\to S$. This defines a morphism $\widetilde{\cC}\to \cR$, and it is straightforward to check that the two morphisms of stacks that we defined are one the inverse of the other.
\end{proof}
A combination of \Cref{lm:cartesian} and \Cref{lm:root} immediately gives us the following.
\begin{Prop}\label{prop:Ctilde is root}
    The stack $\widetilde{\cC}^{\CY}_{2n}$ is a $\bmu_2$-root stack of $\cC^{\CY}_{2n}$ along the divisor of singular curves.
\end{Prop}
\begin{Cor}\label{cor:Ctilde has gms}
    The stack $\widetilde{\cC}^{\CY}_{2n}$ is an algebraic stack of finite type and it admits a projective good moduli space, which is isomorphic to $\cC^{\git}_{2n}$.
\end{Cor}
\begin{proof}
    This follows from \Cref{prop:Ctilde is root} and the fact that $\cC^{\CY}_{2n}$ is an algebraic stack of finite type with projective good moduli space isomorphic to $C^{\git}_{2n}$ \cite{ascher2023moduli}*{Theorem 16.9}.
\end{proof}
\begin{comment}
We need one last ingredient.
\begin{Lemma}\label{lm:fake extended}
    There is a schematically dense, open embedding $i\colon\cB \PGL_ \to \cR$ and a morphism $\pi\colon \cR \to \cB\PGL_2$ such that the composition $\pi\circ i$ is the identity on $\cB \PGL_2$.
\end{Lemma}
\begin{proof}
    The $\PGL_3$-equivariant schematically dense, open embedding $\bP H^0(\bP^2,\cO(2))\smallsetminus\Delta_1 \hookrightarrow \bP H^0(\bP^2,\cO(2))\smallsetminus\Delta_2$ induces a $i$ as in the statement of the lemma.

    To construct $\pi$, let $P\to S$ be a twisted curve with $S=\spec(A)$, and $\overline{P}\to S$ its coarse space. As already done before, we can assume that the locus of points with a singular fiber is the vanishing locus of $f\in A$, where $f$ is not a zero divisor.

    Set $\widetilde{P}\to S$ to be the blow-up of $P$ along the relative singular locus, with exceptional divisor $E$. We claim (1) that $\overline{P}\to S$ has reduced fibers and that (2) there is a canonical $S$-morphism $g\colon\overline{P} \to P'$ with $(P'\to S)$ a smooth rational curve.

    Assuming (1), consider the line bundle $\cF=\omega_{\widetilde{P}/S}^\vee(-E)$: by construction this has degree two on the smooth fibers. Let $R_{\overline{s}} \cup E_{\overline{s}} \cup R'_{\overline{s}}$ be a singular fiber; then the restriction of $\cF$ has multidegree $(0,2,0)$. Let $g\colon P\to P'$ be the restriction onto the image of the $S$-morphism induced by $\cF$; it is immediate to verify that $g$ has the properties described in (2).

    To prove (1), observe that up to passing to an \'{e}tale-local neighbourhood of the relative singular locus, we can substitute $\overline{P}$ with $\spec(A[x,y]/(xy-f^2))$; indeed we must have an $A_1$-singularity, because $\overline{P}$ is the coarse space of a curve with a twisted node.

    A quick computation shows that the blow-up along $(x,y,f)$ of $\spec(A[x,y]/(xy-f^2))$ is covered by the following three charts:
    \begin{itemize}
        \item[(i)] $\spec(A[u,v]/(uv-1))$ with $x=uf$ and $y=vf$;
        \item[(ii)]$\spec(A[w,y]/(f-wy)$ with $x=w^2y$;
        \item [(iii)] $\spec(A[w',x]/(f-w'x))$ with $y=w'^2x$.
    \end{itemize}
    In particular, the fiber over $S_0$ is reduced.

    In this way, we have constructed a morphism $\pi\colon \cR \to \cB\PGL_2$. It is straightforward to check that $\pi\circ i$ is the identity.
\end{proof}
\end{comment}
We are ready to prove the result stated at the beginning of this section. We begin with the following Remark, where we explain the main ideas of the proof of \Cref{thm:CY extended blow up}.
\begin{Remark}\label{remark_why_we_need_twisted_nodes}
    Consider a one parameter family $\pi:(\cC,\Delta)\to \spec(R)$ in $\cC^{\CY}_{2n}$, and assume that the generic fiber of $\pi$ is smooth and the special one is singular. In other terms, the family $\cC\to \spec(R)$ is a degeneration of $\bP^1$ to a nodal union of two $\bP^1$s. To construct a map $\cC^{\CY}_{2n}\to \cC^{\git}_{2n}$, one might want to:
    \begin{enumerate}
        \item first blow-up the singular point on the central fiber of $\cC$ to get $B\to \cC$; let $E$ be the exceptional divisor.
\item Contract the proper transforms of the two irreducible components of the central fiber of $\cC$ in $B$, so that the push forward of the proper transform of $\Delta$ intersects the pushforward of $E$ in two points. 
    \end{enumerate}However, as it is, this approach does not work as the central fiber of $B\to \spec(R)$ is not reduced when, for example, the total space $\cC$ is smooth. On the other hand, a local computation shows that one can arrange $E$ to be reduced whenever the nodal point $x$ on the central fiber of $\cC$ is an $A_{2n+1}$ singularity. So we would like to only consider families when $x$ is an $A_{2n+1}$ singularity. The idea is to force the central fiber to be a \textit{twisted} curve with a $\bmu_2$-stabilizer at the node, so that when we take the corresponding coarse moduli space, the nodal point is indeed always an $A_{2n+1}$ singularity.
\end{Remark}
\begin{proof}[Proof of \Cref{thm:CY extended blow up}]We divide the proof in several steps.

    \noindent
    \textbf{Step 1}: we first need to construct a morphism $\widetilde{\cC}^{\CY}_{2n} \to \cC^{\git}_{2n}$. Let $P\to S$ be a twisted curve with $S=\spec(A)$, and $\overline{P}\to S$ its coarse space. As already done before, we can assume that the locus of points with a singular fiber is the vanishing locus of $f\in A$, where $f$ is not a zero divisor.

    Set $\widetilde{P}\to S$ to be the blow-up of $P$ along the relative singular locus, with exceptional divisor $E$. We claim (1) that $\overline{P}\to S$ has reduced fibers and that (2) there is a canonical $S$-morphism $g\colon\overline{P} \to P'$ with $(P'\to S)$ a smooth rational curve, obtained by contracting the irreducible components of the singular fibers which are not $E$.

    Assuming (1), consider the line bundle $\cF=\omega_{\widetilde{P}/S}^\vee(-E)$: by construction this has degree two on the smooth fibers. Let $R_{\overline{s}} \cup E_{\overline{s}} \cup R'_{\overline{s}}$ be a singular fiber; then the restriction of $\cF$ has multidegree $(0,2,0)$. Let $g\colon P\to P'$ be the restriction onto the image of the $S$-morphism induced by $\cF$; one can verify that $g$ has the properties described in (2). Indeed, $\oH^1(\cF_s)=0$ for every $s\in S$, so the formation of $\Proj(\bigoplus_{n\ge 0}\pi_*(\cF^{\otimes n}))$ commutes with base change from cohomology and base change. Then we can check the desired statement fiberwise, where one can check it easily.

    To prove (1), observe that up to passing to an \'{e}tale-local neighbourhood of the relative singular locus, we can substitute $\overline{P}$ with $\spec(A[x,y]/(xy-f^2))$; we must have such a local description, because $\overline{P}$ is the coarse space of a curve with a twisted node.
    A quick computation shows that the blow-up along $(x,y,f)$ of $\spec(A[x,y]/(xy-f^2))$ is covered by the following three charts:
    \begin{itemize}
        \item[(i)] $\spec(A[u,v]/(uv-1))$ with $x=uf$ and $y=vf$;
        \item[(ii)]$\spec(A[w,y]/(f-wy)$ with $x=w^2y$;
        \item [(iii)] $\spec(A[w',x]/(f-w'x))$ with $y=w'^2x$.
    \end{itemize}
    In particular, the fiber over $S_1$ is reduced.

    Given $(P\to S, D)$ an object of $\Ctilde^{\CY}_{2n}$, we can therefore consider the smooth rational curve $P'\to S$ together with the divisor $g(D)$. It is immediate to check that the latter is still a divisor, it is finite over $S$ of degree $2n$, and if there is a geometric fiber such that $D_s=p_1+\ldots +p_{2n}$ and $p_{i_1}=\ldots=p_{i_m}$, then $m\leq n$. In other words, the pair $(P'\to S, g(D))$ is an object of $\cC^{\git}_{2n}$. In this way, we have defined $\pi\colon \Ctilde^{\CY}_{2n} \to \cC^{\git}_{2n}$.

    \noindent
    \textbf{Step 2}: in order to apply the criterion given in \Cref{cor:criterion 2}, we first observe that there is a schematically dense, open embedding $\cC^{\git}_{2n}\to \cC^{\CY}_{2n}$ \cite{ascher2023moduli}*{Lemma 16.8} and that $\pi$ induces an isomorphism of good moduli spaces. As the image of this embedding does not intersect the divisor of singular curves, we can lift this embedding to $\Ctilde_{2n}^{\CY}$, as the latter is a root stack over $\cC^{\CY}_{2n}$ by \Cref{lm:root}. It is still schematically-dense.

    We know from \Cref{cor:Ctilde has gms} that $\Ctilde^{\CY}_{2n}$ has a projective good moduli space, which is normal because $\Ctilde^{\CY}_{2n}$ is a root stack over $\cC^{\CY}_{2n}$ (see \Cref{prop:Ctilde is root}) and the latter is smooth stack \cite{ascher2023moduli}*{Lemma 16.8}. The morphism $\pi$ induces a morphism of good moduli spaces $C_{2n}^{\CY}\to C_{2n}^{\git}$, which has a section given by the morphism induced by $i$, which is an isomorphism (this follows from how we defined $i$ and \cite{ascher2023moduli}*{Theorem 16.9}). It follows that the morphism $\pi$ induces an isomorphism of good moduli spaces.

    \noindent
    \textbf{Step 3}: to apply \Cref{cor:criterion 2} we need to show that the morphism $\xi:\Ctilde^{\CY}_{2n} \to \cC^{\git}_{2n}\times [\bA^1/\Gm]$, given by $\pi$ and the Cartier divisor $\widetilde{\cE}\subset \Ctilde^{\CY}_{2n}$ of singular curves, is representable. From how $\pi$ is constructed, it suffices to check that $\xi$ is representable at the points of $\widetilde{\cC}_{2n}^\CY$ which correspond to singular curves (namely, $\widetilde{\cE}$). Let $\pi_i$ the two projections of $\cC^{\git}_{2n}\times [\bA^1/\Gm]$. As there is a unique polystable point $p\in \widetilde{\cE}$, namely the curve where the support of the divisor consists of two points, it from \Cref{lemma_one_can_check_rep_at_poli_pts} it suffices to check that $\xi$ is representable at this point. We recall a few results from \cite{DLV}*{\S 1.5} on $\Aut_{\Ctilde^{\CY}_{2n}}(p)$.
    
     Let $(C,E)$ be the pair in $\widetilde{\cE}$ corresponding to $p$, namely a nodal twisted cubic $C$ with a divisor $E$ consisting of two smooth points on $C$, one per irreducible component. If we denote by $D$ the nodal union of two $\bP^1$s attached at the point $[0,1]$, there is an action of $\bmu_2$ on $D$ which preserves the irreducible components and on each irreducible component it acts as $[a,b]\mapsto [-a,b]$. The stack quotient $[D/\bmu_2]$ has three stacky points, two of which are smooth. There is a map $[D/\bmu_2]\to C$ which rigidifies $[D/\bmu_2]$ at the two smooth stacky points; the corresponding points on $C$ are the support of the divisor $E$, and we still denote by $E$ the two smooth points on $D$ fixed by $\bmu_2$. To summarize:
     \begin{itemize}
         \item $D$ is the nodal union of two $\bP^1$, and $E $ are two smooth points on $D$, one per irreducible component,
         \item $\bmu_2$ acts on $D$, non-trivially on each irreducible component, fixing $E$, and
         \item $C$ is the quotient $[D/\bmu_2]$ rigidified at $[E/\bmu_2]$.
     \end{itemize}Then:
    \begin{enumerate}
\item $\Aut(D,E)\cong \Gm^2\rtimes \bmu_2$ with where $\Gm$ acts by scaling each branch, and $\bmu_2$ swaps the two branches. 
        \item There is a map $\Aut(D,E)\to \Aut([D/\bmu_2],[E/\bmu_2])$ that is surjective with kernel $(-1,-1,\Id)$, and an isomorphism $\Aut([D/\bmu_2],[E/\bmu_2])\cong\Aut(C,E)$ \cite{DLV}*{Pag. 10}.
        \item The action of $\Aut_{\widetilde{\cC}^{\CY}_{2n}}(p)=\Aut(C,E)$ on the fiber of $\cO_{\widetilde{\cC}^{\CY}_{2n}}(\widetilde{\cE})$ at $p$ (i.e., on $\Aut_{[\bA^1/\Gm]}(\pi_2\circ \xi (p))$) is computed via the composition $\Gm^2\rtimes \bmu_2\cong \Aut(D,E)\to \Gm$, $(a,b,\sigma)\mapsto ab$ \cite{DLV}*{Proposition A.4 and pag. 12}.
\item The map $\widetilde{\cC}^{\CY}_{2n}\to \cC^\git_{2n}$ sends the twisted conic $C$ to the $\bP^1$ obtained by performing the blow-up of the reduced stacky point of $C$ and taking the coarse moduli space of the exceptional divisor. In a neighbourhood of the node, the blow-up is the $\bmu_2$-quotient of $\Proj(k[x,y,X_0,X_1]/(X_0y-X_1x,xy))$. We study the chart where $X_0\neq 0$; the other chart is similar. In this chart the blow-up is $\spec(k[x,\frac{X_1}{X_0}]/x^2\frac{X_1}{X_0})$ and the action of $\bmu_2$ sends $x\mapsto -x$.
Hence the coarse moduli space is $\spec(k[x^2,\frac{X_1}{X_0}]/x^2\frac{X_1}{X_0})$. The exceptional divisor is $\spec(k[\frac{X_1}{X_0}])$. So the composition \[\Gm^2\rtimes \bmu_2\cong \Aut(D,E)\to \Aut(C,E)=\Aut_{\widetilde{\cC}^{\CY}_{2n}}(p)\xrightarrow{\pi_1\circ \xi} \Aut_{\cC^{\git}_{2n}}(\pi_1\circ \xi(p))\text{ sends }(a,b,\sigma)\mapsto \left(\frac{a}{b},\sigma\right).\]
    \end{enumerate}
Then from points (3) and (4), the composition
\begin{align*}
\Gm^2\rtimes \bmu_2\cong \Aut(D,E)\to &\Aut_{\widetilde{\cC}^{\CY}_{2n}}(p)\to \Aut_{\cC^{\git}_{2n}}(\pi_1\circ\xi(p))\times \Aut_{[\bA^1/\Gm]}(\pi_2\circ\xi(p))\\&\text{sends } (a,b,\sigma)\mapsto \left(\left(\frac{a}{b},\sigma\right),ab\right).
\end{align*}
The kernel of this map is generated by $(-1,-1,\Id)$ so from (2) the map 
\[\Aut(C,E)\cong \Aut_{\widetilde{\cC}^{\CY}_{2n}}(p)\to \Aut_{\cC^{\git}_{2n}}(\pi_1\circ\xi(p))\times \Aut_{[\bA^1/\Gm]}(\pi_2\circ\xi(p))\]
is injective, as desired.  

    \noindent
    \textbf{Step 4}: we are only left with proving that $\Ctilde_{2n}^{\CY}$ contains a proper Deligne-Mumford stack with a coarse moduli space. Consider $\cU_{2n}$ the complement of the closed substack whose objects are pairs $(P\to S, D)$ where on each geometric fiber the restriction of $D$ is supported on only two points.
    Observe finally that $\cU_{2n}$ is Deligne-Mumford, and it has a projective coarse moduli space. Indeed, it is a root stack of the KSBA-moduli space compactifying $(\bP^2,(\frac{1}{n}+\epsilon )D)$  where $D\subseteq \bP^1$ is a smooth divisor of degree $2d$.
\end{proof}

\subsubsection{Moduli of pointed rational curves}\label{subsubsection_ordered_pts}
Let $\Mcal_{0,2n}$ denote the moduli stack of smooth rational curves with $2n$ markings, for $n>1$, which is a quasi-projective variety.

One possible modular compactification of $\Mcal_{0,2n}$ can be constructed as follows: consider the scheme \[\bP {\rm H}^0(\bP^1,\cO(1))^{\times 2n} \simeq (\bP^1)^{\times 2n}\]parametrizing $2n$-uples of points on $\bP^1$. There is a natural $\PGL_2$-action on this variety, which can be linearized via the line bundle $L=\cO(2)\boxtimes\ldots\boxtimes\cO(2)$. 

Set $\Mbar_{0,2n}^\git=[(\bP^1)^{\times 2n}(L)^{ss}/\PGL_2]$: then this is an algebraic stack, with good moduli space the GIT quotient $\overline{M}_{0,2n}^{\git}:=(\bP^1)^{\times 2n}(L)^{ss}/\!\!/\PGL_2$. 

The semistable locus of this linearized $\PGL_2$-action on $(\bP^1)^{\times n}$ consists of the configurations with at most $n$ markings collapsing together, and the stable points are those with at most $n-1$ points collapsing together. The closed orbits of the strictly semistable locus are the semistable configurations with the markings supported on only two points (see the proof of \cite{MFK}*{Proposition 4.1}).

Therefore, the stack $\Mbar_{0,2n}^{\git}$ is a compactification of $\Mcal_{0,2n}$, but it does not have a proper, Deligne-Mumford substack with a coarse moduli space, so we cannot apply the theory of stable quasimaps of \S \ref{sec:qmaps}. 

Consider the stack $\Mtilde_{0,2n}^{\Hass}$ whose objects are tuples $(P\to S,p_1,\ldots,p_{2n})$ where:
\begin{itemize}
    \item the morphism $P\to S$ is a twisted conic and $p_i\colon S \to P$ are sections landing in the smooth locus of $P\to S$;
    \item on every geometric point $s$, each irreducible component of $P_s$ contains at least $n$ sections, and no more than $n$ sections can coincide.
\end{itemize}
We omit the definition of the morphisms.
\begin{Prop}\label{prop:hass to git}
    The following hold true:
    \begin{enumerate}
        \item there is a morphism $\pi\colon \Mtilde^{\Hass}_{0,2n}\to\Mbar^{\git}_{0,2n}$ which realizes the first stack as an extended weighted blow-up of the second;
        \item the stack $\Mtilde^{\Hass}_{0,2n}$ contains a proper Deligne-Mumford stack $\Mbar_{0,2n}^{\Hass}$ admitting a projective coarse moduli space.
    \end{enumerate}
\end{Prop}
The Deligne-Mumford stack $\Mbar_{0,2n}^{\Hass}$ appearing in \Cref{prop:hass to git} is the moduli stack of stable weighted $2n$-marked genus zero curves with weights $(\frac{1+\epsilon}{n},\ldots,\frac{1+\epsilon}{n})$ constructed in \cite{Hassett_w}.

\begin{proof}[Proof of \Cref{prop:hass to git}]
    As before, we divide the proof in several steps.

    \noindent
    \textbf{Step 1}: the morphism $\pi\colon \Mtilde^{\Hass}_{0,2n}\to\Mbar^{\git}_{0,2n}$ is constructed exactly as in the first step of the proof \Cref{thm:CY extended blow up}, with the only difference that one has to take the image of the sections $p_i\colon S\to P$ instead of the image of the divisor $D\subset P$.     
    
    \noindent
    \textbf{Step 2}: we want to apply \Cref{cor:criterion 2}. First observe that there is a schematically-dense open embedding $i\colon \Mbar^{\git}_{0,2n}\to \Mtilde^{\Hass}_{0,2n}$, defined similarly as in the proof of \Cref{thm:CY extended blow up}. It is easy to check that $\pi\circ i = \id$.
    
    Next, we show that $\Mtilde_{0,2n}^{\Hass}$ admits a good moduli space and that $\pi\colon\Mtilde_{0,2n}^{\Hass}\to\Mtilde^{\git}_{0,2n}$ induces an isomorphism at the level of good moduli spaces.

    Observe that there is a natural morphism $\Mtilde^{\Hass}_{0,2n} \to \Ctilde^{\CY}_{2n}$, which is representable and finite: indeed, it is induced by taking the quotient with respect to the natural action of the symmetric group on the labels of the sections.

    As $\Ctilde^{\CY}_{2n} \to C^{\git}_{2n}$ is cohomologically affine by \Cref{thm:CY extended blow up}, so it is the composition $f\colon \Mtilde^{\Hass}_{0,2n} \to C^{\git}_{2n}$. Define $X={\spec}_{C^{\git}_{2n}}(f_*\cO_{\Mtilde^{\Hass}_{0,2n}})$. Applying \Cref{lm:stein of coh aff is gms}, we deduce that $\Mtilde^{\Hass}_{0,2n} \to X$ is a good moduli space.

    Let $\overline{M}_{0,2n}^{\git}$ be the moduli space of $\Mbar^{\git}_{0,2n}$: then the composition $\Mtilde^{\Hass}_{0,2n} \to \Mbar^{\git}_{0,2n} \to \overline{M}_{0,2n}^{\git}$ induces a morphism $f\colon X \to \overline{M}_{0,2n}^{\git}$ which can be easily checked to be birational and bijective. As the target is a normal variety and both the domain and the target are projective, by Zariski's main theorem we deduce that $f$ is an isomorphism.

    \noindent
    \textbf{Step 3}: we show that there is a representable morphism $\Mtilde^{\Hass}_{0,2n} \to \Mbar^{\git}_{0,2n}\times [\bA^1/\Gm]$. The morphism to $[\bA^1/\Gm]$ is induced by the Cartier divisor of singular twisted conics, so that we have a commutative diagram
    \[
    \begin{tikzcd}
        \Mtilde^{\Hass}_{0,2n} \ar[r] \ar[d] & \Mbar^{\git}_{0,2n}\times [\bA^1/\Gm] \ar[d]\\
        \Ctilde^{\CY}_{2n} \ar[r] & \cC^{\git}_{2n}\times [\bA^1/\Gm].
    \end{tikzcd}
    \]
    As the vertical arrows are representable, and as we have seen in the proof of \Cref{thm:CY extended blow up} that the bottom horizontal arrow is representable, we deduce that also the top horizontal arrow is representable.

    \noindent
    \textbf{Step 4}: we can apply the criterion given in \Cref{cor:criterion 2}. To conclude, observe that $\Mtilde^{\Hass}_{0,2n}$ contains the Hassett moduli stack $\Mbar^{\Hass}_{0,2n}$ of weighted $2n$-marked rational curves, with weights $(\frac{1+\epsilon}{n},\ldots,\frac{1+\epsilon}{n})$ constructed in \cite{Hassett_w}.
\end{proof}
\begin{Remark}
    It follows from the results in this section that one can compactify the moduli space of maps $C\to \cC_{2n}^{\git}$ as in \Cref{thm_intro_simplified} by considering the enlargement $\cC_{2n}^{\git}\subseteq \widetilde{\cC}^\CY_{2n}$ and via \Cref{teo_intro_cX_contains_open_proper_dm}. In particular, the locus $\cU$ parametrizing maps $\bP^1\to \cC_{2n}^{\git}$ which send the generic point of $\bP^1$ to a strictly stable point, and such that the composition $\bP^1\to\cC^\git_{2n}\to C^\git_{2n}$ with the good moduli space is finite, can be compactified $\cU\subseteq \overline{\cU}$ in a way such that the boundary parametrizes maps $\cC\to \widetilde{\cC}^\git_{2n}$ from a rational twisted curve, and which are stable in the sense of \Cref{def_stable_qmap}. A similar picture goes through if we consider instead $\overline{\cM}^\git_{0,2n}$.
\end{Remark}


\bibliographystyle{amsalpha}
\bibliography{v19}

\end{document}
